
%%%%%%%%%%%%%%%%%%%%%%%%%%%%%%%%%%%%%%%%%%%%%%%%%%%%%%%%%%%%%%%%%%%%%%%%%%%%%
%
% Copyright 2007/2008 by Christian Feuersaenger.
%
% This program is free software: you can redistribute it and/or modify
% it under the terms of the GNU General Public License as published by
% the Free Software Foundation, either version 3 of the License, or
% (at your option) any later version.
%
% This program is distributed in the hope that it will be useful,
% but WITHOUT ANY WARRANTY; without even the implied warranty of
% MERCHANTABILITY or FITNESS FOR A PARTICULAR PURPOSE.  See the
% GNU General Public License for more details.
%
% You should have received a copy of the GNU General Public License
% along with this program.  If not, see <http://www.gnu.org/licenses/>.
%
%
%%%%%%%%%%%%%%%%%%%%%%%%%%%%%%%%%%%%%%%%%%%%%%%%%%%%%%%%%%%%%%%%%%%%%%%%%%%%%

\pdfminorversion=5 % to allow compression
\pdfobjcompresslevel=2
\documentclass[a4paper,doc2]{ltxdoc}

%%%%%%%%%%%%%%%%%%%%%%%%%%%%%%%%%%%%%%%%%%%%%%%%%%%%%%%%%%%%%%%%%%%%%%%%%%%%%
%
% Package pgfplots.sty documentation.
%
% Copyright 2007/2008 by Christian Feuersaenger.
%
% This program is free software: you can redistribute it and/or modify
% it under the terms of the GNU General Public License as published by
% the Free Software Foundation, either version 3 of the License, or
% (at your option) any later version.
%
% This program is distributed in the hope that it will be useful,
% but WITHOUT ANY WARRANTY; without even the implied warranty of
% MERCHANTABILITY or FITNESS FOR A PARTICULAR PURPOSE.  See the
% GNU General Public License for more details.
%
% You should have received a copy of the GNU General Public License
% along with this program.  If not, see <http://www.gnu.org/licenses/>.
%
%
%%%%%%%%%%%%%%%%%%%%%%%%%%%%%%%%%%%%%%%%%%%%%%%%%%%%%%%%%%%%%%%%%%%%%%%%%%%%%


\usepackage{makeidx}

% 中文支持：xelatex + ctex
\usepackage[UTF8]{ctex}

% DON't let hyperref overload the format of index and glossary.
% I want to do that on my own in the stylefiles for makeindex...
\makeatletter
\let\@old@wrindex=\@wrindex
\makeatother

\usepackage{ifpdf}
\usepackage[pdfborder=0 0 0]{hyperref}
    \hypersetup{%
        colorlinks=true,    % use true to enable colors below:
        linkcolor=blue,%red,
        filecolor=blue,%magenta,
        %pagecolor=blue,%red,
        urlcolor=blue,%cyan,
        citecolor=blue,
        %frenchlinks=false,    % small caps instead of colors
        pdfborder=0 0 0,    % PDF link-darstellung, falls colorlinks=false. 0 0 0: nix. 0 0 1: default.
        %plainpages=false,    % Das ist notwendig, wenn die Seitenzahlen z.T. in Arabischen und z.T. in römischen Ziffern gemacht werden.
        pdftitle=Package PGFPLOTS manual,
        pdfauthor=Dr. Christian Feuersänger,
        %pdfsubject=,
        pdfkeywords={pgfplots pgf tikz tex latex},
    }

\makeatletter
\let\@wrindex=\@old@wrindex
\makeatother

\date{\small{Revision \pgfplotsrevision\space (\pgfplotsrevisiondate)}}

\usepackage{pdftexcmds}

\usepackage{luatex85}

\makeatletter
% disables colorlinks for all following \ref commands
\def\pgfplotsmanualdisablecolorforref{%
    \pgfutil@ifundefined{pgfplotsmanual@oldref}{%
        \let\pgfplotsmanual@oldref=\ref
    }{}%
    \def\ref##1{%
        \begingroup
        \let\Hy@colorlink=\pgfplots@disabled@Hy@colorlink
        \let\Hy@endcolorlink=\pgfplots@disabled@Hy@endcolorlink
        \pgfplotsmanual@oldref{##1}%
        \endgroup
    }%
}%
\def\pgfplots@disabled@Hy@colorlink#1{\begingroup}%
\def\pgfplots@disabled@Hy@endcolorlink{\endgroup}%
\makeatother

% Formatiere Seitennummern im Index:
\newcommand{\indexpageno}[1]{%
    {\bfseries\hyperpage{#1}}%
}

\def\C{\mathbb{C}}
\newcommand{\R}{\mathbb{R}}
\newcommand{\N}{\mathbb{N}}
\newcommand{\Z}{\mathbb{Z}}

\long\def\COMMENTLOWLEVEL#1\ENDCOMMENT{}
\def\ENDCOMMENT{}

\usepackage{textcomp}
\usepackage{booktabs}

\usepackage{calc}
\usepackage[formats]{listings}
%\usepackage{courier} % don't use it -- the '^' character can't be copy-pasted in courier

\usepackage{array}
\lstset{%
    basicstyle=\ttfamily,
    language=[LaTeX]tex, % Seems as if \lstset{language=tex} must be invoked BEFORE loading tikz!?
    tabsize=4,
    breaklines=true,
    breakindent=0pt,
}

\ifpdf
%    \pdfinfo {
%        /Author    (Christian Feuersaenger)
%    }

\else
%    \def\pgfsysdriver{pgfsys-dvipdfm.def}
\fi
%\def\pgfsysdriver{pgfsys-pdftex.def}
\usepackage{pgfplots}
\usepackage{pgfplotstable}

\ifpdf
    % this allows to disable the clickable lib from command line using
    % \pdflatex '\def\pgfplotsclickabledisabled{1}
%%%%%%%%%%%%%%%%%%%%%%%%%%%%%%%%%%%%%%%%%%%%%%%%%%%%%%%%%%%%%%%%%%%%%%%%%%%%%
%
% Package pgfplots.sty documentation.
%
% Copyright 2007/2008 by Christian Feuersaenger.
%
% This program is free software: you can redistribute it and/or modify
% it under the terms of the GNU General Public License as published by
% the Free Software Foundation, either version 3 of the License, or
% (at your option) any later version.
%
% This program is distributed in the hope that it will be useful,
% but WITHOUT ANY WARRANTY; without even the implied warranty of
% MERCHANTABILITY or FITNESS FOR A PARTICULAR PURPOSE.  See the
% GNU General Public License for more details.
%
% You should have received a copy of the GNU General Public License
% along with this program.  If not, see <http://www.gnu.org/licenses/>.
%
%
%%%%%%%%%%%%%%%%%%%%%%%%%%%%%%%%%%%%%%%%%%%%%%%%%%%%%%%%%%%%%%%%%%%%%%%%%%%%%
% SEE pgfplots-macros.tex as well!
%\pdfminorversion=5 % to allow compression
%\pdfobjcompresslevel=2
\documentclass[a4paper,openany]{book}

% Promote `Missing character' reports to full errors
% NOTE: lualatex turns this into a fatal \shipout error (no PDF produced);
% keeping it commented keeps both engines compiling.
%\tracinglostchars=3

\let\bookmaketitle=\maketitle

% -----------------------------------
% this here is from ltxdoc:
\usepackage{doc}[=v2]
\AtBeginDocument{\MakeShortVerb{\|}}
% 缓解超宽 verb/键名引用导致的 Overfull hbox（允许断行时轻微伸展行宽）
\setlength{\emergencystretch}{10em}
%\setlength{\textwidth}{355pt}
%\addtolength\marginparwidth{30pt}
%\addtolength\oddsidemargin{20pt}
%\addtolength\evensidemargin{20pt}
\def\cmd#1{\cs{\expandafter\cmd@to@cs\string#1}}
\def\cmd@to@cs#1#2{\char\number`#2\relax}
\DeclareRobustCommand\cs[1]{\texttt{\char`\\#1}}
\providecommand\marg[1]{%
  {\ttfamily\char`\{}\meta{#1}{\ttfamily\char`\}}}
\providecommand\oarg[1]{%
  {\ttfamily[}\meta{#1}{\ttfamily]}}
\providecommand\parg[1]{%
  {\ttfamily(}\meta{#1}{\ttfamily)}}
\raggedbottom

% -----------------------------------

%%%%%%%%%%%%%%%%%%%%%%%%%%%%%%%%%%%%%%%%%%%%%%%%%%%%%%%%%%%%%%%%%%%%%%%%%%%%%
%
% Package pgfplots.sty documentation.
%
% Copyright 2007/2008 by Christian Feuersaenger.
%
% This program is free software: you can redistribute it and/or modify
% it under the terms of the GNU General Public License as published by
% the Free Software Foundation, either version 3 of the License, or
% (at your option) any later version.
%
% This program is distributed in the hope that it will be useful,
% but WITHOUT ANY WARRANTY; without even the implied warranty of
% MERCHANTABILITY or FITNESS FOR A PARTICULAR PURPOSE.  See the
% GNU General Public License for more details.
%
% You should have received a copy of the GNU General Public License
% along with this program.  If not, see <http://www.gnu.org/licenses/>.
%
%
%%%%%%%%%%%%%%%%%%%%%%%%%%%%%%%%%%%%%%%%%%%%%%%%%%%%%%%%%%%%%%%%%%%%%%%%%%%%%


\usepackage{makeidx}

% 中文支持：xelatex + ctex
\usepackage[UTF8]{ctex}

% DON't let hyperref overload the format of index and glossary.
% I want to do that on my own in the stylefiles for makeindex...
\makeatletter
\let\@old@wrindex=\@wrindex
\makeatother

\usepackage{ifpdf}
\usepackage[pdfborder=0 0 0]{hyperref}
    \hypersetup{%
        colorlinks=true,    % use true to enable colors below:
        linkcolor=blue,%red,
        filecolor=blue,%magenta,
        %pagecolor=blue,%red,
        urlcolor=blue,%cyan,
        citecolor=blue,
        %frenchlinks=false,    % small caps instead of colors
        pdfborder=0 0 0,    % PDF link-darstellung, falls colorlinks=false. 0 0 0: nix. 0 0 1: default.
        %plainpages=false,    % Das ist notwendig, wenn die Seitenzahlen z.T. in Arabischen und z.T. in römischen Ziffern gemacht werden.
        pdftitle=Package PGFPLOTS manual,
        pdfauthor=Dr. Christian Feuersänger,
        %pdfsubject=,
        pdfkeywords={pgfplots pgf tikz tex latex},
    }

\makeatletter
\let\@wrindex=\@old@wrindex
\makeatother

\date{\small{Revision \pgfplotsrevision\space (\pgfplotsrevisiondate)}}

\usepackage{pdftexcmds}

\usepackage{luatex85}

\makeatletter
% disables colorlinks for all following \ref commands
\def\pgfplotsmanualdisablecolorforref{%
    \pgfutil@ifundefined{pgfplotsmanual@oldref}{%
        \let\pgfplotsmanual@oldref=\ref
    }{}%
    \def\ref##1{%
        \begingroup
        \let\Hy@colorlink=\pgfplots@disabled@Hy@colorlink
        \let\Hy@endcolorlink=\pgfplots@disabled@Hy@endcolorlink
        \pgfplotsmanual@oldref{##1}%
        \endgroup
    }%
}%
\def\pgfplots@disabled@Hy@colorlink#1{\begingroup}%
\def\pgfplots@disabled@Hy@endcolorlink{\endgroup}%
\makeatother

% Formatiere Seitennummern im Index:
\newcommand{\indexpageno}[1]{%
    {\bfseries\hyperpage{#1}}%
}

\def\C{\mathbb{C}}
\newcommand{\R}{\mathbb{R}}
\newcommand{\N}{\mathbb{N}}
\newcommand{\Z}{\mathbb{Z}}

\long\def\COMMENTLOWLEVEL#1\ENDCOMMENT{}
\def\ENDCOMMENT{}

\usepackage{textcomp}
\usepackage{booktabs}

\usepackage{calc}
\usepackage[formats]{listings}
%\usepackage{courier} % don't use it -- the '^' character can't be copy-pasted in courier

\usepackage{array}
\lstset{%
    basicstyle=\ttfamily,
    language=[LaTeX]tex, % Seems as if \lstset{language=tex} must be invoked BEFORE loading tikz!?
    tabsize=4,
    breaklines=true,
    breakindent=0pt,
}

\ifpdf
%    \pdfinfo {
%        /Author    (Christian Feuersaenger)
%    }

\else
%    \def\pgfsysdriver{pgfsys-dvipdfm.def}
\fi
%\def\pgfsysdriver{pgfsys-pdftex.def}
\usepackage{pgfplots}
\usepackage{pgfplotstable}

\ifpdf
    % this allows to disable the clickable lib from command line using
    % \pdflatex '\def\pgfplotsclickabledisabled{1}\input{pgfplots.tex}'
    \expandafter\ifx\csname pgfplotsclickabledisabled\endcsname\relax
        \usepgfplotslibrary{clickable}
    \fi
\fi

%\usepackage{fp}
% ATTENTION:
% this requires pgf version NEWER than 2.00:
%\usetikzlibrary{fixedpointarithmetic}

\usepgfplotslibrary{dateplot,units,groupplots,statistics}

\pgfplotsset{
    table/search path={.,plotdata}
}

\usepackage[a4paper,left=2.25cm,right=2.25cm,top=2.5cm,bottom=2.5cm]{geometry}
\usepackage{amsmath,amssymb}
\usepackage{xxcolor}

\usepackage{ifluatex}

% tlmgr install zapfding symbol
\usepackage{pifont}

\def\lua{\textsc{lua}}

\ifluatex
    %\usepackage[utf8]{luainputenc}
\else
    \expandafter\ifx\csname XeTeXrevision\endcsname\relax
        \usepackage[utf8]{inputenc}
    \fi
\fi
\usepackage{amsmath}

% tlmgr eurosym
\usepackage{eurosym}

% tlmgr install units
\usepackage{nicefrac}
\input{pgfplots-macros}

\makeatletter
\@ifpackageloaded{tex4ht}{
}{%
    \IfFileExists{ocg.sty}{%
        \usepackage{ocg}%
    }{%
        \usepackage{pgfplots_ocg_copy}%
    }
}%
\makeatother

\usepackage{nicefrac}

\graphicspath{{figures/}}

\edef\preambleconfig{width=7cm,compat=\pgfkeysvalueof{/pgfplots/compat/mostrecent}}

% generate the following file such that it does not need to be updated
% manually for each release:
\makeatletter
\immediate\openout\w@pgf@writea pgfplots.preliminaries.compatcurrent.tmp
\immediate\write\w@pgf@writea{%
\string\begin{codeexample}[code only]^^J%
\string\usepackage{pgfplots}^^J%
\string\pgfplotsset{compat=\pgfkeysvalueof{/pgfplots/compat/mostrecent}}^^J%
\string\end{codeexample}%
}%
\immediate\closeout\w@pgf@writea


\expandafter\pgfplotsset\expandafter{\preambleconfig}

\pgfplotsset{lua debug=off and silent}

\def\pgfplotsmanualcurlibrary{}

% And now, invoke
%     /codeexample/typeset listing/.add={% Preamble:\pgfplotsset{\preambleconfig}}{}}
% since listings are VERBATIM, I need to do some low-level things
% here to get the correct \catcodes:
\pgfkeys{/codeexample/typeset listing/.add code={%
        \ifcode@execute
            \begingroup
            \def\b@show@header{0}%
            \pgfutil@in@{axis}{#1}%
            \ifpgfutil@in@
                \def\b@show@header{1}%
            \fi
            %
            \pgfutil@in@{addplot}{#1}%
            \ifpgfutil@in@
                \def\b@show@header{1}%
            \fi
            %
            \pgfutil@in@{smithchart}{#1}%
            \ifpgfutil@in@
                \def\b@show@header{1}%
            \fi
            \if1\b@show@header
                {\tiny
                    \% Preamble: \pgfmanualpdfref{\textbackslash pgfplotsset}{\pgfmanual@pretty@backslash pgfplotsset}%
                        \pgfmanual@pretty@lbrace \expandafter\pgfmanualprettyprintpgfkeys\expandafter{\preambleconfig}\pgfmanual@pretty@rbrace
                        \ifx\pgfplotsmanualcurlibrary\pgfutil@empty
                        \else
                            \pgfmanual@pretty@backslash usepgfplotslibrary\pgfmanual@pretty@lbrace \pgfplotsmanualcurlibrary\pgfmanual@pretty@rbrace
                        \fi
                }%
            \fi
            \endgroup
        \fi
    }{},%
    %/codeexample/typeset listing/.show code,
}%
\makeatother

\pgfplotsset{
    %every axis/.append style={width=7cm},
    filter discard warning=false,
}

\pgfqkeys{/codeexample}{%
    every codeexample/.append style={
        width=8cm,
        /pgfplots/legend style={fill=graphicbackground},
        /pgfplots/contour/every contour label/.append style={
            every node/.append style={fill=graphicbackground}
        },
    },
    tabsize=4,
}

\usetikzlibrary{backgrounds,patterns}
% Global styles:
\tikzset{
  shape example/.style={
    color=black!30,
    draw,
    fill=yellow!30,
    line width=.5cm,
    inner xsep=2.5cm,
    inner ysep=0.5cm}
}

\newcommand{\FIXME}[1]{\textcolor{red}{(FIXME: #1)}}

% fuer environment 'sidewaysfigure' bspw
% \usepackage{rotating}

\newcommand\Tikz{Ti\textit kZ}
\newcommand\PGF{\textsc{pgf}}
\newcommand\PGFPlots{\pgfplotsmakefilelinkifuseful{pgfplots}{\textsc{pgfplots}}}
\newcommand\PGFPlotstable{\pgfplotsmakefilelinkifuseful{pgfplotstable}{\textsc{PgfplotsTable}}}

\makeindex

% Fix overfull hboxes automatically:
\tolerance=2000
\emergencystretch=10em

\tikzset{prefix=gnuplot/pgfplots_} % prefix for 'plot function'

\author{%
    Dr.\ Christian Feuersänger\\
    {\footnotesize\texttt{cfeuersaenger@users.sourceforge.net}}}%


\usepackage{todonotes}
\newcommand\todosp[2][]{%        % Stefan Pinnow
    \todo[
        disable,
        color=blue!40,
        #1]{#2}
}

% ---------------------------------------------------------------
% 修复：doc 宏包 v2 的 \ltxdoc@meta 在 TeX Live 2026 中未定义，
% 导致 \marg/\oarg/\parg（内部调用 \meta）写 aux/idx 时报
% "Undefined control sequence"。
% 对策（双保险）：
%  1) 兜底提供 \ltxdoc@meta（doc 的 \meta 依赖它）；
%  2) 在 \AtBeginDocument 末尾覆盖 \meta 为鲁棒命令
%     （doc v2 会在 begin document 时重定义 \meta，必须最后覆盖）。
\makeatletter
\providecommand\ltxdoc@meta[1]{%
    {\normalfont\itshape\textlangle #1\textrangle}%
}
\AtBeginDocument{%
    \DeclareRobustCommand\meta[1]{%
        {\normalfont\itshape\textlangle #1\textrangle}%
    }%
}
\makeatother


\makeatletter
% I want a two-column index just as in pgfmanual styles. This here
% was the best way to get one:
\def\index@prologue{\section*{索引}\addcontentsline{toc}{chapter}{索引}
}
\makeatother

%\RequirePackage[german,english,francais]{babel}

\def\matlabcolormaptext{此色图与 Matlab$^\text{\textregistered}$ 中以类似名称发布的色图相似。}

\IfFileExists{tikzlibraryspy.code.tex}{%
\usetikzlibrary{spy}
}{%
    \message{ERROR: tikz SPY library NOT available. The manual will only compile partially.^^J}%
}%

\usepackage{xparse}% for colorbrewer manual

\usetikzlibrary{
    decorations.markings,
    decorations.footprints,
    shapes.arrows,
    matrix,
    positioning,
}

\usepgfplotslibrary{
    fillbetween,
    ternary,
    smithchart,
    patchplots,
    polar,
    colormaps,
    colorbrewer,
    colortol,           % docu not ready yet
}
\pgfqkeys{/codeexample}{%
    every codeexample/.append style={
        /pgfplots/every ternary axis/.append style={
            /pgfplots/legend style={fill=graphicbackground},
        }
    },
    tabsize=4,
}

\pgfplotsmanualenableexternalizationofexpensive

%\usetikzlibrary{external}
%\tikzexternalize[prefix=figures/]{pgfplots}

\title{%
    \PGFPlots{} 宏包手册\\
    {\small 在 \LaTeX{} 中绘制二维/三维图形，版本 \pgfplotsversion}\\
    {\small\url{https://github.com/pgf-tikz/pgfplots}}
    %\\{\small Attention: you are using an unstable development version.}
}

\makeatletter
\long\def\abstractsmuggle{%
    \centering
    \textbf{摘要}\\[0.5cm]

    \begin{minipage}{12cm}
        \PGFPlots{} 直接在 \TeX{} 中以常规或对数坐标绘制高质量函数图，
        界面友好。用户只需提供坐标轴标签、图例条目以及一个或多个图的坐标数据，
        \PGFPlots{} 即可完成坐标轴缩放、对数与刻度计算并绘制图形。
        它支持线图、散点图、分段常数图、柱状图、面积图、网格图与曲面图、
        面片图、等高线图、矢量图、直方图、箱线图、极坐标轴、三元相图、
        史密斯圆图等多种图形。它基于 Till Tantau 开发的 \PGF{}/\Tikz{} 宏包。
    \end{minipage}

}%

\expandafter\date\expandafter{\@date\\[2cm]
    \abstractsmuggle
}%
\makeatother

%\includeonly{pgfplots.reference}


\begin{document}

\def\plotcoords{%
\addplot coordinates {
(5,8.312e-02)    (17,2.547e-02)   (49,7.407e-03)
(129,2.102e-03)  (321,5.874e-04)  (769,1.623e-04)
(1793,4.442e-05) (4097,1.207e-05) (9217,3.261e-06)
};

\addplot coordinates{
(7,8.472e-02)    (31,3.044e-02)    (111,1.022e-02)
(351,3.303e-03)  (1023,1.039e-03)  (2815,3.196e-04)
(7423,9.658e-05) (18943,2.873e-05) (47103,8.437e-06)
};

\addplot coordinates{
(9,7.881e-02)     (49,3.243e-02)    (209,1.232e-02)
(769,4.454e-03)   (2561,1.551e-03)  (7937,5.236e-04)
(23297,1.723e-04) (65537,5.545e-05) (178177,1.751e-05)
};

\addplot coordinates{
(11,6.887e-02)    (71,3.177e-02)     (351,1.341e-02)
(1471,5.334e-03)  (5503,2.027e-03)   (18943,7.415e-04)
(61183,2.628e-04) (187903,9.063e-05) (553983,3.053e-05)
};

\addplot coordinates{
(13,5.755e-02)     (97,2.925e-02)     (545,1.351e-02)
(2561,5.842e-03)   (10625,2.397e-03)  (40193,9.414e-04)
(141569,3.564e-04) (471041,1.308e-04)
(1496065,4.670e-05)
};
}%


\bookmaketitle
\tableofcontents

\chapter{简介}

本宏包提供了方便地生成图和带标签坐标轴的工具。它可以绘制二维和三维的普通图、
对数图以及半对数图。坐标轴刻度、标签、图例（在叠加多幅图时）都可以通过键值
选项来添加。它支持线图、散点图、分段常数图、柱状图、面积图、网格图与曲面图、
补丁图、等高线图、矢量场图、直方图、箱线图、极坐标轴、三元图、史密斯图，
以及更多类型。它还可以循环使用一组预定义的线/标记/颜色规格。

总之，其目的是简化高质量函数图和/或数据图的生成，并解决以下问题：
%
\begin{itemize}
    \item 文档、字体类型与字号的一致性，
    \item 在坐标轴描述中直接使用 \TeX{} 数学模式，
    \item 数据与图的一致性（无需第三方工具），
    \item 通过导言区配置和样式实现文档之间的
        一致性。
\end{itemize}
%
尽管并非必需，也可以使用作为 \Tikz{} 一部分而开发的 |external| 库来生成
独立的 |.pdf| 或 |.eps| 图形。

欢迎使用 \PGFPlots{} 对二维和三维的中等规模数据集进行可视化。

它基于 Till Tantau 开发的 \PGF{}/\Tikz{} 宏包。


\chapter[关于 PGFPlots：预备知识]{关于 {\normalfont\PGFPlots{}}：预备知识}

本节包含关于升级、开发团队、安装（如果你需要手动安装的话）以及
故障排查的信息。除了升级说明之外，你可以完全跳过本节。

\PGFPlots{} 完全构建在 \Tikz{}/\PGF{} 之上。了解 \Tikz{} 会简化使用
\PGFPlots{} 的工作，尽管这并非必需。

\section{组成部分}

\PGFPlots{} 由两个组件构成：
%
\begin{enumerate}
    \item 绘图组件（即你当前正在阅读的部分），以及
    \item \PGFPlotstable{} 组件，它简化数字表格的数字格式化
        与后处理。它作为一个独立的
        宏包提供，并拥有自己的手册
        \href{file:pgfplotstable.pdf}{pgfplotstable.pdf}。
\end{enumerate}


\section{升级说明}

本版本提供了大量改进，感兴趣的读者可以在 \texttt{ChangeLog} 中找到
全部细节。不过，以下几项变更值得特别注意。

\PGFPlots{} 中一个共通的设置是键 |compat|：强烈建议始终把它写进你的
|.tex| 文件。虽然这个键给最终用户带来了一些工作，但它也解决了一个常见需求：
它确保你的 |.tex| 文件总是产生相同的输出，即使你安装了新版本的
\PGFPlots{}。另一方面，它让我们作为维护者能够修复软件缺陷并引入行为变更，
前提是这些变更只影响具有合适 |compat| 级别的文档。|compat| 键的确切影响、
可选取值及其含义将在以下各节中说明。


\subsection{新增可选功能}

\PGFPlots{} 在编写时就考虑了向后兼容：旧的 \TeX{}
文件应当无需修改即可编译，外观也不发生变化。不过，新功能偶尔会导致不同的行为。
在这种情况下，\PGFPlots{} 会停用新功能。\footnote{如果出现破坏向后兼容的情况，
我们深表歉意——并请你提交 bug 报告。我们会负责处理。}

任何导致向后兼容问题的新功能或错误修复都需要\emph{手动}且\emph{显式地}
启用。为此，你应当使用
%
\input pgfplots.preliminaries.compatcurrent.tmp

%
\noindent 放在你的导言区中。这将配置兼容层。

你至少应使用 |compat=1.3|。建议的取值会在运行 \TeX{}
之后打印到 |.log| 文件中。

以下是一份变更清单，除非使用合适的
|compat| 级别，否则这些变更处于停用状态：
%
\begin{enumerate}
    \item \PGFPlots{} 1.18 和 1.17 没有额外约束，在
        兼容级别方面与 1.16 相同。
    \item \PGFPlots{} 1.16 没有额外约束，在
        兼容级别方面与 1.15 相同。
    \item \PGFPlots{} 1.15 启用 |3d log sampling|，并修复
        柱状图的 |clip limits| 问题。
    \item \PGFPlots{} 1.14 改变了系统处理非均匀色图的方式，
        并启用高级 |colormap| 操作（参见
        |of colormap|）。
    \item \PGFPlots{} 1.13 修复极坐标轴中的坐标轴标签，并确保
        |shader=flat| 所选颜色与 |z buffer| 和
        |mesh/ordering| 无关。此外，它为所有堆叠柱状图启用
        |stack negative=separate|。保留较旧的兼容级别是为了
        保留最终用户的临时变通方案。
    \item \PGFPlots{} 1.12 启用 |lua backend| 并定义
        |boxplot/estimator=Excel|。
    \item \PGFPlots{} 1.11 改变了 |axis cs|：它现在是默认
        坐标系。如果你写 |\draw (1,2) -- (2,2);|，\PGFPlots{}
        会自动把它视为
        |\draw (axis cs:1,2) -- (axis cs:2,2);|。
    \item \PGFPlots{} 1.10 在兼容性方面与 1.9 没有区别。
    \item \PGFPlots{} 1.9 附带一个预设，用于组合 |ybar stacked| 和
        |nodes near coords|。此外，它会抑制堆叠柱状图中的空增量。
        要启用这个新预设，你必须
        使用 |compat=1.9| 或更高版本。
    \item \PGFPlots{} 1.8 对标签与刻度比例标签的对齐方式进行了
        新的修订。此外，它改进了
        |hide axis| 的包围盒。该修订通过 |compat=1.8| 或
        更高版本启用。

        配置 |compat=1.8| 是\emph{修复}
        三维坐标轴中 |axis lines=center| 的\emph{必要}条件。
    \item \PGFPlots{} 1.7 新增了以坐标轴单位定义柱宽的选项。
        这些通过 |compat=1.7| 或更高版本启用。
    \item \PGFPlots{} 1.6 新增了用于更精确缩放的选项，以及
        针对 |\addplot3 graphics| 的更多缩放选项。这些通过
        |compat=1.6| 或更高版本启用。
    \item \PGFPlots{} 1.5.1 将圆和椭圆的半径解释为
        \PGFPlots{} 坐标（旧版本使用 \pgfname{} 单位向量，
        它们与 \PGFPlots 没有直接关系）。换句话说：从 1.5.1 版开始，
        可以在坐标轴内部编写 |\draw circle[radius=5]|。这需要
        |\pgfplotsset{compat=1.5.1}| 或更高版本。

        若没有此兼容设置，圆和椭圆将像早期版本那样使用
        \pgfname{} 的底层画布单位。
    \item \PGFPlots{} 1.5 默认使用 |log origin=0|（这会影响
        对数柱状图或堆叠对数图）。较旧的版本
        保持 |log origin=infty|。这需要 |\pgfplotsset{compat=1.5}| 或
        更高版本。
    \item \PGFPlots{} 1.4 修复了若干较小的 bug，
        与早期版本相比可能产生约 $1$--$2\text{pt}$ 的差异。
        这需要 |\pgfplotsset{compat=1.4}| 或更高版本。
    \item \PGFPlots{} 1.3 带来了用户界面上的改进。
        旧版本中“行为选项”与“样式
        选项”之间的技术区分不再必要（不过仍
        完全支持）。

        这总是启用的。
    \item \PGFPlots{} 1.3 有一个新功能，允许自动\emph{把坐标轴标签
        紧贴刻度标签移动}。强烈推荐使用。它需要
        |\pgfplotsset{compat=1.3}| 或更高版本。

        由于这会影响间距，默认情况下并不启用。
    \item \PGFPlots{} 1.3 支持反向坐标轴。不再需要
        使用带负单位的变通方案。\pgfkeys{/pdflinks/search key
        prefixes in/.add={/pgfplots/,}{}}

        请查看 |x dir=reverse| 键。

        现有的变通方案仍能正常工作。使用
        |\pgfplotsset{compat=1.3}| 或更高版本，与 |x dir=reverse| 一起
        切换到新版本。
\end{enumerate}


\subsection{可能需要注意的旧功能}

\begin{enumerate}
    \item |scatter/classes| 功能从 1.3 版起会生成恰当的图例。
        这可能会改变已有使用 |scatter/classes| 的图的图例
        外观。
    \item 从 \PGFPlots{} 1.1 开始，\emph{不应再使用} |\tikzstyle|
        来设置 \PGFPlots{} 选项。

        尽管 |\tikzstyle| 仍被某些较旧的 \PGFPlots{}
        选项支持，但你应当把任何出现的 |\tikzstyle| 替换为
        |\pgfplotsset{|\meta{样式名称}|/.style={|\meta{键值列表}|}}|
        或相应的 |/.append style| 变体。更多细节参见
        Section~\ref{sec:styles}。
\end{enumerate}
%
对于这些变更带来的任何不便，我深表歉意。

\begin{pgfplotskey}{%
    compat=\mchoice{
        1.18,
        1.17,
        1.16,
        1.15,
        1.14,
        1.13,
        1.12,
        1.11,
        1.10,
        1.9,
        1.8,
        1.7,
        1.6,
        1.5.1,
        1.5,
        1.4,
        1.3,
        pre 1.3,
        default
    } (initially default)}
    导言区配置
    %
\input pgfplots.preliminaries.compatcurrent.tmp
    %
    允许在向后兼容与最新功能之间进行选择。
    这个键被设计为文档中遇到的第一个 \PGFPlots{} 键，因为它会准备
    全局选项。

    有时，你可能想在同一
    文档中使用不同的版本。这时，提供
    %
\begin{codeexample}[code only]
\begin{figure}
    \pgfplotsset{compat=1.4}
    ...
    \caption{...}
\end{figure}
\end{codeexample}
    %
    \noindent 以便把兼容设置限制在当前
    上下文中（本例中即 |figure| 环境）。

    查看你的 |.log| 文件输出，即可得到 |compat| 的建议取值。

    使用 |\pgfplotsset{compat=default}| 恢复出厂设置。

    尽管通常没有必要，也可以只启用
    选定的变更，而在总体上保持与旧版本的兼容：
    %
    \begin{pgfplotskeylist}{%
        compat/path replacement=\meta{version},
        compat/labels=\meta{version},
        compat/scaling=\meta{version},
        compat/scale mode=\meta{version},
        compat/empty line=\meta{version},
        compat/plot3graphics=\meta{version},
        compat/bar nodes=\meta{version},
        compat/BB=\meta{version},
        compat/bar width by units=\meta{version},
        compat/pgfpoint substitution=\meta{version},
        compat/general=\meta{version}
    }
        假设我们有一个使用 |\pgfplotsset{compat=1.3}|
        的文档，并且你想保持这种方式。

        此外，你发现 1.5.1 版支持圆和
        椭圆。这时，使用
\begin{codeexample}[]
% preamble:
\pgfplotsset{
    compat=1.3,
    compat/path replacement=1.5.1,
}
\begin{tikzpicture}
\begin{axis}[
    extra x ticks={-2,2},
    extra y ticks={-2,2},
    extra tick style={grid=major}]
    \addplot {x};
    \draw (axis cs:0,0) circle[radius=2];
\end{axis}
\end{tikzpicture}
\end{codeexample}

        所有这些键都接受 |compat| 键的可选值。

        |compat/path replacement| 键控制圆和
        椭圆半径的解释方式。

        |compat/labels| 键控制坐标轴标签的对齐方式：要么
        紧邻刻度，要么使用绝对偏移。从 |1.8| 起，它还
        启用了一套全新的坐标轴标签样式修订。在大多数
        情况下，你看不出区别——但它修复了
        三维坐标轴中的 |axis lines=center|。

        |compat/scaling| 键控制 1.4 和 1.6 版中引入的
        一些 bug 修复：为了提高精度，它们可能带来轻微的
        缩放差异。

        |compat/plot3graphics| 控制
        |\addplot3 graphics| 的新功能。

        |compat/scale mode| 允许启用/禁用警告“The
        content of your 3d axis has CHANGED compared to previous versions”，
        因为 |axis equal| 和 |unit vector ratio| 功能在
        1.6 之前的所有版本中都是坏的，并在 1.6 中得到了修复。

        |compat/empty line| 允许在输入文件中写入空行以
        产生一次跳跃。这需要 |compat=1.4| 或更新版本。详见
        |empty line|。

        |compat/BB| 改变包围盒，使其即使在
        |hide axis| 的情况下也紧凑贴合。

        |compat/bar width by units| 允许用坐标轴单位来表示
        |bar width=1|（即以坐标轴单位计）。

        |compat/bar nodes| 启用 |ybar stacked| 和
        |nodes near coords| 的预设。此外，它为
        堆叠柱状图启用 |stacked ignores zero|。

        |compat/general| 键控制 |log origin|、|lua backend|、
        |enable tick line clipping| 和 |boxplot/estimator|。

        |compat/pgfpoint substitution| 键决定 |axis cs| 是否为
        默认坐标系（自 1.11 起）。

        详细效果可参见本节开头。
    \end{pgfplotskeylist}

    取值 \meta{version} 可以是 |default|、一个版本号，或 |newest|。
    取值 |default| 与 |pre 1.3| 相同（除无关紧要的差异外）。
    强烈\emph{不建议}使用 |newest|：它可能根据 \PGFPlots{}
    的当前版本而改变你的文档。请查看你的 |.log| 文件，
    以获得关于最佳版本的建议。
\end{pgfplotskey}


\section{开发团队}

\PGFPlots{} 主要由 Christian Feuersänger 编写，Pascal Wolkotte 和
Nick Papior Andersen 作为业余项目贡献了许多改进。我们希望它有用，
并能提供有价值的图。

如果你有兴趣编写一些东西但不知从何入手，可以参阅
随 \PGFPlots{} 提供的辅助手册
\href{file:TeX-programming-notes.pdf}{TeX-programming-notes.pdf}。它
远非完整，但也许是一个好的
起点（至少能引导你找到更多文献）。


\section{致谢}

感谢上帝赐予我享受编程的无数时光。我感谢 Pascal Wolkotte 和
Nick Papior Andersen 作为开发团队一员所付出的编程努力与贡献。我感谢 Francesco Poli 为 |contour lua| 实现 |lua| 算法。
我感谢 Jürnjakob Dugge 贡献了
|hist/density|、用于 \verbpdfref{\addplot3} |graphics| 的 matlab 脚本、出色的
用户论坛帮助以及有价值的 bug 报告。我感谢 Stefan Tibus，他贡献了
|\addplot shell| 功能。我感谢 Tom Cashman 贡献了
|reverse legend| 功能。特别感谢 Stefan Pinnow 持续
测试 \PGFPlots{}、讨论需求、提出功能请求和
bug 修复（从而带来众多质量改进），以及移植并
集成 |colorbrewer| 库。此外，我感谢 Schweitzer 教授
富有成效的讨论，以及他最初鼓励我开始这样一个
宏包。感谢 Meine 博士的想法和建议。特别感谢
Markus Böhning 校对了 \PGF{}、\PGFPlots{} 和
\PGFPlotstable{} 的所有手册。感谢 Vincent A. Traag 把 |colorbrewer| 配色
带给 \PGFPlots{}。也感谢众多国际贡献者，他们
提出了功能请求、发现了 bug，或仅仅是改进了
手册！

最后但同样重要的是，我感谢 Till Tantau 和 Mark Wibrow 打造了出色的
图形（及更多）宏包 \PGF{} 和 \Tikz{}，它们正是
\PGFPlots{} 的基础。


\input{pgfplots.install}

% main=pgfplots.tex

\chapter{循序渐进教程}
{%
\tikzset{external/figure name/.add={}{tutorials_}}%


\section{简介}

为了分析和传达研究成果、论文结果，或者仅仅是某种结果，数据可视化往往
是必要且方便的。

\PGFPlots{} 是一个可视化工具。设计 \PGFPlots{} 的动机在于：作为最终用户，你
只需提供数据和描述作为输入，而 \PGFPlots{} 会处理其余一切，例如选择合适的
缩放因子、缩放到指定的目标尺寸、选择合适的显示范围、分配刻度位置、绘制
坐标轴并把描述文字放置在合适的位置。

\PGFPlots{} 解决了 \LaTeX{} 中一个长期存在的可视化问题：它的
描述文字使用与嵌入正文相同的字体，字号也完全相同。它直接嵌入 \LaTeX{}，
使得尽可能方便地使用 \LaTeX{} 强大的数学模式：从各种坐标轴描述一直到
用户自定义的标注。它提供文档级统一的线型、配色方案、标记\ldots{}
所有这些共同保证了一致性。

\PGFPlots{} 提供高质量的输出。同时，它是一种嵌入式解决方案：
它在很大程度上独立于第三方工具，不过也提供导入功能，以利用现有的工具。

它的主要目标是：你提供数据和描述——然后
\PGFPlots{} 无需更多输入即可运行。如果你愿意，也可以按自己的喜好进行
定制。


\input{pgfplots.tutorial1.tex}
\input{pgfplots.tutorial2.tex}
\input{pgfplots.tutorial3.tex}
\input{pgfplots.tutorial4.tex}
}%

% main=manual.tex

\chapter{参考手册}
\label{cha:pgfplots:reference}

{
\tikzset{external/figure name/.add={}{reference_}}

\input{pgfplots.reference.texdialects.tex}
\input{pgfplots.reference.axis-addplot.tex}
\input{pgfplots.reference.preliminaryoptions.tex}
\input{pgfplots.reference.2dplots.tex}
\input{pgfplots.reference.3dplots.tex}
\input{pgfplots.reference.markers-meta.tex}
\input{pgfplots.reference.axisdescription.tex}
\input{pgfplots.reference.scaling.tex}
\input{pgfplots.reference.3dconfiguration.tex}
\input{pgfplots.reference.errorbars.tex}
\input{pgfplots.reference.numberformatting.tex}
\input{pgfplots.reference.specifyrange.tex}
\input{pgfplots.reference.tickoptions.tex}
\input{pgfplots.reference.gridoptions-axiscoordinates.tex}
\input{pgfplots.reference.styleoptions.tex}
\input{pgfplots.reference.alignment.tex}
\input{pgfplots.reference.bb-clip.tex}
\input{pgfplots.reference.symbolic-transformations.tex}
\input{pgfplots.reference.coordfiltering.tex}
\input{pgfplots.reference.transformations.tex}
\input{pgfplots.reference.linefitting.tex}
\input{pgfplots.reference.miscellaneous.tex}
\input{pgfplots.reference.tikzinteroperability.tex}
\input{pgfplots.reference.layers.tex}
\input{pgfplots.reference.technicalinternals.tex}

}


\chapter{相关库}

本章介绍一些随 \PGFPlots{} 提供的库，但它们多少有些特殊，需要单独
激活。
\pgfmanualpdflabel{\textbackslash usepgfplotslibrary}{}

% \input{pgfplots.libs.statistics.tex}
% \endinput

\input{pgfplots.libs.clickable.tex}
\input{pgfplots.libs.colorbrewer.tex}
\input{pgfplots.libs.colormaps.tex}
%\input{pgfplots.libs.colortol.tex}         % not ready yet
\input{pgfplots.libs.dateplot.tex}
\input{pgfplots.libs.decorations.softclip.tex}
\input{pgfplots.libs.external.tex}
\input{pgfplots.libs.fillbetween.tex}
%\section{FIXME} \endinput
\input{pgfplots.libs.groupplots.tex}
\input{pgfplots.libs.patchplots.tex}
\input{pgfplots.libs.polar.tex}
\input{pgfplots.libs.smithchart.tex}
\input{pgfplots.libs.statistics.tex}
\input{pgfplots.libs.ternary.tex}
\input{pgfplots.libs.units.tex}


\chapter{内存与速度考量}

{
\tikzset{external/figure name/.add={}{memspeed_}}

\section{\TeX{} 的内存上限}
\label{sec:pgfplots:optimization}

只要正确配置 \TeX{} 的内存上限，\PGFPlots{} 就能排版包含数千个数据点的图。其运行时间大致与输入点数成正比。\footnote{事实上，运行时间是伪线性的：当点数大约达到 $100{,}000$ 时，会退化为二次增长。这一限制同样适用于 \PGF{} 路径的路径长度。此外，对堆叠图（stacked plots）而言，线性运行时间目前尚不可实现。}
%
\pgfplotsexpensiveexample
\begin{codeexample}[]
\begin{tikzpicture}
\begin{axis}[
    enlargelimits=0.01,
    title style={yshift=5pt},
    title=Scatter plot with $2250$ points,
]
    \addplot [
        blue,
        mark=*,
        only marks,
        mark options={scale=0.3},
    ] file [skip first]
        {plotdata/pgfplots_scatterdata3.dat};
\end{axis}
\end{tikzpicture}
\end{codeexample}

\pgfplotsexpensiveexample
\begin{codeexample}[]
\begin{tikzpicture}
\begin{axis}[
    enlarge x limits=0.03,
    title=Ornstein--Uhlenbeck sample
        ($13000$ time steps),
    xlabel=$t$,
]
    \addplot [blue] file {plotdata/ou.dat};
\end{axis}
\end{tikzpicture}
\end{codeexample}

\pgfplotsexpensiveexample
\begin{codeexample}[]
% huger graphs are possible; consider lualatex
\begin{tikzpicture}
\begin{axis}[
    title=$80 \times 80$ Smooth Surface,
    xlabel=$x$,
    ylabel=$y$,
]
  \addplot3 [surf,samples=80,shader=interp,domain=0:1]
        {sin(deg(8*pi*x))* exp(-20*(y-0.5)^2)
        + exp(-(x-0.5)^2*30
            - (y-0.25)^2 - (x-0.5)*(y-0.25))};
\end{axis}
\end{tikzpicture}
\end{codeexample}

\PGFPlots{} 完全依赖 \TeX{} 完成所有排版工作。它使用 \PGF{} 的前端层（front-end layer）和基础层（basic layer）来执行所有绘图操作。对于复杂的图，这可能需要花费一些时间，你可以参阅第~\ref{cha:pgfplots:importexport} 章，了解如何将单个图导出到外部图形文件。外部化（externalization）是减少排版时间的最佳方式。

然而，对于包含大量数据点的大规模图，很容易触及 \TeX{} 自身能力的上限。


\section{内存限制}

大多数 \TeX{} 发行版的默认设置都相当保守，因此可能需要调整它们。

通常，日志文件或最终的错误消息会包含一段关于已使用资源的摘要，从而提示需要增大哪一个参数。


\subsection{Lua\LaTeX{}}

一个效果相当不错的解决方案是切换 \LaTeX{} 可执行程序：如果你使用的是一个较完善的 \TeX{} 发行版，那么它应该已经附带了 |lualatex| 可执行程序。它采用动态内存分配，因此通常对任何 \PGFPlots{} 坐标轴都有足够的内存可用。

Lua\LaTeX{} 可执行程序 |lualatex| 与 |pdflatex| 几乎完全兼容。

这种方法适用于任何平台。


\subsection{MiK\TeX{}}

如果你正在使用 MiK\TeX{}，并且不想（或不能）切换到 |lualatex|，可以按如下步骤操作。

对于 MiK\TeX{}，可以通过两种方式增大内存上限。第一种是使用命令行开关：
%
\begin{codeexample}[code only]
pdflatex
    --stack-size=n --save-size=n
    --main-memory=n --extra-mem-top=n --extra-mem-bot=n
    --pool-size=n --max-strings=n
\end{codeexample}
%
\noindent 如果 MiK\TeX{} 出现内存不足，请调整这些设置。通常，我们不会手动调用 |pdflatex|：有一个开发辅助工具会负责所有调用，因此需要调整的是它。

有时，永久修改 MiK\TeX{} 配置文件可能更好，例如可以避免每次重新配置 \TeX{} 开发程序。这可以通过下面的命令实现
%
\begin{codeexample}[code only]
initexmf --edit-config-file=pdflatex
\end{codeexample}
%
\noindent 该命令既可以在 Windows 的命令提示符中输入，也可以通过「开始 $\gg$ 运行」执行。执行后，系统会用编辑器打开相应的配置文件。一个示例配置文件可以是
%
\begin{codeexample}[code only]
main_memory=90000000
save_size=80000
\end{codeexample}
%
或者也可以输入下面列出的任意一项配置。感谢 ``LeSpocky'' 提供的文档：

\url{http://blog.antiblau.de/2009/04/21/speicherlimits-von-miktex-erhoehen}.


\subsection{\TeX{}Live 或类似发行版}

除了切换到 |lualatex| 之外，你还可以按如下方式操作，以保留现有的 |dvips| 或 |pdflatex| 工作流。

对于 Unix 安装，需要调整配置文件。可以按如下步骤进行：
%
\begin{enumerate}
    \item 在你的系统上找到 |texmf.cnf|。在我的 Ubuntu 安装中，它位于

        |/usr/share/texmf/web2c/texmf.cnf|。
    \item 可以直接修改 |texmf.cnf|，也可以把它复制到某个方便的位置。如果你选择复制，步骤如下：
        %
        \begin{itemize}
            \item 在你本地的副本中只保留被修改的条目，以减少冲突。\TeX{} 总是会读取其搜索路径中找到的 \emph{所有} 配置文件。
            \item 调整搜索路径，使其能找到你的本地副本。这可以通过环境变量 |TEXMFCNF| 来完成。假设你的本地副本位于 |~/texmf/mytexcnf/texmf.cnf|，则可以写
                %
\begin{codeexample}[code only]
export TEXMFCNF=~/texmf/mytexcnf:
\end{codeexample}
                %
                这样就会先在你的目录中搜索，然后再搜索所有其他系统目录。
        \end{itemize}
    \item 你应当修改以下条目
        %
\begin{codeexample}[code only]
main_memory = n
extra_mem_top = n
extra_mem_bot = n
max_strings = n
param_size = n
save_size = n
stack_size = n
\end{codeexample}
        %
        日志文件通常会包含关于需要增大哪一个参数的信息。
\end{enumerate}
%
下面给出一个此类配置文件的示例。它会修改内存上限。
%
\begin{enumerate}
    \item 创建文件 |~/texmf/mytexcnf/texmf.cnf|（如必要，也一并创建相应的路径）。
        %
\begin{codeexample}[code only]
% newly created file ~/texmf/mytexcnf/texmf.cnf:
% If you want to change some of these sizes only for a certain TeX
% variant, the usual dot notation works, e.g.,
% main_memory.hugetex = 20000000
main_memory = 230000000 % words of inimemory available; also applies to inimf&mp
extra_mem_top = 10000000     % extra high memory for chars, tokens, etc.
extra_mem_bot = 10000000     % extra low memory for boxes, glue, breakpoints, etc.
save_size = 150000    % for saving values outside current group
stack_size = 150000    % simultaneous input sources

% Max number of characters in all strings, including all error messages,
% help texts, font names, control sequences.  These values apply to TeX and MP.
%pool_size = 1250000
% Minimum pool space after TeX/MP's own strings; must be at least
% 25000 less than pool_size, but doesn't need to be nearly that large.
%string_vacancies = 90000
% Maximum number of strings.
%max_strings = 100000
% min pool space left after loading .fmt
%pool_free = 47500
\end{codeexample}
        %
    \item 运行 |texhash|，让 \TeX{} 更新它的 |~/texmf/ls-R| 数据库。
    \item 创建环境变量 |TEXMFCNF|，并将其值设为 `|~/texmf/mytexcnf:|'（注意末尾的 `|:|'！）。在我的 Linux 系统上，可以通过向 |~/.bashrc| 添加
        %
\begin{codeexample}[code only]
export TEXMFCNF=~/texmf/mytexcnf:
\end{codeexample}
        %
        来实现。
\end{enumerate}

遗憾的是，\TeX{} 不允许任意设置内存上限：可执行程序中硬编码了一个上界。


\section{减少排版时间}

\PGFPlots{} 会进行大量计算，从抽象的坐标计算一直到由 \PGF{} 实现的底层 |.pdf| 绘图命令。对于复杂的图，这可能需要相当长的时间——尤其是 3D 图。


\subsection{LUA}

如果你使用 |compat=1.12|（或更新版本），并通过 |lualatex| 编译文档，\PGFPlots{} 就会启用其 |lua backend|。这一开关能减少生成输出文件所需的时间，对 3D 图尤其明显。

\begin{pgfplotskey}{lua backend=\mchoice{true,false} (initially true)}
    如果 |lua backend| 处于活动状态，并且文档是通过 |lualatex| 处理的，\PGFPlots{} 就会启用可扩展性与性能改进；它产生的输出与不启用时相同，但耗时更少，并且具有更智能的内存管理。

    该特性依赖于用快速脚本语言 \lua 对 \PGFPlots{} 的部分重新实现。要从中受益，你需要
    %
    \begin{enumerate}
        \item 在导言区写入 |compat=1.12|（或更新版本），并且
        \item 使用 |lualatex| 来编译你的 |.tex| 文件（或者至少编译那些包含 \PGFPlots{} 图的文件）。
    \end{enumerate}

    使用该特性后，生成曲面图所需的时间可以缩减到纯 \TeX{} 实现（即 |pdflatex|）所需时间的 25\%--50\%。其他类型的图同样会从中受益。

    |lua backend| 以两种方式工作：第一，用更快的 \lua 对应实现替换若干独立的子程序。这相对通用。第二，用等价的 \lua 实现替换整个处理步骤。%
    %
    \footnote{
        第一种方式的一个例子是：每当 \PGFPlots{} 在其 |colormap| 中进行查找时，它就把控制权转交给 |lua backend|，并将结果立即传回给 \TeX{}。主控制流仍然位于较慢的 \TeX{} 实现中。第二种方式的一个例子是 \texttt{\textbackslash addplot expression}：\PGFPlots{} 会把数学表达式以及所有相关的输入参数（包括 \texttt{samples} 和 \texttt{domain}）拷贝到 \texttt{lua backend}。然后，\texttt{lua backend} 会执行所有循环并收集坐标，最终把这些坐标交给 \pgfname{} 的（\TeX{}）实现。这样，控制流就位于快速的 \lua{} 解释器中。
    }
    \PGFPlots{} 1.12 为重要且耗时的操作提供了 \lua{} 实现。只要可能，就会使用这一实现。然而，它目前只覆盖了部分功能：\PGFPlots{} 的某些特性尚不被 \lua{} 后端支持。在这种情况下，它会自动关闭，并回退使用 \TeX{} 实现。

    |lua backend| 还会继续扩展，以获得更多性能提升并覆盖更多特性。如果你想查看 |lua backend| 是否在你的图中被（去）激活，请参考 |lua debug|。

    请注意，|lua backend| 使用了一个不同的、精度更高的数学引擎。因此，与 |pdflatex| 生成的输出相比，图的外观可能会出现一些无关紧要的像素差异。一般建议固定使用其中一种方式来生成文档，即要么使用 |lualatex|，\emph{要么}使用 |pdflatex|。由于 |lua backend| 的存在，从一种方式迁移到另一种方式可能会改变你文档的外观。

    最终，\PGFPlots{} 可能会加入一些计算开销较大的特性，这些特性只能作为 |lua backend| 的一部分实现，没有它将不可用。

    请注意，|lua backend| 需要一个至少支持 \lua~5.2 的 \TeX{} 发行版（例如 \TeX{}Live 2014）。如果你的 |lualatex| 附带的是更老版本的 \lua{} 解释器，它会自动停用。

    下面是一些如何从 |lua backend| 受益的指南：
    %
    \begin{itemize}
        \item 涉及宏定义的数学表达式在 |lua backend| 中不可用。最佳实践：优先使用 |declare function| 而不是宏常量：
            %
\begin{codeexample}[]
\begin{tikzpicture}
\begin{axis}[
    small,
    title=Do,
    samples=2,
]
    \addplot+ [
        declare function={C=4;}
    ] {C*x};
\end{axis}
\end{tikzpicture}
\end{codeexample}

\begin{codeexample}[]
\begin{tikzpicture}
\begin{axis}[
    small,
    title=Don't,
    samples=2,
]
        \def\constantC{4}
    \addplot {\constantC*x};
\end{axis}
\end{tikzpicture}
\end{codeexample}

            两者在语义上是等价的，但由于 \lua{} 无法解释 \TeX{} 宏，它会拒绝处理「Don't」这一情形。「Don't」情形会产生一条日志消息
            %
\begin{verbatim}
Package pgfplots info on input line 16: Deactivating LUA version of
plot expression for plot 0 (type 'pgfplothandlerlineto'):
y expression '\constantC *x' contains a TeX macro.
\end{verbatim}
            %
        \item 任何需要「原生」\TeX{} 代码的操作在 |lua backend| 中都不可用。这包括 |x filter/.code={...}|，因为 |/.code| 无法映射到 \lua。最佳实践：优先使用 |x filter/.expression| 而不是 |x filter/.code|：
            %
\begin{codeexample}[]
\begin{tikzpicture}
    \begin{axis}[small,grid=major,title=Do]
    \addplot+ [
        unbounded coords=jump,
        y filter/.expression={y==3 ? nan : y},
    ] table {
        x   y
        1   2
        1.5 2.5
        2   3
        2.5 3.5
        3   4
    };
    \end{axis}
\end{tikzpicture}
\end{codeexample}
            %
        \item 绘图表达式可以完全在 \lua 中处理；其他所有输入坐标类型仍由 \TeX{} 读取数值后再交给 \lua。
        \item 还有若干操作计划在未来版本中加入 |lua backend| 支持。
    \end{itemize}
\end{pgfplotskey}

\begin{pgfplotskeylist}{%
    lua debug,
    lua debug=\mchoice{false,off,off and silent,verbose,compileerror}%
}
    通常，|lua backend| 是静默工作的：\PGFPlots{} 自行决定它是否处于活动状态、是否支持当前操作，并据此继续执行。如果两个条件都满足，就会把控制权转交给 |lua backend|。如果其中一个不满足，它会自动回退到旧的 \TeX{} 实现。因此，你无法立即感知到它的存在，只有编译图片所需的时间会有所不同。

    该键用于控制：当 |lua backend| 处于活动状态、但由于遇到了不支持的特性而无法使用时，\PGFPlots{} 应如何处理。

    选项 \declaretext{false} 和 \declaretext{off} 会关闭 \lua{} 调试。在这种情况下，\PGFPlots{} \emph{仅}在使用了 \TeX{} 实现而不是 |lua backend| 时才写入日志消息。此外，它只会把消息写入 |.log| 文件（而不会写入控制台输出）。

    选项 \declaretext{off and silent} 既不写信息消息，也不写调试消息：它与 |off| 相同，只是在使用了 \TeX{} 实现时不再写日志消息。这是信息量最少的选项。

    选项 \declaretext{verbose} 还会把成功消息写入 |.log| 文件（不写入控制台）。

    选项 \declaretext{compileerror} 会在 |lua backend| 处于活动状态、但不支持当前图或绘图时中止编译。

    请注意，\PGFPlots{} 可能附带一些开销较大的受限操作，这些操作只能在 |lua backend| 中使用。如果由于某种原因无法使用 |lua backend|，这些条目将导致编译错误。
\end{pgfplotskeylist}


\subsection{每个图只编译一次}

减少排版时间的一种方法是：让 \PGF{} 在一次运行中为一部分（或全部）图形生成独立的临时 |.pdf|（或 |.eps|）文档，并在后续运行中复用这些临时图片。对于 \PGFPlots{} 而言，这是减少较大文档排版时间最有效的方式。可以通过第~\ref{sec:pgfplots:export} 节中描述的 |external| 库来实现。

}


\chapter{从其他格式导入/导出}
\label{cha:pgfplots:importexport}

{
\tikzset{external/figure name/.add={}{importexport_}}

本章介绍如何把单个图导出为独立的 \pdf{} 图形文件（或 \eps{} 图形文件）。此外，还介绍了一个 Matlab$^\text{\textregistered}$ 脚本，可用于把 Matlab 图转换为 \PGFPlots{}。


\section[导出为 PDF/EPS]{导出为 {\normalfont\pdf{}/\eps{}}}
\label{sec:pgfplots:export}

可以使用 \pgfname{} 和/或 \Tikz{} 的例程，把图导出为独立的 \pdf{} 文档。


\subsection{使用 \Tikz 的自动外部化框架}

\begin{pgfplotslibrary}{external}
\pgfkeys{/pdflinks/search key prefixes in/.add={/tikz/external/,}{}}
    |external| 库提供了一种便捷的方法，把每个独立的 |tikzpicture| 导出为单独的~|.pdf|（或~|.eps|）。后续的 \LaTeX{} 运行将直接包含这些图形，从而显著减少排版时间。

    该库也可以用于把文档提交给那些甚至没有安装 \PGFPlots{} 或 \Tikz{} 的作者。


    \paragraph{技术前言：}

    |external| 库由 Christian Feuersänger（\PGFPlots 的作者）编写。它已作为通用库贡献给 \Tikz{}，因此完整的参考文档以及所有调节选项可参见~\cite[Section ``Externalization Library'']{tikz}。命令 |\usepgfplotslibrary{external}| 实际上只是一个包装，它加载 |\usetikzlibrary{external}|；如果因为所安装的 \pgfname{} 版本至多为 $2.00$ 而该库尚不存在，则加载 \PGFPlots{} 自带的一个副本。

    |external| 库的设计使得对文档本身\emph{无需做任何改动}。其思路如下：
    %
    \begin{enumerate}
        \item 每一个 |\begin{tikzpicture}| $\dotsc$ |\end{tikzpicture}| 都会得到一个文件名。文件名既可以用 |\tikzsetnextfilename|\marg{output file name} 手动指定，也可以自动生成；自动生成时使用 \meta{tex file name}|-figure|\meta{number}，其中 \meta{number} 递增。

        \item 该库通过 \TeX{} 的 write18 系统调用，自动以形如 |pdflatex --jobname |\meta{output file name} 的命令写出生成的图片。这与调用 |gnuplot| 所用的是同一套框架。
    \end{enumerate}
    %
    唯一必要的步骤是使用

    \pgfmanualpdflabel{\textbackslash tikzexternalize}{}%
    |\usepgfplotslibrary{external}|

    |\tikzexternalize|

    \noindent 在文档导言区的某处（与系统相关的配置选项见下文）。文档无需再做其他修改。假设我们有一个名为 |test.tex| 的文件：
    %
\begin{codeexample}[code only]
\documentclass{article}
\usepackage{pgfplots}
    \usepgfplotslibrary{external}
    \tikzexternalize% activate externalization!
\begin{document}
    \begin{figure}
        \begin{tikzpicture}
            \begin{axis}
                \addplot {x^2};
            \end{axis}
        \end{tikzpicture}
        \caption{Our first external graphics example}
    \end{figure}

    \begin{figure}
        \begin{tikzpicture}
            \begin{axis}
                \addplot {x^3};
            \end{axis}
        \end{tikzpicture}
        \caption{A second graphics}
    \end{figure}
\end{document}
\end{codeexample}
    %
    \noindent 要启用系统调用，我们输入
    %
\begin{codeexample}[code only]
pdflatex -shell-escape test
\end{codeexample}
    %
    \noindent 之后 \LaTeX{} 会自动生成所需的图形文件 |test-figure0.pdf| 和 |test-figure1.pdf|。之后每次调用 |pdflatex| 都会直接使用 |\includegraphics|，不再重新处理 |tikzpicture| 本身（对于 MiK\TeX{}，需要不同的命令行开关，参见 |shell escape| 选项）。

    如果某个图需要重新生成，只需删除全部或部分图形文件并重新生成即可。另一种方式是在文档的某处使用 |\tikzset{external/force remake}|，之后每一个图都会自动重新生成。

    外部化图的文件名有三种修改方式：
    %
    \begin{itemize}
        \item 使用 |prefix| 修改整体文件名前缀，
        \item 使用 |\tikzsetnextfilename| 修改单个图的文件名，
        \item 使用 |figure name| 修改一组图的文件名。
    \end{itemize}

    \begin{key}{/tikz/external/prefix=\marg{file name prefix} (initially empty)}
        |\tikzsetexternalprefix|\marg{file name prefix} 的快捷方式，见下文。
    \end{key}

    \begin{command}{\tikzsetexternalprefix\marg{file name prefix}}
        为所有文件名指定一个公共前缀。例如，
        %
\begin{codeexample}[code only]
\tikzsetexternalprefix{figures/}
\end{codeexample}
        %
        会把 |figures/| 加到每一个外部图形文件名之前。
    \end{command}

    \begin{command}{\tikzsetnextfilename\marg{file name}}
        为\emph{下一个} \tikzname{} 图或 |\tikz| 短命令设置文件名。它\emph{只}对下一个图生效。

        没有显式指定文件名的图会自动生成文件名。

        请注意，|prefix| 仍然会被加到 \meta{file name} 之前。
        %
\begin{codeexample}[code only]
\documentclass{article}
% main document, called main.tex
\usepackage{tikz}
    \usepgfplotslibrary{external}
    \tikzexternalize[prefix=figures/]% activate with a name prefix
\begin{document}
\tikzsetnextfilename{firstplot}
\begin{tikzpicture} % will be written to 'figures/firstplot.pdf'
    \begin{axis}
        \addplot {x};
    \end{axis}
\end{tikzpicture}

\begin{tikzpicture} % will be written to 'figures/main-figure0.pdf'
   \draw [help lines] (0,0) grid (5,5);
\end{tikzpicture}
\end{document}
\end{codeexample}

\begin{codeexample}[code only]
pdflatex -shell-escape main
\end{codeexample}
        %
    \end{command}

    \begin{key}{/tikz/external/figure name=\marg{name}}
        与 |\tikzsetfigurename|\marg{name} 相同。
    \end{key}

    \begin{command}{\tikzsetfigurename\marg{name}}
        修改\emph{之后所有}图的名称。在文档编写过程中，可以通过
        |\tikzset{external/figure name|=\marg{name}|}| 来修改 |figure name|。每一个不同的 \meta{name} 都会使用一个独立的计数器，每个计数器都从
        $0$ 开始。\footnote{这些计数器存储在不同的\emph{宏}中。换句话说：不需要使用 \TeX{} 寄存器。}

        |prefix| 的值会在 |figure name| 求值之后再应用。
        %
\begin{codeexample}[code only]
\documentclass{article}
% main document, called main.tex
\usepackage{tikz}
    \usepgfplotslibrary{external}
    \tikzexternalize% activate externalization!
\begin{document}
% will be written to 'main-figure0.pdf'
\begin{tikzpicture}
    \begin{semilogyaxis}
        \addplot {exp(x)};
    \end{semilogyaxis}
\end{tikzpicture}

{
    \tikzset{external/figure name={subset_}}
    A simple image is \tikz \fill (0,0) circle(5pt);. % will be written to 'subset_0.pdf'

    \begin{tikzpicture} % will be written to 'subset_1.pdf'
        \begin{axis}
            \addplot {x^2};
        \end{axis}
    \end{tikzpicture}
}% here, the old file name will be restored:

\begin{tikzpicture} % will be written to 'main-figure1.pdf'
    \begin{axis}
        \addplot[domain=1e-3:100] {1/x};
    \end{axis}
\end{tikzpicture}
\end{document}
\end{codeexample}
        %
        |figure name| 的作用域在下一个右花括号处结束（与 |\tikzset| 设置的所有值一样）。

            \medbreak
        备注：使用
        |\tikzset{external/figure name/.add=|\marg{prefix}\marg{suffix}|}| 可以在 |figure name| 的实际值前添加一个 \meta{prefix}、在末尾追加一个 \meta{suffix}。这对于类似下面这样的场景可能有用：
        %
\begin{codeexample}[code only]
\tikzset{external/figure name=main}

% uses main_0.pdf, main_1.pdf, ...

\chapter{The first section}
{\tikzset{external/figure name/.add={}{_firstsection}}
    ...
    % uses main_firstsection_0.pdf, main_firstsection_1.pdf, ...
}

\chapter{The second section}
{\tikzset{external/figure name/.add={}{secondsection_}}
    ...
    % uses main_secondsection_0.pdf, main_secondsection_1.pdf, ...
    \section{Second subsection}
    {\tikzset{external/figure name/.add={}{sub_}}
        ...
        % uses main_secondsection_sub_0.pdf, main_secondsection_sub_1.pdf, ...
    }
    % uses main_secondsection_2.pdf, main_secondsection_3.pdf, ...
}
\end{codeexample}
    \end{command}

    \begin{command}{\tikzappendtofigurename\marg{suffix}}
        在 |figure name| 的实际值之后追加 \meta{suffix}。

        它是
        |\tikzset{external/figure name/.add={}|\marg{suffix}|}| 的快捷方式（即使没有安装 \tikzname{} 也支持此快捷方式，见下文）。
    \end{command}


    \paragraph{\textsc{eps} 输出或 MiK\TeX{} 的配置选项：}

    由于 |external| 库通过系统调用来工作，因此需要根据本地系统进行调整。对于 MiK\TeX{} 这是必要的，因为它使用不同的选项来启用这些系统调用。对于 \textsc{eps} 输出也是必要的，因为这涉及到一组不同的工具。

    请注意，\emph{最重要的部分}是启用系统调用。通常的做法是用 |pdflatex -shell-escape| 或 |pdflatex -enable-write18|（MiK\TeX{}）来排版你的文档。这些选项\emph{需要在你的 \TeX{} 编辑器中配置}。除这一步之外，你可能还希望配置系统调用：

    \begin{key}{/tikz/external/system call=\marg{template}}
    \label{extlib:systemcall:option}
        用于生成系统调用的模板字符串。在 \meta{template} 内部，宏 |\image| 可以作为占位符，表示即将生成的图片；而 |\texsource| 包含主文件名（实际上它包含 |\input|\marg{main file name}，但这无关紧要）。

        默认值为
        %
\begin{codeexample}[code only]
\tikzset{external/system call={pdflatex \tikzexternalcheckshellescape -halt-on-error
    -interaction=batchmode -jobname "\image" "\texsource"}
\end{codeexample}
        %
        \noindent 其中 \declareandlabel{\tikzexternalcheckshellescape} 仅在当前文档是用 |-shell-escape| 排版时，才会插入配置键 |shell escape| 的值。\footnote{请注意，对于默认配置而言这总是成立的。这一安全考虑主要针对 \texttt{mode=list and make}，它\emph{不}需要 shell escape 也能工作。}

        对于 |eps| 输出，你可以（并且需要）使用
        %
\begin{codeexample}[code only]
\tikzset{external/system call={latex \tikzexternalcheckshellescape -halt-on-error
    -interaction=batchmode -jobname "\image" "\texsource" &&
    dvips -o "\image".ps "\image".dvi}}
\end{codeexample}

        参数 \meta{template} 会用 |\edef| 展开，因此所有控制序列都会被展开。在这一求值过程中，`|\\|' 会变成一个普通反斜杠 `|\|'。此外，如果某些宏包把双引号 `|"|'、单引号 `|'|'、分号和短横线 `|-|' 用作宏，它们也会被转换为普通字符。例如，这保证了与 |german| 宏包的兼容性。
    \end{key}

    \begin{key}{/tikz/external/shell escape=\marg{command-line arg} (initially -shell-escape)}
        包含启用 |\write18| 特性的 |latex| 命令行选项。

        对于 \TeX{}Live，是 |-shell-escape|。对于 MiK\TeX{}，系统使用 |-enable-write18|。
    \end{key}


    \paragraph{对外部文件中标签与引用的支持}

    |external| 库对外部文件中的 |\label| 和 |\ref|（以及其他通常会在 |.aux| 文件中保存信息的命令）提供了额外的支持。

    然而，当你尝试使用下面两种情况时，仍有几点需要注意：
    %
    \begin{enumerate}
        \item[a)] 在外部化图形中引用主文档中的某个对象；或
        \item[b)] 在外部化图形中使用 |\label|，而该标签在主文档中被引用。
    \end{enumerate}

    对于情况 a)，外部化图形中的 |\ref| \emph{仅}在你\emph{手动}发出所需的系统调用、或通过 |make| 时才能工作。初始配置 |mode=convert with system call| \emph{不}支持 |\ref|。但你可以复制并粘贴由 |mode=convert with system call| 生成的系统调用，然后手动执行。原因在于 |\ref| 信息存储在主 |.aux| 文件中——但在调用 |mode=convert with system call| 时，该辅助文件还没有完全写完（存在竞争条件）。请注意，|\pageref| 不被支持（抱歉）。因此：如果外部图形中含有 |\ref|，请考虑使用 |mode=list and make|，或复制并粘贴该图的系统调用并手动执行。

    情况 b) 由外部库自动实现。具体来说，外部化图形中的 |\label| 会使外部库为每一张外部图生成单独的辅助文件。这些文件名为 \meta{imagename}|.dpth|。扩展名 |.dpth| 表示该文件还包含图片的深度（即 \tikzname 的 |baseline| 键）。此外，本应写入 |.aux| 文件的任何内容都会被重定向到 |.dpth| 文件——但仅限发生在外部化 |tikzpicture| 环境内部的内容。当主文档加载该图时，会把 |.dpth| 文件的内容拷贝到主 |.aux| 文件中。之后对主文档的后续编译就会包含外部的 |\label| 信息。换句话说，外部图形中的 |\label| 需要遵循下面的工作流程：
    %
    \begin{enumerate}
        \item 外部图形需要连同其 |.dpth| 一起生成（通常由 \tikzname 自动完成）。
        \item 主文档包含该外部图形，并把 |.dpth| 的内容拷贝到自己的主 |.aux| 文件中。
        \item 主文档需要再编译一次，以重新读取其 |.aux| 文件。\footnote{请注意，在 \LaTeX{} 中，无法在 \texttt{\textbackslash begin\{document\}} 之后激活辅助文件的内容。}
    \end{enumerate}
    %
    还有一种特殊情况：当 |\label|/|\ref| 是由 |tikzpicture| 绘制出来的时候。例如图例中的 |\ref| 图片，或 |\pgfplotslegendfromname| 特性。在这种情况下，你需要按照情况 a) 的方式处理，因为 |mode=convert with system call| 本身无法处理这些内容。

    换句话说：外部文档中的 |\label| 可以自动工作，只要把主文档多编译几次即可。|\ref| 则可能需要像上面情况 a) 那样手动调整。


    \paragraph{操作模式}

    \begin{key}{/tikz/external/mode=\mchoice{convert with system call,list and make,$\dotsc$} (initially convert with system call)}
        该键用于修改默认的操作模式。有若干种选项，全部详见
        ~\cite[Section ``Externalization Library'']{tikz}。其中最常用的大概是初始配置 |convert with system call| 以及专门选项 |list and make|。

        选项 |list and make| 配置该库去检查外部图形是否已经存在，并使用已存在的图形。如果图形不存在，该库会\emph{跳过}该图。但它同时还会生成一个 |makefile| 用于生成这些图形，以及一份所有必需图形文件的清单。

        不强制要求使用 |make|：该库只期待你以某种方式生成图片，并不关心「如何生成」。使用
        |make -f |\meta{name-of-tex-file}|.makefile -j 2| 可以并行执行，这确实可能是一个不错的选择。此外，makefile 还支持文件依赖：如果你的某个数据表已更新，外部图形会自动重新生成。\PGFPlots{} 会把任何文件依赖（输入文件和数据表）告知外部库。

        这两种模式具有如下特点：
        %
        \begin{enumerate}
            \item |convert with system call| 是全自动的，会即时完成所有工作。但它\emph{无法}处理外部图片中的 |\ref| 和/或 |\label| 信息。
            \item |list and make| 需要手动（手动发出系统调用）或半自动转换（使用生成的 \meta{main}|.makefile|），并多次运行 |pdflatex|。生成的 Makefile 可以并行处理。此外，|list and make| 对 |\ref| 和 |\label| 提供\emph{完整支持}：在外部化图形中定义的任何 |\label| 仍然可供主文档使用。

                如果你有使用 |legend to name| 的图例，或使用 |\label|/|\ref|，那么你需要先生成定义该 |\label|（或 |legend to name|）的图形，然后在主文档上运行 |pdflatex| 两次。之后就可以把该图例图形外部化了。
        \end{enumerate}
    \end{key}

    完整的参考文档以及其余选项记录在~\cite[``Externalization Library'']{tikz} 中。该参考还包含关于以下内容的信息：
    %
    \begin{itemize}
        \item 如何使用
            |\tikzset{external/|\declareandlabel{force remake}|}| 和
            |\tikzset{external/|\declareandlabel{remake next}|}| 来重新生成选定的图，
        \item 如何用
            |\tikzset{external/|\declareandlabel{export}|=false}| 部分禁用外部化，或用
            |\tikzexternaldisable| 完全禁用，
        \item 如何使用
            |\tikzset{external/optimize command away=\myExpensiveMacro}| 优化转换过程的速度，
        \item 如何用
            |\tikzpicturedependsonfile|\marg{name} 和/或
            |\tikzexternalfiledependsonfile|\marg{external file}\marg{name} 添加更多重新生成的依赖，
        \item 如何启用 |png| 导出的示例，
        \item 如何在未安装 \pgfname{} 的情况下排版此类文档，或
        \item 如何应对 |.pdf| 图片与边界框限制的变通方法。 \index{External Graphics!Bounding Box Issues}
            \index{Bounding Box Control!Image Externalization Problems}
    \end{itemize}


    \paragraph{在未安装 {\normalfont\pgfname} 或 {\normalfont\PGFPlots} 的情况下使用该库}

    有一个小型替换宏包 \declareandlabel{tikzexternal.sty}，可以在每个图都导出之后使用。思路是注释掉 |\usepackage{tikz}| 和 |\usepackage{pgfplots}|，改为写 |\usepackage{tikzexternal}|：
    %
\begin{codeexample}[code only]
% \usepackage{tikz}
% \usepackage{pgfplots}
\usepackage{tikzexternal}
    \tikzexternalize% activate externalization
\begin{document}
    \begin{tikzpicture}
        ...
    \end{tikzpicture}
...
\end{document}
\end{codeexample}
    %
    你不需要安装 \pgfname{}、\tikzname{} 或 \PGFPlots{}。你需要的是 |tikzexternal.sty| 以及所有生成的图（包括图片文件 `|.pdf|'，以及包含 |baseline| 选项信息的 `|.dpth|' 文件）。|tikzexternal.sty| 随 \pgfname{} 一起发布，位于目录
    %
\begin{codeexample}[code only]
latex/pgf/utilities/tikzexternal.sty
\end{codeexample}
    %
    同时 \PGFPlots{} 也附带了一份副本：
    %
\begin{codeexample}[code only]
tex/generic/pgfplots/oldpgfcompatib/pgfplotsoldpgfsupp_tikzexternal.sty
\end{codeexample}
    %
    只需把该文件拷贝到你的目录下，并将其重命名为 |tikzexternal.sty|。


    \paragraph{注意：}

    这个小型替换宏包不支持键值接口。因此，必须用 |\tikzsetexternalprefix| 代替 |prefix| 选项，用 |\tikzsetfigurename| 代替 |figure name| 选项，因为在这种上下文中 |\tikzset| 不可用。另外，如果你之前用过 |\pgfplotsset|，可能还需要定义一个空的宏 |\pgfplotsset|。
\end{pgfplotslibrary}


\subsection[手动使用 PGF 外部化框架]{手动使用 {\normalfont\pgfname} 的外部化框架}

把 \TeX{} 图导出为独立图形文件的另一种方法是使用 \pgfname{} 的外部化框架。它需要更多工作，但比 |external| 库更通用。基本思路是用如下命令把希望导出的部分封装起来：

\declareandlabel{\beginpgfgraphicnamed}\marg{output file name}

\meta{picture contents}

\declareandlabel{\endpgfgraphicnamed}.

\noindent 此外，还需要在导言区用如下命令告知 \pgfname{} 主文档的名称：

\declareandlabel{\pgfrealjobname}\marg{the real job's name}

\noindent 在导言区。这会启用两种不同的模式：
%
\begin{enumerate}
    \item 第一种是正常排版模式。\LaTeX{} 会检查是否存在名为 \meta{output file name}、且带有某种可接受扩展名的文件——如果存在，则用 |\pgfimage| 包含该图形文件，并跳过 \meta{picture contents}。如果该文件不存在，则正常排版 \meta{picture contents}。当 |\jobname| 等于 \meta{the real job's name} 时，使用此模式。
    \item 第二种模式在 |\jobname| 等于 \meta{output file name} 时启用，它进入「转换模式」，用于写出图形文件 \meta{output file name}。此时，\emph{只}把 \meta{picture contents} 写入 |\jobname|，\LaTeX{} 的其余部分照常处理，但会被静默丢弃。

        对于每一个外部化图形文件，都需要手动用
        |pdflatex --jobname |\meta{output file name} 来启动此模式。
\end{enumerate}
%
下面是一个完整的示例。
%
\begin{codeexample}[code only]
\documentclass{article}
\usepackage{pgfplots}
    \pgfrealjobname{test}
\begin{document}
    \begin{figure}
        \beginpgfgraphicnamed{testfigure}
            \begin{tikzpicture}
                \begin{axis}
                    \addplot {x^2};
                \end{axis}
            \end{tikzpicture}
        \endpgfgraphicnamed
        \caption{Our first external graphics example}
    \end{figure}

    \begin{figure}
        \beginpgfgraphicnamed{testfigure2}
            \begin{tikzpicture}
                \begin{axis}
                    \addplot {x^3};
                \end{axis}
            \end{tikzpicture}
        \endpgfgraphicnamed
        \caption{A second graphics}
    \end{figure}
\end{document}
\end{codeexample}
%
\noindent 该文件名为 |test.tex|，我们用如下命令处理它（例如）：
%
\begin{codeexample}[code only]
pdflatex test
\end{codeexample}
%
\noindent 现在我们输入
%
\begin{codeexample}[code only]
pdflatex --jobname testfigure test
pdflatex --jobname testfigure2 test
\end{codeexample}
%
\noindent 进入转换模式。最后这两次调用\emph{只}会把我们命名的图形环境内容写入各自的输出文件 |testfigure.pdf| 和 |testfigure2.pdf|：一个对应 \meta{testfigure}，另一个对应 \meta{testfigure2}。

总之，需要 |\pgfrealjobname|，并对每一个外部化图形环境调用
|pdflatex --jobname |\meta{graphics file}。请注意，必须使用上面给出的语法，\emph{不要}使用 |\begin{pgfgraphicnamed}|。

这些步骤在~\cite{tikz} 的 Section``Externalizing Graphics'' 中有更详细的说明。


\paragraph{注意：}

不要忘记正确的 |\pgfrealjobname| 语句！如果缺失，外部化将根本无法工作。如果写错，每次调用 \LaTeX{} 都会生成空的输出文件。

需要注意的是，这种图片外部化方法并不局限于 \Tikz{} 图环境。事实上，它会把起始与结束语句之间的所有内容都收集到外部文件中。这里隐含的假设是被封装的内容是一个盒子，但你也可以用类似 \LaTeX{} minipage 的方式（或使用 |\vbox|，它功能没那么强大，但对文档其余部分影响也更小）封装完整的段落。

\begin{key}{/pgf/images/aux in dpth=\mchoice{true,false} (initially false)}
    如果该布尔键被设为 |true|，则在外部图片内部生成的任何 |\label| 信息都会存储到前面提到的 |.dpth| 文件中。这样主文档就可以引用文档中被外部化部分的标签信息（尽管你可能需要多次运行 |latex|）。

    对 |\ref| 提供了标签支持，可能也包括 |\cite|。|\pageref| 命令只被部分支持。
\end{key}


\paragraph{在未安装 {\normalfont\pgfname} 的情况下使用该库}

只需把宏包 |\usepackage{tikz}| 和 |\usepackage{pgfplots}| 取消注释，改为使用
%
\begin{codeexample}[code only]
\long\def\beginpgfgraphicnamed#1#2\endpgfgraphicnamed{%
    \begingroup
    \setbox1=\hbox{\includegraphics{#1}}%
    \openin1=#1.dpth
    \ifeof1 \box1
    \else
        \read1 to\pgfincludeexternalgraphicsdp \closein1
        \dimen0=\pgfincludeexternalgraphicsdp\relax
        \hbox{\lower\dimen0 \box1 }%
    \fi
    \endgroup
}
\end{codeexample}
%
代替。这将包含生成的图形文件（并且会遵循 |.dpth| 文件中存储的 |baseline| 信息）。这样一来，你就不需要安装 \pgfname{} 或 \PGFPlots{}。详见~\cite{tikz} 的 Section``Externalizing Graphics''。


\section{从 Matlab 导入}

\subsection{从 Matlab 导入网格数据到 PGFPlots}
\label{sec:pgfplots:import:matlab}

虽然把 Matlab 向量写入文件很容易（使用 |save P.dat data -ASCII|），但导出网格数据就要复杂得多。

主要问题在于如何把网格结构告诉 \PGFPlots{}。

下面是实现这一任务的一个例子：在 Matlab 中，我们有网格数据 |X|、|Y| 和 |Z|，它们是大小相同的矩阵。例如，假设我们有
%
\begin{codeexample}[code only]
[X,Y] = meshgrid( linspace(-1,1,5), linspace(4,5,10) );
Z = X + Y;
surf(X,Y,Z)
\end{codeexample}
%
\noindent 作为数据。然后，我们可以生成一个 $N \times 3$ 的表，按列优先顺序包含所有单独的元素：
%
\begin{codeexample}[code only]
data = [ X(:) Y(:) Z(:) ]
% or -ascii
save P.dat data -ASCII
size(X)
ans =
   10.00    5.00
\end{codeexample}
%
\noindent 其中第二条命令把 $N \times 3$ 的表存入 |P.dat|。最后，我们可以在 \PGFPlots{} 中使用

|\addplot3[surf,mesh/rows=10,mesh/ordering=colwise,shader=interp] file {P.dat};|

来读取这些数据。我们需要提供行数（这里是 $10$）或列数——以及排序方式（对于 Matlab 矩阵是 |colwise|）。

在 \PGFPlots{} 中更快的另一种做法是：在 Matlab 中转置矩阵，并告诉 \PGFPlots{} 它们是按 |rowwise| 顺序排列的。这样，最后一步就变为
%
\begin{codeexample}[code only]
XX=X'; YY=Y'; ZZ=Z';
data = [ XX(:) YY(:) ZZ(:) ]
save P.dat data -ASCII
\end{codeexample}
%
\noindent 对应的 \PGFPlots{} 命令为

|\addplot3[surf,mesh/cols=10,mesh/ordering=rowwise,shader=interp] file {P.dat};|.


\subsection{matlab2pgfplots.m}

这是一个 Matlab$^\text{\textregistered}$ 脚本，尝试把 Matlab 图转换为 \PGFPlots{}。它要求 Matlab 版本 7.4（或更高）。


\paragraph{注意：}

该脚本已经严重过时，只支持 \PGFPlots{} 非常小的一个子集。你可以看看 |matlab2tikz|，这是 Nico Schlömer 编写的转换脚本，地址为

\url{http://www.mathworks.com/matlabcentral/fileexchange/22022-matlab2tikz}

\noindent 它同样使用 \PGFPlots{} 进行 \LaTeX{} 转换。

\medskip
|matlab2pgfplots.m| 的思路是
%
\begin{itemize}
    \item 以完整的 Matlab 图作为输入，
    \item 获取坐标轴标签、坐标轴缩放（对数或普通）以及图例条目，
    \item 获取所有绘图坐标
\end{itemize}
%
并写出一个等价的 \texttt{.pgf} 文件，用 \PGFPlots{} 来排版该图。

其目的\emph{不是}模拟 Matlab。它只是转换工作的第一步。在 Matlab 命令行输入
%
\begin{lstlisting}
> help matlab2pgfplots
\end{lstlisting}
%
在 Matlab 提示符下，可获得关于其特性与局限的更多信息。

该脚本仍处于实验阶段。


\subsection{matlab2pgfplots.sh}

一个 \texttt{bash} 脚本，简单地启动 Matlab 并运行
%
\begin{lstlisting}
    f=hgload( 'somefigure.fig' );
    matlab2pgfplots( 'outputfile.pgf', 'fig', f );
\end{lstlisting}
%
参见上面的 matlab2pgfplots.m。


\subsection{从 Matlab 导入颜色映射表}

有时你可能希望在 \PGFPlots{} 中复用 Matlab 的 |colormap|。下面是一个把 Matlab 颜色映射表转换为 \PGFPlots{} 的小脚本：
%
\begin{codeexample}[code only]
C = colormap;  % gets data of the current colormap.
% C = colormap(jet) % gets data of "jet"
eachnth = 1;
I = 1:eachnth:size(C,1); % this is nonsense for eachnth=1 -- but perhaps you don't want each color.
CC = C(I,:);
TeXstring = [ ...
 sprintf('\\pgfplotsset{\n\tcolormap={matlab}{\n') ...
 sprintf('\t\trgb=(%f,%f,%f)\n',CC') ...
 sprintf('\t}\n}\n') ]
\end{codeexample}


\section{SVG 输出}

可以把每一个 \Tikz{} 图都写入一个可缩放矢量图形（\texttt{.svg}）文件。这与 \PGFPlots{} 无关，它是 \PGF{} 的一个独立驱动。请参见~\cite[Section ``Producing HTML / SVG
Output'']{tikz}。


\section{在 Python 中生成 \PGFPlots{} 图形}

Mario Orne DÍAZ ANADÓN 贡献了一个小型 python 脚本 |pgfplots.py|，它提供了一个简单的接口，可以在 python 中生成 \PGFPlots{} 图。该脚本位于 \PGFPlots{} 安装目录下的
|pgfplots/scripts/pgfplots/pgfplots.py|；其文档可在该文件中找到。

}


\chapter{工具与基础层命令}
\label{cha:pgfplots:lowlevel}

本章记录的命令提供了对 \PGFPlots{} 更基础元素的访问。其中大多数与 \pgfname{} 的基础层密切相关，尤其是各种与坐标轴相关的点命令。还有一些是通用工具，例如循环。

不过，本节中的大多数元素只对高级用户有意义——而且可能只在特殊情况下才用得到。


\section{工具命令}

\begin{command}{\foreach \meta{variables} |in| \meta{list} \marg{commands}}
    \Tikz{} 提供的一个功能强大的循环命令，参见~\cite[Section
    ``Utilities'']{tikz}。
    %
\begin{codeexample}[]
\foreach \x in {1,2,...,4} {Iterating \x. }%
\end{codeexample}

    一个与 \PGFPlots{} 相关的例子可以是
    %
\begin{codeexample}[code only]
\foreach \i in {1,2,...,10} {\addplot table {datafile\i}; }%
\end{codeexample}
\end{command}

\begin{command}{\pgfplotsforeachungrouped \meta{variable} |in| \meta{list} \marg{command}}
    |\foreach| 的一个特殊变体，它有两点不同：在执行 \meta{command} 时不引入额外的分组；并且允许对（简单的！）
    \meta{$x_0$}|,|\meta{$x_1$}|,...,|\meta{$x_n$} 表达式调用数学解析器。
    %
\begin{codeexample}[]
\def\allcollected{}
\pgfplotsforeachungrouped \x in {1,2,...,4} {Iterating \x. \edef\allcollected{\allcollected, \x}}%
All collected = \allcollected.
\end{codeexample}

    一个更有用的例子是与数据表配合使用。下面的例子取自 \PGFPlotstable{}：
    %
\begin{codeexample}[code only]
\pgfplotsforeachungrouped \i in {1,2,...,10} {%
    \pgfplotstablevertcat{\output}{datafile\i} % appends `datafile\i' -> `\output'
}%
% since it was ungrouped, \output is still defined (would not work
% with \foreach)
\end{codeexample}

    \paragraph{备注：}

    特殊语法
    \meta{list}=\meta{$x_0$}|,|\meta{$x_1$}|,...,|\meta{$x_n$}，即以两个元素开头、中间为省略号、末尾为一个元素，会为循环调用数学解析器。因此，当 |/pgf/fpu| 处于活动状态时，它支持的数值范围比其他任何语法都更大。在其他情况下，
    |\pgfplotsforeachungrouped| 会调用 |\foreach|，并在不使用 \TeX{} 分组的情况下提供结果。

    请记住：在坐标轴环境内，所有循环构造（包括自定义循环、|\foreach| 和 |\pgfplotsforeachungrouped|）都需要小心处理：循环变量只能在被立即求值的地方使用；但 \PGFPlots{} 会推迟许多宏的求值。例如，要遍历某些内容并生成形如 |\node at (axis cs:\i,0.5)...| 的坐标轴描述，循环宏 |\i| 会在 |\end{axis}| 处才被求值——但那时循环已经结束，它的值已经丢失。正确的做法是\emph{显式展开}循环变量。例如：
    %
\begin{codeexample}[code only]
\pgfplotsforeachungrouped \i/\j in {
    1 / a,
    2 / b,
    3 / c
}{
    \edef\temp{\noexpand\node at (axis cs: \i,0.5) {\j};}
    % \show\temp % lets TeX show you what \temp contains
    \temp
}
\end{codeexample}
    %
    该示例生成三次循环迭代：|\i=1|、|\j=a|；然后 |\i=2|、|j=b|；然后 |\i=3|、|\j=c|。在循环体内部，它把它们展开，并用「展开定义」|\edef| 把结果赋给一个宏。结果中不再包含 |\i| 或 |\j|（因为它们已经被展开）。然后，调用得到的宏。关于 \TeX{} 命令 |\edef| 以及展开控制的详细信息，请参见随 \PGFPlots{} 一起提供的文档
    \href{file:TeX-programming-notes.pdf}{TeX-programming-notes.pdf}。
\end{command}

\begin{command}{\pgfplotsinvokeforeach\marg{list} \marg{command}}
    |\pgfplotsforeachungrouped| 的一个变体（因而也是 |\foreach| 的变体），它会对 \meta{list} 中的每一个元素，把 \meta{command} 内出现的 |#1| 各替换一次。因此，它实际上假设 \marg{command} 像一个 |\newcommand| 的命令体。

    换句话说，对 \meta{list} 的每一个元素都会调用一次 \meta{command}。\meta{list} 的当前元素以 |#1| 的形式可用。

    与 |\pgfplotsforeachungrouped| 一样，本命令\emph{不}引入额外的作用域（即它也是不分组的）。

    它与 |\foreach \x in |\meta{list}\marg{command} 的区别很细微：|\x| \emph{不会}被展开，而 |#1| 会。
    %
\begin{codeexample}[]
\pgfkeys{
    otherstyle a/.code={[a]},
    otherstyle b/.code={[b]},
    otherstyle c/.code={[c]},
    otherstyle d/.code={[d]}}
\pgfplotsinvokeforeach{a,b,c,d}
    {\pgfkeys{key #1/.style={otherstyle #1}}}
Invoke them:
\pgfkeys{key a} \pgfkeys{key b}
\pgfkeys{key c} \pgfkeys{key d}
\end{codeexample}
相反的例子则使用一个宏（这里是 |\x|）作为循环参数：
\begin{codeexample}[]
\pgfkeys{
    otherstyle a/.code={[a]},
    otherstyle b/.code={[b]},
    otherstyle c/.code={[c]},
    otherstyle d/.code={[d]}}
\pgfplotsforeachungrouped \x in {a,b,c,d}
    {\pgfkeys{key \x/.style={otherstyle \x}}}
Invoke them:
\pgfkeys{key a} \pgfkeys{key b}
\pgfkeys{key c} \pgfkeys{key d}
\end{codeexample}

    \paragraph{限制：}

    目前还不能嵌套使用此命令（因为它不通过作用域提供保护）。
\end{command}

\begin{command}{\pgfmathparse\marg{expression}}
    对 \meta{expression} 调用 \pgfname{} 的数学解析器，并把
    \declareandlabel{\pgfmathresult} 定义为结果。
    %
\begin{codeexample}[]
\pgfmathparse{1+41}

The result is `\pgfmathresult'.
\end{codeexample}
    %
    \noindent \pgfname{} 中的数学引擎通常使用 \TeX{} 内部的算术。也就是说：它适合处理范围在
    $[-16384,16384]$ 内的数字，精度为 $5$ 位。

    对于绘图应用而言，这个数值范围通常太小。\PGFPlots{} 通过
    |\pgfkeys{/pgf/fpu}\pgfmathparse{1+41}| 启用「浮点单元」（fpu），从而扩展数值范围，并让之后所有运算都以浮点数进行。

    在 \PGFPlots{} 中，键 |/pgfplots/use fpu| 默认是打开的，这意味着所有坐标运算都用 |fpu| 进行。然而，所有与 \pgfname{} 相关的绘图操作仍然使用标准数学引擎。

    如果你曾经需要处理扩展精度的数字，可以使用
    %
\begin{codeexample}[]
\pgfkeys{/pgf/fpu}%
\pgfmathparse{1000*1000}

The result is `\pgfmathprintnumber{\pgfmathresult}'.
\end{codeexample}
    %
    请注意，|fpu| 的结果通常不是人类可读格式，因此 |\pgfmathprintnumber| 是排版这类数字的首选方式。

    更多细节请参见 \cite{tikz}。
\end{command}

\begin{pgfplotskey}{use fpu=\mchoice{true,false} (initially true)}
    \PGFPlots{} 提供了多种计算数学表达式的方式，其中 |use fpu| 是最强大的。它的数学运算要么在 |lua backend| 中实现，要么在纯 \TeX{} 实现中完成，数值范围大、精度也足够。

    然而，存储在 |\pgfmathresult| 中的值是晦涩难懂的，需要通过专门的宏来处理。|use fpu| 这一开关只在该数值格式导致问题时才有用——也就是说，它是一个调试开关，在正常操作中不应使用。
\end{pgfplotskey}

\begin{key}{/pgf/declare function=\meta{function definitions}}
    允许定义一个或多个函数。

    参数 \meta {function definitions} 可以包含一个或多个定义，每一个\emph{必须}以分号结尾：
    %
\begin{codeexample}[]
\begin{tikzpicture}
\begin{axis}[
    declare function={
        C=4;
        square(\t)=(\t)^2 + C;
    },
]
    \addplot+ [samples=2] {C*x};

    \addplot {square(x)};
\end{axis}
\end{tikzpicture}
\end{codeexample}
    %
    这些定义本身的形式为 \meta{function}\meta{argument list} =
    \meta{definition}，其中 \meta{argument list} 是用逗号分隔的参数列表，例如 |\t| 或 |\t,\a,\b|。\meta{definition} 是一个数学表达式，它使用这些参数。

    更多细节请参见 \cite{tikz}。
\end{key}

\begin{command}{\pgfplotstableread\marg{file}}
    请参见 \PGFPlotstable{} 的手册 |pgfplotstable.pdf|，它是 \PGFPlots{} 套件的一部分。
\end{command}

\begin{command}{\pgfplotstabletypeset\marg{\textbackslash macro}}
    请参见 \PGFPlotstable{} 的手册 |pgfplotstable.pdf|，它是 \PGFPlots{} 套件的一部分。
\end{command}

\begin{command}{\pgfplotsiffileexists\marg{filename}\marg{true code}\marg{false code}}
    如果 \meta{filename} 存在则调用 \meta{true code}，否则调用 \meta{false code}。例如，可以在循环宏中使用它，依次绘制每一个数据文件，直到没有更多文件为止。
\end{command}

\begin{command}{\pgfplotsutilifstringequal\marg{first}\marg{second}\marg{true code}\marg{false code}}
    一个简单的「字符串比较」工具：当 \meta{first}
    $=$\meta{second} 时调用 \meta{true code}，否则调用 \meta{false code}。它不展开宏。
\end{command}

\begin{commandlist}{\pgfkeys,\pgfeov,\pgfkeysvalueof,\pgfkeysgetvalue}
    这些命令是 \Tikz{} 指定选项方式的一部分，属于其子宏包 |pgfkeys|。命令 |\pgfplotsset| 实际上只是对 |\pgfkeys| 的一层包装。

    关于 |\pgfkeys| 的简短介绍可参见~\cite{keyvalintro}，完整的参考当然是 \Tikz{} 手册~\cite{tikz}。

    键 |\pgfkeysvalueof|\marg{key name} 展开为某个键的值；
    |\pgfkeysgetvalue|\marg{key name}\marg{\textbackslash macro} 把 \meta{key name} 的值存入 \meta{\textbackslash macro}。宏 |\pgfeov| 用于在 |\pgfkeys| 中为代码键分隔参数，请参见上面提到的参考。
\end{commandlist}


\section[PGFPlots 坐标轴内部的命令]
        {{\normalfont\PGFPlots{}} 坐标轴内部的命令}

\begin{command}{\autoplotspeclist}
    此命令不应再使用，但会作为技术实现细节保留下来。请使用 `|cycle list|' 选项，见第~\ref{sec:cycle:list} 节。
\end{command}

\begin{command}{\logten}
    展开为常数 $\log(10)$。对于对数图很有用，因为 $\log(10^i)
    = i\log(10)$。此命令仅在 \Tikz{} 图内部可用。
\end{command}

\begin{command}{\pgfmathprintnumber\marg{number}}
    为 \meta{number} 生成排版精美的输出\footnote{此方法以前叫
    \texttt{\textbackslash prettyprintnumber}。它的功能已并入 \PGF{}，旧命令现已弃用。}。每个刻度标签都使用此方法。

    数字按当前数字打印选项输出；关于不同的数字样式、舍入精度与舍入方法，请参见随本宏包提供的 \PGFPlotstable{} 手册。
\end{command}

\begin{command}{\numplots}
    在任一坐标轴环境、相关样式、选项或命令内部，|\numplots| 展开为图的总数。
\end{command}

\begin{command}{\numplotsofactualtype}
    与 |\numplots| 类似，但此宏返回具有相同绘图处理器的图的总数。因此，如果启用了 |sharp plot|，它返回所有 |sharp plots| 的数量；如果启用了 |ybar|，它返回 |ybar| 图的数量，依此类推。
\end{command}

\begin{command}{\plotnum}
    在 |\addplot| 或任何相关样式、选项或命令内部，|\plotnum| 展开为当前图的编号，从~$0$ 开始。
\end{command}

\begin{command}{\plotnumofactualtype}
    与 |\plotnum| 类似，但它返回在所有同类型图中的编号。所有此类图的总数可用 |\numplotsofactualtype| 获取。
\end{command}

\begin{command}{\coordindex}
    在 |\addplot| 命令内部，此宏展开为当前坐标的编号（从~$0$ 开始）。

    它与 |x filter| 或 |y filter| 配合使用，用于（反）选坐标。
\end{command}


\section{路径操作}

\begin{commandlist}{\path,\draw,\fill,\node,\matrix}
    这些都是 \Tikz{} 绘图命令，全部记录在~\cite{tikz} 中。它们用于绘制或填充路径、生成文本节点或对齐的文本矩阵。它们分别等价于
    \pgfmanualpdflabel{/tikz/draw}{}|\path[draw]|、
    \pgfmanualpdflabel{/tikz/fill}{}|\path[fill]|、
    \pgfmanualpdflabel{/tikz/node}{}|\path[node]|、
    \pgfmanualpdflabel{/tikz/matrix}{}|\path[matrix]|。
\end{commandlist}

\begin{pathoperation}{--}{\meta{coordinate}}
    一个 \Tikz{} 路径操作，用直线连接当前点（即 |--| 之前的最后一个点）与 \meta{coordinate}。
\end{pathoperation}

{\catcode`\|=12
\begin{pathoperation}[noindex]{|-}{\meta{coordinate}}
\pgfmanualpdflabel[\catcode`\|=12 ]{|-}{}%
    一个 \Tikz{} 路径操作，用\emph{两}段直线连接当前点与
    \meta{coordinate}：先竖直，后水平。
\end{pathoperation}

\begin{pathoperation}[noindex]{-|}{\meta{coordinate}}
\pgfmanualpdflabel[\catcode`\|=12 ]{-|}{}%
    一个 \Tikz{} 路径操作，用\emph{两}段直线连接当前点与
    \meta{coordinate}：先水平，后竖直。
\end{pathoperation}
}

\begin{keylist}{/tikz/xshift=\marg{dimension},/tikz/yshift=\marg{dimension}}
    这些 \Tikz{} 键允许把某个对象沿 \meta{dimension} 平移，\meta{dimension} 可以是任何 \TeX{} 长度（或表达式）。
\end{keylist}

\begin{command}{\pgfplotsextra\marg{low-level path commands}}
    在 \PGFPlots{} 坐标轴中执行 \meta{low-level path commands} 的命令。由于坐标轴内部的任何绘图命令都必须推迟到坐标轴完成、缩放初始化之后才能执行，所以不能简单地直接绘制路径。取而代之，必须在坐标轴完成时再绘制它们。对每一个 \Tikz{} 路径这会自动完成——如果你写 |\pgfplotsextra|\marg{commands}，也是手动触发这一过程。
    %
\begin{codeexample}[]
\begin{tikzpicture}
\begin{axis}[xmin=0,xmax=3,ymin=0,ymax=5]
    \pgfplotsextra{
        \pgfpathmoveto{\pgfplotspointaxisxy{1}{2}}
        \pgfpathlineto{\pgfplotspointaxisxy{2}{4}}
        \pgfusepath{stroke}
    }
\end{axis}
\end{tikzpicture}
\end{codeexample}
    %
    上面的例子初始化了一个坐标轴，并在坐标轴就绪后执行基础层路径命令。多个 |\path|、|\addplot| 和 |\pgfplotsextra| 命令的执行顺序与它们在环境中出现的顺序相同。\footnote{堆叠图例外，其顺序可能相反，参见键 \texttt{reverse stack plots}。}
\end{command}

\begin{command}{\pgfplotspathaxisoutline}
    生成一条与当前坐标轴外轮廓相似的路径。该路径用于裁剪路径和背景路径（如果存在）。
\end{command}


\section{指定基础坐标}
\label{sec:basic:coordinates}

\begin{commandlist}{%
    \pgfplotspointaxisxy\marg{x coordinate}\marg{y coordinate},
    \pgfplotspointaxisxyz\marg{x coordinate}\marg{y coordinate}\marg{z coordinate}%
}
    类似 |\pgfpointxy| 的点命令，接受逻辑的绝对坐标，返回一个底层点。所有从用户变换到对数的变换都会被应用。

    由于变换是在坐标轴完成后才初始化的，因此此命令需要被推迟（参见 |\pgfplotsextra|）。

    此命令是基础层版本的 |axis cs:|\meta{x
    coordinate}|,|\meta{y coordinate}|,|\meta{z coordinate}。

    请注意，这也是可视化阶段的默认坐标系；换句话说：如果你写 |\draw (1,2) -- (1,4)|，\PGFPlots{} 会自动使用 |(axis cs:1,2) -- (axis cs:1,4)|。
\end{commandlist}

\begin{commandlist}{%
    \pgfplotspointaxisdirectionxy\marg{x coordinate}\marg{y coordinate},
    \pgfplotspointaxisdirectionxyz\marg{x coordinate}\marg{y coordinate}\marg{z coordinate}%
}
    类似 |\pgfpointxy| 的点命令，接受逻辑\emph{相对}坐标，返回一个底层点。所有从用户变换到对数的变换都会被应用。与
    |\pgfplotspointaxisxy| 的区别在于：这里会跳过线性变换的平移部分（对比 |disabledatascaling|）。

    此命令是基础层版本的 |axis direction cs:|\meta{x
    coordinate}|,|\meta{y coordinate}|,|\meta{z coordinate}。更多细节请参见 |axis direction cs| 的文档。

    当需要提供具有\emph{相对}性质的量（如方向或长度）时，使用此命令。一个用例是绘制椭圆：
    %
\begin{codeexample}[]
\begin{tikzpicture}
\begin{axis}[
    xmin=-3,   xmax=3,
    ymin=-3,   ymax=3,
    extra x ticks={-1,1},
    extra y ticks={-2,2},
    extra tick style={grid=major},
]
    \draw [red] \pgfextra{
        \pgfpathellipse{\pgfplotspointaxisxy{0}{0}}
            {\pgfplotspointaxisdirectionxy{1}{0}}
            {\pgfplotspointaxisdirectionxy{0}{2}}
    % see also the documentation of
    % 'axis direction cs' which
    % allows a simpler way to draw this ellipse
    };
    \draw [blue] \pgfextra{
        \pgfpathellipse{\pgfplotspointaxisxy{0}{0}}
            {\pgfplotspointaxisdirectionxy{1}{1}}
            {\pgfplotspointaxisdirectionxy{0}{2}}
    };
    \addplot [only marks,mark=*] coordinates
        { (0,0) };
\end{axis}
\end{tikzpicture}
\end{codeexample}

    由于变换是在坐标轴完成后才初始化的，因此此命令需要放在 \tikzname{} |\path| 命令内部（如上面例子中的 |\draw|），或者放在 |\pgfplotsextra| 内部。
\end{commandlist}

\begin{commandlist}{%
    \pgfplotspointrelaxisxy\marg{rel x coordinate}\marg{rel y coordinate},
    \pgfplotspointrelaxisxyz\marg{rel x coordinate}\marg{rel y coordinate}\marg{rel z coordinate}%
}
    接受\emph{相对}坐标的点命令：$x=0$ 对应坐标轴的\emph{下}限，$x=1$ 对应\emph{上}限。

    这些命令用于 |rel axis cs|。

    请注意，变换只有在坐标轴完成后才会初始化！这意味着你需要使用 |\pgfplotsextra|。
\end{commandlist}

\begin{commandlist}{%
    \pgfplotspointdescriptionxy\marg{$x$ fraction}\marg{$y$ fraction},
    \pgfplotsqpointdescriptionxy\marg{$x$ fraction}\marg{$y$ fraction}%
}
    点命令：|{0}{0}| 是坐标轴包围盒的左下角，|{1}{1}| 是右上角；其余值位于二者之间。`|q|' 版本更快，因为它不对参数调用数学解析器。

    它们用于 |axis description cs|，见
    Section~\ref{pgfplots:sec:axis:description:cs}。
\end{commandlist}

\begin{commandlist}{\pgfplotspointaxisorigin}
    原点处的点坐标，即 $(0,0,0)$。如果原点不在坐标轴范围内，则返回边界上最近的点。

    这与 |origin| 锚点返回的坐标相同。
\end{commandlist}

\begin{commandlist}{%
    \pgfplotstransformcoordinatex\marg{x coordinate of an axis},
    \pgfplotstransformcoordinatey\marg{y coordinate of an axis},
    \pgfplotstransformcoordinatez\marg{z coordinate of an axis}%
}
    把 |\pgfmathresult| 定义为与输入参数对应的 \PGF{} 底层坐标。

    该命令会应用任何 |[xyz] coord trafo| 键、数据缩放和/或对数，以及 \PGFPlots{} 用于把输入坐标映射为内部坐标的任何其他处理。

    结果可以在 |\pgfpointxy| 语句内部使用（也就是说，它仍然需要与相应的 \PGF{} 单位向量相乘缩放）。
    %
\begin{codeexample}[]
\begin{tikzpicture}
\pgfplotsset{compat/pgfpoint substitution=1.3}
\begin{axis}[xmin=0,xmax=2,ymin=0,ymax=5]
    \pgfplotsextra{
        \pgfplotstransformcoordinatex{1}
        \let\xcoord=\pgfmathresult
        \pgfplotstransformcoordinatey{1}
        \let\ycoord=\pgfmathresult
        \pgfpathcircle
            {\pgfqpointxy{\xcoord}{\ycoord}}
            {5pt}
        \pgfusepath{fill}
    }
\end{axis}
\end{tikzpicture}
\end{codeexample}
    %
    请注意，\PGFPlots{} 默认会把 |\pgfqpointxy| 替换为 |\pgfplotspointaxisxyz|——而此命令本身已经隐式完成了坐标变换。为了看出区别，前面的例子先通过
    |compat/pgfpoint substitution=1.3| 关闭了这种坐标系的自动替换。
    %
    此命令的结果也可以作为数学方法
    |transformcoordinatex| 使用（参见 |axis cs| 的文档）。

    请注意，变换只有在坐标轴完成后才会初始化。这意味着你需要像上面例子那样使用 |\pgfplotsextra|。
\end{commandlist}

\begin{commandlist}{%
    \pgfplotstransformdirectionx\marg{x direction of an axis},
    \pgfplotstransformdirectiony\marg{y direction of an axis},
    \pgfplotstransformdirectionz\marg{z direction of an axis}%
}
    把 |\pgfmathresult| 定义为 \PGF{} 底层\emph{方向向量分量}。

    方向向量需要与某个坐标\emph{相加}才能得到一个坐标，参见
    |\pgfplotspointaxisdirectionxy| 和 |axis direction cs| 的文档。

    参数 \meta{x direction of an axis} 的处理方式（几乎）与处理绝对位置的宏
    |\pgfplotstransformcoordinatex| 相同。唯一区别是：\emph{方向}不需要平移变换。

    此命令的结果也可以作为数学方法
    |transformdirectionx| 使用（参见 |axis direction cs| 的文档）。

    关于此命令的细节与示例，请参见 |axis direction cs|。
\end{commandlist}

% this command is for internal use only:
%--------------------------------------------------
% \begin{command}{\pgfplotsconvertunittocoordinate\marg{x, y or z}\marg{dimension}}
%     Converts a dimension (with unit!) to a corresponding $x$-, $y$- or $z$-coordinate. The result will be written to |\pgfmathresult| (without units).
%
%     It is possible to use the result as arguments for the |\pgfpointxyz| commands.
%
%     The effect is to multiply \meta{dimension} with the inverse length of the unit vector for the specified axis. These lengths are precomputed in \PGFPlots{} so the operation is fast.
% \begin{codeexample}[code only]
% \pgfplotsconvertunittocoordinate{x}{5pt}
% % now, the command uses exactly 5pt in x direction:
% \pgfqpointxyz{\pgfmathresult}{4}{3}
% \end{codeexample}
% \end{command}
%--------------------------------------------------

\begin{commandlist}{%
    \pgfplotspointunitx,
    \pgfplotspointunity,
    \pgfplotspointunitz%
}
    返回画布 $x$、$y$ 或 $z$ 单位向量的底层点命令。

    |\pgfplotspointunitx| 是 $x$ 方向上的 \pgfname{} 单位向量。

    这些向量本质上分别等同于\linebreak
    |\pgfqpointxyz{1}{0}{0}|、|\pgfqpointxyz{0}{1}{0}|\linebreak
    和 |\pgfqpointxyz{0}{0}{1}|。

    单位 $z$ 向量只在三维坐标轴中才有定义。
\end{commandlist}

\begin{commandlist}{%
    \pgfplotsunitxlength,
    \pgfplotsunitylength,
    \pgfplotsunitzlength,
    \pgfplotsunitxinvlength,
    \pgfplotsunityinvlength,
    \pgfplotsunitzinvlength%
}
    展开为相应单位向量 $x_i$ 的长度 $\lVert x_i \rVert$，或其倒数 $1/\lVert x_i \rVert$ 的宏。这些宏可以在 |\pgfmathparse| 内部使用。

    $x_i$ 即 |\pgfplotspointunitx| 系列。
\end{commandlist}

\begin{command}{\pgfplotsqpointoutsideofaxis%
        \marg{three-char-string}\marg{coordinate}\marg{normal distance}%
}
    对于二维图，返回四条可用坐标轴之一上的一个点坐标；对于三维图，则返回十二条可用坐标轴之一上的点坐标。

    所指定的坐标轴由一个三字符字符串唯一标识，作为命令的第一个参数传入。三个字符中的第一个为 `|0|'，表示该坐标轴的 $x$ 坐标经过坐标轴的下限；为 `|1|'，表示经过上限；为 `|2|'，表示经过原点。第二个字符同样是 |0|、|1| 或 |2|，表示 $y$ 轴上的位置。第三个字符对应第三维，即 $z$ 轴；二维图应保留为 `|0|'。但是，三个字符中\emph{必须}有一个为 `|v|'，表示该坐标轴\underline{是变化的}。例如，|v01| 表示 $\{ (x,y_{\min},z_{\max}) \vert x \in \R \}$。

    第二个参数 \meta{coordinate} 是该坐标轴上的逻辑坐标。由于该坐标轴的两个坐标是固定的，\meta{coordinate} 指的是坐标轴中\underline{变化}的那一分量。它必须是一个不带单位的数字；此处不支持数学表达式。

    第三个参数 \meta{normal distance} 是一个长度，例如 |10pt|。它把坐标沿外法向量方向从指定坐标轴向外平移。外法向量总是指向远离坐标轴的方向。它由 |\pgfplotspointouternormalvectorofaxis| 计算。

    此命令有若干变体，都记录在源代码中。其中一个特别有用：
\end{command}

\begin{command}{\pgfplotsqpointoutsideofaxisrel%
        \allowbreak\marg{three-char-string}\allowbreak\marg{axis fraction}\allowbreak\marg{normal distance}%
}
    此点坐标是 |\pgfplotsqpointoutsideofaxis| 的一个变体，它允许传入 \meta{axis fraction} 而不是绝对坐标。该分数是 $0$（坐标轴下限）到 $1$（坐标轴上限）之间的一个数，即以坐标轴总长的百分比表示。允许使用负值或大于 $1$ 的值。

    |\pgfplotsqpointoutsideofaxisrel| 命令在思路上类似于
    |rel axis cs|。

    与反向坐标轴配合使用时有一点特殊：如果坐标轴已通过 |x dir=reverse| 反向，且 |allow reversal of rel axis cs| 为 true，那么 $0$ 表示\emph{上}限，$1$ 表示\emph{下}限。这样，仅因坐标轴反向，坐标值不会发生改变。
        \index{allow reversal of rel axis cs}%
\end{command}

\begin{command}{\pgfplotspointouternormalvectorofaxis\marg{three-char-string}}
    一个点命令，返回相应坐标轴的外法向量。该法向量长度为 $1$（用
    |\pgfpointnormalised| 计算）。它与 |\pgfplotsqpointoutsideofaxis| 及其变体中使用的是同一个法向量。

    此命令的输出会被缓存，并在坐标轴的整个生命周期内复用。
\end{command}

\begin{command}{\pgfplotsticklabelaxisspec\marg{x, y or z}}
    展开为指定坐标轴（\meta{x}、\meta{y} 或 \meta{z}）所在的、用于承载刻度标签的坐标轴的三字符标识。
\end{command}

\begin{command}{\pgfplotsvalueoflargesttickdimen\marg{x, y or z}}
    展开为刻度位置到其刻度标签包围盒沿外单位法向量方向的最大距离。它还包括 |ticklabel shift| 键的值。

    此值用于 |ticklabel cs|。
\end{command}

\begin{commandlist}{
    \pgfplotsmathfloatviewdepthxyz\marg{x}\marg{y}\marg{z},
    \pgfplotsmathviewdepthxyz\marg{x}\marg{y}\marg{z}%
}
    这两个宏都把 |\pgfmathresult| 定义为三维点 $\bar x = (x,y,z)$ 的「深度」。深度定义为 $\bar x$ 与当前坐标轴视图方向 $\vec d$ 的内积。

    对于 |\pgfplotsmathfloatviewdepthxyz|，参数按浮点数解析，结果也以浮点数编码。可以用
    |\pgfmathfloattofixed{\pgfmathresult}| 生成定点表示。

    对于 |\pgfplotsmathviewdepthxyz|，内积使用 \TeX{} 算术完成，结果以定点数赋值。它略快，但数值范围小得多。

    这两个命令都只能在三维 \PGFPlots{} 坐标轴\emph{内部}使用（坐标轴初始化之后，参见
    |\pgfplotsextra|）。
\end{commandlist}

\begin{texif}{pgfplotsthreedim}
    一个 \TeX{} |\if|：如果坐标轴是三维的，求 \meta{true code}；否则求 \meta{else code}。
\end{texif}


\section{访问坐标轴范围}

也可以在可视化阶段，即在 |\end{axis}| 期间，访问坐标轴范围。请参见
第~\pageref{page:access:limits} 页上 |xmin| 的参考文档。


\section{访问点坐标的值}

在可视化阶段，\PGFPlots{} 提供了对当前正在处理的坐标及其值的访问。

这种访问需要调用特定的宏。这些宏把坐标值写入一些公开的键值对中。然后，就可以访问当前点的 $x$、$y$、$z$ 和颜色数据。

\begin{commandlist}{\pgfplotspointgetcoordinates,\pgfplotspointgetcoordinates\marg{point}}

    调用此宏之后，会设置以下键：

    \declaretext{/data point/x} 将包含当前点的 $x$ 坐标。

    \declaretext{/data point/y} 将包含当前点的 $y$ 坐标。

    \declaretext{/data point/z} 将包含当前点的 $z$ 坐标（如果适用）。

    \declaretext{/data point/meta} 将包含当前点的
    |point meta| 值（如果适用）。

    \declaretext{/data point/index} 将包含当前点在坐标流中的索引。它实际上与 |\coordindex| 相同。

    此命令实际上支持两种工作模式：
    %
    \begin{enumerate}
        \item 不带参数。此时，它返回当前绘图处理器即将处理的点的值。
        \item 带花括号参数。此时，它接受一个坐标并相应地为各键赋值。请注意，此命令仅支持二维坐标轴，并且只赋值
            |/data point/x| 和 |/data point/y|。
    \end{enumerate}

    返回的值与你在坐标轴上读到的相同，也与你写进 |axis cs| 的相同。

    这意味着任何 |x coord inv trafo| 都已经作用在该值上。也就是说，即使内部坐标是以对数形式存储的，这里也已经调用了指数函数。

    对于任何 |scatter| 图（包括 |nodes near coords|），此函数都会被隐式调用。这样就可以一次性访问\emph{所有}坐标值：
    %
\begin{codeexample}[]
\begin{tikzpicture}
\begin{axis}
    \addplot+ [
        domain=0:360,
        samples=9,
        only marks,
        nodes near coords={%
            \footnotesize
            $(\pgfmathprintnumber
                {\pgfkeysvalueof{/data point/x}},
               \pgfmathprintnumber
                {\pgfkeysvalueof{/data point/y}})$%
        },
    ] {sin(x)};
\end{axis}
\end{tikzpicture}
\end{codeexample}
    %
    这个例子之所以能工作，是因为 |\pgfplotspointgetcoordinates| 是 |nodes near coords| 标准实现的一部分；得到的值可以直接使用。请注意，如果只打印一个值，前面的例子还可以更简单：|nodes near coords| 会退回到使用 |point meta|。而 |point meta| 默认情况下本来就包含 $y$ 坐标。

    一个更高级的例子是：在 |ybar| 图中，希望把节点放在坐标轴下端，并从节点到相应柱条画一些虚线：
    %
\begin{codeexample}[]
\begin{tikzpicture}
\begin{axis}[
    ybar,
    nodes near coords,
    %
    % we want to provide absolute 'at' values
    % for the nodes:
    scatter/position=absolute,
    every node near coord/.style={
        at={(\pgfkeysvalueof{/data point/x},-1)},
        % pretty printing:
        anchor=north,
        /pgf/number format/fixed,
        /pgf/number format/precision=1,
        % assign a name which can be referenced below:
        name=NNC\pgfkeysvalueof{/data point/index},
    },
    % ... draw a dotted line between
    % the marker and the bar:
    /pgfplots/scatter/@post marker code/.add code={}{
        \draw [dotted,help lines]
            (NNC\pgfkeysvalueof{/data point/index})
            -- (\pgfkeysvalueof{/data point/x},
            {min(0,\pgfkeysvalueof{/data point/y})});
    },
    % assign suitable tick labels:
    xtick=data,
]
    % some dummy data:
    \addplot+ [
        domain=0:360,
        bar width=360/9,
        samples=9,
    ] {sin(x)};
\end{axis}
\end{tikzpicture}
\end{codeexample}
    %
    同样，该命令是作为 |nodes near coords| 的一部分隐式使用的，因此在例子本身中并未出现。

    \paragraph{另请参见} 在线相关示例：
    \url{https://tex.stackexchange.com/a/141006}。它演示了如何根据数值（放在柱条内部或柱条上方）放置由 |nodes near coords| 生成的节点。

    下面的例子使用一个花括号参数，我们希望获取该参数的坐标值：
    %
\begin{codeexample}[]
% requires \usetikzlibrary{intersections}
\begin{tikzpicture}
\begin{axis}
    \addplot [name path=A] {x^3};

    \draw [red,name path=HelperLine]
        (axis cs:-6,-100) -- (axis cs:6,130);

    \draw [
        font=\footnotesize,
        name intersections={of=A and HelperLine},
    ]
    node [pin={right:
      \pgfplotspointgetcoordinates{(intersection-1)}
        $(\pgfmathprintnumber[fixed]{
            \pgfkeysvalueof{/data point/x}},
          \pgfmathprintnumber[fixed]{
            \pgfkeysvalueof{/data point/y}})$
      }
    ] at (intersection-1) {}
    node [pin={
      \pgfplotspointgetcoordinates{(intersection-2)}
        $(\pgfmathprintnumber[fixed]{
            \pgfkeysvalueof{/data point/x}},
          \pgfmathprintnumber[fixed]{
            \pgfkeysvalueof{/data point/y}})$
      }
    ] at (intersection-2) {};
\end{axis}
\end{tikzpicture}
\end{codeexample}
    %
    该例子计算并命名了 |A| 与 |HelperLine| 的交点。下面的代码在交点处生成 pin。为此，它使用
    |\pgfplotspointgetcoordinates|\marg{point}，该宏会定义 |/data point/x|
    和 |/data point/y|。然后，用 |\pgfmathprintnumber| 把这些值格式化输出。

    在第二种模式下，|\pgfplotspointgetcoordinates|\marg{point} 期望 \meta{point} 为以下二者之一：
    %
    \begin{enumerate}
        \item 一个基础层 \PGF{} 点，例如 |\pgfpointanchor{A}{center}|；或
        \item 一个 \Tikz{} 点，例如 |(A.base)| 或 |(3,5)|。
    \end{enumerate}
\end{commandlist}

\begin{command}{\pgfplotspointgetnormalizedcoordinates}
    与 |\pgfplotspointgetcoordinates| 非常相似的一个宏。同样，它也应在可视化阶段调用。

    它赋值的输出宏完全相同，但值不同。更准确地说，它定义的宏为：

    \declaretext{/data point/x} 将包含当前点的\emph{归一化} $x$ 坐标。

    \declaretext{/data point/y} 将包含当前点的\emph{归一化} $y$ 坐标。

    \declaretext{/data point/z} 将包含当前点的\emph{归一化} $z$ 坐标（如果适用）。

    \declaretext{/data point/meta} 将包含当前点的
    |point meta| 值（如果适用）。

    \declaretext{/data point/index} 将包含当前点在坐标流中的索引。它实际上与 |\coordindex| 相同。

    关键字 \declaretext{normalized} 表示这些值采用适合坐标轴使用的数值形式。更具体地说：任何用户 |x coord inv trafo| 都会被\emph{忽略}。一个重要的例子是 |symbolic x coords|：归一化坐标将是某些关联数字，而不是符号。而
    |\pgfplotspointgetcoordinates| 返回的结果则是符号。对于对数坐标轴，归一化值是对数。

    通常，当你希望执行某种数学运算（如求平均或做差）时，归一化值会更有用。

    此函数需要显式调用。目前 |ybar stacked| 用它来对齐 |nodes near coords|。
\end{command}


\section{图层访问}

\begin{command}{\pgfplotsonlayer\marg{layer name}}
    一个底层命令，它会检查当前坐标轴是否启用了图层支持；如果启用，则调用 |\pgfonlayer|\marg{layer name}。

    必须有一个 |\endpgfplotsonlayer| 来结束该环境。
\end{command}

\begin{command}{\endpgfplotsonlayer}
    |\pgfplotsonlayer| 的结束。
\end{command}

\begin{command}{\pgfonlayer\marg{layer name}}
    \PGF{} 的一个底层命令，它会把直到匹配的 |\endpgfonlayer| 之前的所有内容收集到图层 \meta{layer name} 中。

    \meta{layer name} 必须是活动的，即它必须是 |set layers| 图层名集合的一部分。

    唯一的特殊情况是：如果你在某处调用 |\pgfdeclarelayer{discard}|，这个特殊图层具有一个「神奇的名字」，当用 |\pgfonlayer{discard}| 启用时，它充当 |/dev/null|：它不必是活动的，分配给该图层的所有内容都会被丢弃——只要它不在图层名配置之中。

    必须有一个 |\endpgfonlayer| 来结束该环境。
\end{command}

\begin{command}{\endpgfonlayer}
    |\pgfonlayer| 的结束。
\end{command}


\begin{command}{\pgfsetlayers\marg{layer list}}
    这是 \PGF{} 的一个底层命令。在撰写本文时，它是告诉 \PGF{} 当前/下一张图应使用哪些图层的唯一方式。它由 |set layers| 隐式使用。
\end{command}


\printindex

\bibliographystyle{abbrv} %gerapali} %gerabbrv} %gerunsrt.bst} %gerabbrv}% gerplain}
\nocite{pgfplotstable}
\nocite{programmingnotes}
\bibliography{pgfplots}
\end{document}

% -----------------------------------------------------------------------------
% For Stefan Pinnow as reminder on what to look for when editing the manual
% -----------------------------------------------------------------------------
% There should be no line breaks in the following environments
% - |...|
% - \declareandlabel{...}
% - \verbpdfref{...}
% If "MakeTikzPictures" isn't running through check one of the externalized
% LOG files what is the cause of that.
% -----------------------------------------------------------------------------
'
    \expandafter\ifx\csname pgfplotsclickabledisabled\endcsname\relax
        \usepgfplotslibrary{clickable}
    \fi
\fi

%\usepackage{fp}
% ATTENTION:
% this requires pgf version NEWER than 2.00:
%\usetikzlibrary{fixedpointarithmetic}

\usepgfplotslibrary{dateplot,units,groupplots,statistics}

\pgfplotsset{
    table/search path={.,plotdata}
}

\usepackage[a4paper,left=2.25cm,right=2.25cm,top=2.5cm,bottom=2.5cm]{geometry}
\usepackage{amsmath,amssymb}
\usepackage{xxcolor}

\usepackage{ifluatex}

% tlmgr install zapfding symbol
\usepackage{pifont}

\def\lua{\textsc{lua}}

\ifluatex
    %\usepackage[utf8]{luainputenc}
\else
    \expandafter\ifx\csname XeTeXrevision\endcsname\relax
        \usepackage[utf8]{inputenc}
    \fi
\fi
\usepackage{amsmath}

% tlmgr eurosym
\usepackage{eurosym}

% tlmgr install units
\usepackage{nicefrac}

\def\eps{\textsc{eps}}

% Copyright 2006 by Till Tantau
%
% This file may be distributed and/or modified
%
% 1. under the LaTeX Project Public License and/or
% 2. under the GNU Free Documentation License.
%
% See the file doc/generic/pgf/licenses/LICENSE for more details.

% $Header: /cvsroot/pgfplots/pgfplots/doc/latex/pgfplots/pgfmanual-en-macros.tex,v 1.4 2009/03/05 22:23:22 ludewich Exp $


% This file is part of the pgf manual; I reuse it here with
% permission.
%
% I have changed very few things (marked with 'CF')

\providecommand\href[2]{\texttt{#1}}


\colorlet{examplefill}{yellow!80!black}
\definecolor{graphicbackground}{rgb}{0.96,0.96,0.8}
\definecolor{codebackground}{rgb}{0.8,0.8,1}

\newenvironment{pgfmanualentry}{%
  \pgfmanualresetlabel
  \list{}{\leftmargin=2em\itemindent-\leftmargin\def\makelabel##1{\hss##1}}%
}{\endlist}
\newcounter{pgfmanualentry}
\newcommand\pgfmanualentryheadline[1]{%
  \itemsep=0pt\parskip=0pt\relax
  {\raggedright\item\refstepcounter{pgfmanualentry}\pgfmanualstorelabel
   \strut{#1}\par}%
  \topsep=0pt}
\newcommand\pgfmanualbody{\parskip3pt}

\makeatletter
\newif\if@pgfmanualfirstlabel

\newcommand\pgfmanualresetlabel{%
  \@pgfmanualfirstlabeltrue}

% Store info of first label in \pgfmanual@labelinfo for latter use.
\newcommand\pgfmanualstorelabel{%
  \if@pgfmanualfirstlabel
    \begingroup
    \label@hook
    \xdef\pgfmanual@labelinfo
      {{\@currentlabel}{\thepage}{\@currentlabelname}{\@currentHref}{}}%
    \endgroup
    \global\@pgfmanualfirstlabelfalse
  \fi
}

% Used like a prifix of \label, \labelprefix makes the following 
% \label use label info stored in \pgfmanual@labelinfo. Based on the 
% 
% Usage:
%     \labelprefix\label{<label name>}
% Possible extension:
%     \labelprefix[<offset>]\label{<label name>}
%
\newcommand\labelprefix[2]{%
  \def\@pgfmanualtemp{#1}%
  \ifx\@pgfmanualtemp\pgfmanual@labelcmd
    \@bsphack
    \begingroup
      \def\label@name{#2}%
      % \label@hook
      \protected@write\@auxout{}{%
        \string\newlabel{#2}{\pgfmanual@labelinfo}}%
    \endgroup
    \@esphack
  \else
    \PackageError{pgfmanual-macros}{\string\labelprefix should use as prefix of \string\label.}{}%
  \fi
}
\def\pgfmanual@labelcmd{\label}
\makeatother

\newenvironment{pgflayout}[1]{
  \begin{pgfmanualentry}
    \pgfmanualentryheadline{\texttt{\string\pgfpagesuselayout\char`\{\declare{#1}\char`\}}\oarg{options}}
    \index{#1@\protect\texttt{#1} layout}%
    \index{Page layouts!#1@\protect\texttt{#1}}%
    \pgfmanualbody
}
{
  \end{pgfmanualentry}
}


\newenvironment{command}[1]{
  \begin{pgfmanualentry}
    \extractcommand#1\@@
    \pgfmanualbody
}
{
  \end{pgfmanualentry}
}

%% MW: START MATH MACROS
\def\mvar#1{{\rmfamily\textit{#1}}}

\makeatletter

\def\extractmathfunctionname#1{\extractmathfunctionname@#1(,)\tmpa\tmpb}
\def\extractmathfunctionname@#1(#2)#3\tmpb{\def\mathname{#1}}

\def\extractmathoperatorname{\begingroup\def\mvar##1{}\def\ {}\extractmathoperatorname@}
\def\extractmathoperatorname@#1{\xdef\mathname{#1}\endgroup}

\makeatother

\newenvironment{math-function}[1]{
    \begin{pgfmanualentry}
        \extractmathfunctionname{#1}
        \pgfmanualentryheadline{\texttt{#1}}%
        \index{\mathname @\protect\texttt{\mathname} math function}%
        \index{Math functions!\mathname @\protect\texttt{\mathname}}
        \pgfmanualbody
}
{
    \end{pgfmanualentry}
}

\newenvironment{math-operator}[1]{
    \begin{pgfmanualentry}
        \extractmathoperatorname{#1}
        \pgfmanualentryheadline{\texttt{#1}}%
        \index{\mathname @\protect\texttt{\mathname} math operator}%
        \index{Math operators!\mathname @\protect\texttt{\mathname}}
        \pgfmanualbody
}
{%
    \end{pgfmanualentry}
}

\newenvironment{math-constant}[1]{
    \begin{pgfmanualentry}
        \pgfmanualentryheadline{\texttt{#1}}%
        \index{#1@\protect\texttt{#1} math constant}%
        \index{Math constants!#1@\protect\texttt{#1}}
        \pgfmanualbody
}
{
    \end{pgfmanualentry}
}
\def\calcname{\textsc{calc}}
%% MW: END MATH MACROS


\def\extractcommand#1#2\@@{%
  \pgfmanualentryheadline{\declare{\texttt{\string#1}}#2}%
  \removeats{#1}%
  \index{\strippedat @\protect\myprintocmmand{\strippedat}}}


% \begin{environment}{{name}\marg{arguments}}
\renewenvironment{environment}[1]{
  \begin{pgfmanualentry}
    \extractenvironement#1\@@
    \pgfmanualbody
}
{
  \end{pgfmanualentry}
}

\def\extractenvironement#1#2\@@{%
  \pgfmanualentryheadline{{\ttfamily\char`\\begin\char`\{\declare{#1}\char`\}}#2}%
  \pgfmanualentryheadline{{\ttfamily\ \ }\meta{environment contents}}%
  \pgfmanualentryheadline{{\ttfamily\char`\\end\char`\{\declare{#1}\char`\}}}%
  \index{#1@\protect\texttt{#1} environment}%
  \index{Environments!#1@\protect\texttt{#1}}}


\newenvironment{plainenvironment}[1]{
  \begin{pgfmanualentry}
    \extractplainenvironement#1\@@
    \pgfmanualbody
}
{
  \end{pgfmanualentry}
}

\def\extractplainenvironement#1#2\@@{%
  \pgfmanualentryheadline{{\ttfamily\declare{\char`\\#1}}#2}%
  \pgfmanualentryheadline{{\ttfamily\ \ }\meta{environment contents}}%
  \pgfmanualentryheadline{{\ttfamily\declare{\char`\\end#1}}}%
  \index{#1@\protect\texttt{#1} environment}%
  \index{Environments!#1@\protect\texttt{#1}}}


\newenvironment{contextenvironment}[1]{
  \begin{pgfmanualentry}
    \extractcontextenvironement#1\@@
    \pgfmanualbody
}
{
  \end{pgfmanualentry}
}

\def\extractcontextenvironement#1#2\@@{%
  \pgfmanualentryheadline{{\ttfamily\declare{\char`\\start#1}}#2}%
  \pgfmanualentryheadline{{\ttfamily\ \ }\meta{environment contents}}%
  \pgfmanualentryheadline{{\ttfamily\declare{\char`\\stop#1}}}%
  \index{#1@\protect\texttt{#1} environment}%
  \index{Environments!#1@\protect\texttt{#1}}}


\newenvironment{shape}[1]{
  \begin{pgfmanualentry}
      \pgfmanualentryheadline{形状 {\ttfamily\declare{#1}}}%
    \index{#1@\protect\texttt{#1} shape}%
    \index{Shapes!#1@\protect\texttt{#1}}
    \pgfmanualbody
}
{
  \end{pgfmanualentry}
}


\newenvironment{handler}[1]{
  \begin{pgfmanualentry}
    \extracthandler#1\@nil%
    \pgfmanualbody
}
{
  \end{pgfmanualentry}
}

\def\gobble#1{}
\def\extracthandler#1#2\@nil{%
  \pgfmanualentryheadline{键处理器 \meta{key}{\ttfamily/\declare{#1}}#2}%
  \index{\gobble#1@\protect\texttt{#1} handler}%
  \index{Key handlers!#1@\protect\texttt{#1}}
}


\makeatletter


\newenvironment{stylekey}[1]{
  \begin{pgfmanualentry}
    \def\extrakeytext{样式， }
    \extractkey#1\@nil%
    \pgfmanualbody
}
{
  \end{pgfmanualentry}
}


% \begin{key}{/path/x=value}
% \begin{key}{/path/x=value (initially XXX)}
% \begin{key}{/path/x=value (default XXX)}
\newenvironment{key}[1]{
  \begin{pgfmanualentry}
    \def\extrakeytext{}
    %\def\altpath{\emph{\color{gray}or}}%
    \extractkey#1\@nil%
    \pgfmanualbody
}
{
  \end{pgfmanualentry}
}

\def\extractkey#1\@nil{%
  \pgfutil@in@={#1}%
  \ifpgfutil@in@%
    \extractkeyequal#1\@nil
  \else%
    \pgfutil@in@{(initial}{#1}%
    \ifpgfutil@in@%
      \extractequalinitial#1\@nil%
    \else
      \pgfmanualentryheadline{{\ttfamily\declarekey{#1}}\hfill(\extrakeytext{}无值)}%
      \def\mykey{#1}%
      \def\mypath{}%
      \def\myname{}%
      \firsttimetrue%
      \decompose#1/\nil%
    \fi
  \fi%
}

\pgfkeyssetvalue{/pgfmanual/gray key prefixes}{/pgfplots/,/tikz/}

\def\declarekey#1{%
    \begingroup
    \gdef\declarekey@didit{0}%
    \pgfkeysgetvalue{/pgfmanual/gray key prefixes}\declarekey@gray
    \def\declarekey@loop{\foreach \prefix in }%
    \expandafter\declarekey@loop\expandafter{\declarekey@gray} {%
        \expandafter\pgfutil@in@\expandafter{\prefix}{#1}%
        \ifpgfutil@in@
            \expandafter\def\expandafter\declarekey@handleprefix\prefix####1\relax{%
                \keyprefix{\prefix}\declare{####1}%
            }%
            \declarekey@handleprefix#1\relax
            \gdef\declarekey@didit{1}%
            \breakforeach
        \fi
    }%
    \if0\declarekey@didit
        \declare{#1}%
    \fi
    \endgroup
}%
\def\keyprefix#1{\textcolor{gray}{#1}}

\def\extractkeyequal#1=#2\@nil{%
  \pgfutil@in@{(default}{#2}%
  \ifpgfutil@in@%
    \extractdefault{#1}#2\@nil%
  \else%
    \pgfutil@in@{(initial}{#2}%
    \ifpgfutil@in@%
      \extractinitial{#1}#2\@nil%
    \else
%     \pgfmanualentryheadline{{\ttfamily\declarekey{#1}=}#2\hfill(\extrakeytext no default)}%
      \pgfmanualentryheadline{{\ttfamily\declarekey{#1}=}#2\hfill}%
    \fi%
  \fi%
  \def\mykey{#1}%
  \def\mypath{}%
  \def\myname{}%
  \firsttimetrue%
  \decompose#1/\nil%
}

\def\extractdefault#1#2(default #3)\@nil{%
  \pgfmanualentryheadline{{\ttfamily\declarekey{#1}\opt{=}}\opt{#2}\hfill (\extrakeytext{}默认 {\ttfamily#3})}%
}

\def\extractinitial#1#2(initially #3)\@nil{%
%  \pgfmanualentryheadline{{\ttfamily\declarekey{#1}=}#2\hfill (\extrakeytext no default, initially {\ttfamily#3})}%
  \pgfmanualentryheadline{{\ttfamily\declarekey{#1}=}#2\hfill (\extrakeytext{}初始 {\ttfamily#3})}%
}

\def\extractequalinitial#1 (initially #2)\@nil{%
  \pgfmanualentryheadline{{\ttfamily\declarekey{#1}}\hfill (\extrakeytext{}初始 {\ttfamily#2})}%
  \def\mykey{#1}%
  \def\mypath{}%
  \def\myname{}%
  \firsttimetrue%
  \decompose#1/\nil%
}

% Introduces a key alias '/#1/<name of current key>'
% to be used inside of \begin{key} ... \end{key}
\def\keyalias#1{\vspace{-3pt}\item{\small 别名 {\ttfamily/#1/\myname}}\vspace{-2pt}\par}

\newif\iffirsttime

\makeatother

\def\decompose/#1/#2\nil{%
  \def\test{#2}%
  \ifx\test\empty%
    % aha.
    \index{#1@\protect\texttt{#1} key}%
    \index{\mypath#1@\protect\texttt{#1}}%
    \def\myname{#1}%
  \else%
    \iffirsttime
%    CF : disabled for /pgfplots/ prefix.
%        \def\mypath{#1@\protect\texttt{/#1/}!}%
%        \firsttimefalse
        \def\pgfplotslocTMPa{pgfplots}%
        \edef\pgfplotslocTMPb{#1}%
        \ifx\pgfplotslocTMPb\pgfplotslocTMPa
            \def\mypath{}%
        \else
            \def\mypath{#1@\protect\texttt{/#1/}!}%
        \fi
        \firsttimefalse
    \else
      \expandafter\def\expandafter\mypath\expandafter{\mypath#1@\protect\texttt{#1/}!}%
    \fi
    \def\firsttime{}
    \decompose/#2\nil%
  \fi%
}

\newenvironment{predefinedmethod}[1]{
  \begin{pgfmanualentry}
    \extractpredefinedmethod#1\@nil
    \pgfmanualbody
}
{
  \end{pgfmanualentry}
}
\def\extractpredefinedmethod#1(#2)\@nil{%
  \pgfmanualentryheadline{方法 \declare{\ttfamily #1}\texttt(#2\texttt) \hfill（为所有类预定义）}
  \index{#1@\protect\texttt{#1} method}%
  \index{Methods!#1@\protect\texttt{#1}}
}


\newenvironment{ooclass}[1]{
  \begin{pgfmanualentry}
    \def\currentclass{#1}
    \pgfmanualentryheadline{类 \declare{\texttt{#1}}}
    \index{#1@\protect\texttt{#1} class}%
    \index{Class #1@Class \protect\texttt{#1}}%
    \index{Classes!#1@\protect\texttt{#1}}
    \pgfmanualbody
}
{
  \end{pgfmanualentry}
}

\newenvironment{method}[1]{
  \begin{pgfmanualentry}
    \extractmethod#1\@nil
    \pgfmanualbody
}
{
  \end{pgfmanualentry}
}
\def\extractmethod#1(#2)\@nil{%
  \def\test{#1}
  \ifx\test\currentclass
    \pgfmanualentryheadline{构造函数 \declare{\ttfamily #1}\texttt(#2\texttt)}
  \else
    \pgfmanualentryheadline{方法 \declare{\ttfamily #1}\texttt(#2\texttt)}
  \fi
  \index{#1@\protect\texttt{#1} method}%
  \index{Methods!#1@\protect\texttt{#1}}
  \index{Class \currentclass!#1@\protect\texttt{#1}}%
}

\iffalse
\newenvironment{attribute}[1]{
  \begin{pgfmanualentry}
    \extractattribute#1\@nil
    \pgfmanualbody
}
{
  \end{pgfmanualentry}
}
\fi
\def\extractattribute#1=#2;\@nil{%
  \def\test{#2}%
  \ifx\test\@empty
    \pgfmanualentryheadline{私有属性 \declare{\ttfamily #1} \hfill（初始为空）}
  \else
    \pgfmanualentryheadline{私有属性 \declare{\ttfamily #1} \hfill（初始为 {\ttfamily #2}）}
  \fi
  \index{#1@\protect\texttt{#1} attribute}%
  \index{Attributes!#1@\protect\texttt{#1}}
  \index{Class \currentclass!#1@\protect\texttt{#1}}%
}



\newenvironment{predefinednode}[1]{
  \begin{pgfmanualentry}
    \pgfmanualentryheadline{预定义节点 {\ttfamily\declare{#1}}}%
    \index{#1@\protect\texttt{#1} node}%
    \index{Predefined node!#1@\protect\texttt{#1}}
    \pgfmanualbody
}
{
  \end{pgfmanualentry}
}

\newenvironment{coordinatesystem}[1]{
  \begin{pgfmanualentry}
    \pgfmanualentryheadline{坐标系 {\ttfamily\declare{#1}}}%
    \index{#1@\protect\texttt{#1} coordinate system}%
    \index{Coordinate systems!#1@\protect\texttt{#1}}
    \pgfmanualbody
    \pgfmanualpdflabel{#1}{}%
}
{
  \end{pgfmanualentry}
}

\newenvironment{snake}[1]{
  \begin{pgfmanualentry}
    \pgfmanualentryheadline{蛇形 {\ttfamily\declare{#1}}}%
    \index{#1@\protect\texttt{#1} snake}%
    \index{Snakes!#1@\protect\texttt{#1}}
    \pgfmanualbody
}
{
  \end{pgfmanualentry}
}

\newenvironment{decoration}[1]{
  \begin{pgfmanualentry}
    \pgfmanualentryheadline{装饰 {\ttfamily\declare{#1}}}%
    \index{#1@\protect\texttt{#1} decoration}%
    \index{Decorations!#1@\protect\texttt{#1}}
    \pgfmanualbody
}
{
  \end{pgfmanualentry}
}


\def\pgfmanualbar{\char`\|}
\makeatletter
\newenvironment{pathoperation}[3][]{
  \begin{pgfmanualentry}
    \pgfmanualentryheadline{\textcolor{gray}{{\ttfamily\char`\\path}\
        \ \dots}
      \declare{\texttt{#2}}#3\ \textcolor{gray}{\dots\texttt{;}}}%
    \def\pgfmanualtest{#1}%
    \ifx\pgfmanualtest\@empty%
      \index{#2@\protect\texttt{#2} path operation}%
      \index{Path operations!#2@\protect\texttt{#2}}%
      \pgfmanualpdflabel{#2}{}%
    \fi%
    \pgfmanualbody
}
{
  \end{pgfmanualentry}
}
\makeatother

\def\extractcommand#1#2\@@{%
  \pgfmanualentryheadline{\declare{\texttt{\string#1}}#2}%
  \removeats{#1}%
  \index{\strippedat @\protect\myprintocmmand{\strippedat}}}

\def\doublebs{\texttt{\char`\\\char`\\}}


\newenvironment{package}[1]{
  \begin{pgfmanualentry}
    \pgfmanualentryheadline{{\ttfamily\char`\\usepackage\char`\{\declare{#1}\char`\}\space\space \char`\%\space\space  \LaTeX}}
    \index{#1@\protect\texttt{#1} package}%
    \index{Packages and files!#1@\protect\texttt{#1}}%
    \pgfmanualentryheadline{{\ttfamily\char`\\input \declare{#1}.tex\space\space\space \char`\%\space\space  plain \TeX}}
    \pgfmanualentryheadline{{\ttfamily\char`\\usemodule[\declare{#1}]\space\space \char`\%\space\space  Con\TeX t}}
    \pgfmanualbody
}
{
  \end{pgfmanualentry}
}


\newenvironment{pgfmodule}[1]{
  \begin{pgfmanualentry}
    \pgfmanualentryheadline{{\ttfamily\char`\\usepgfmodule\char`\{\declare{#1}\char`\}\space\space\space
        \char`\%\space\space  \LaTeX\space and plain \TeX\space and pure pgf}}
    \index{#1@\protect\texttt{#1} module}%
    \index{Modules!#1@\protect\texttt{#1}}%
    \pgfmanualentryheadline{{\ttfamily\char`\\usepgfmodule[\declare{#1}]\space\space \char`\%\space\space  Con\TeX t\space and pure pgf}}
    \pgfmanualbody
}
{
  \end{pgfmanualentry}
}

\newenvironment{pgflibrary}[1]{
  \begin{pgfmanualentry}
    \pgfmanualentryheadline{{\ttfamily\char`\\usepgflibrary\char`\{\declare{#1}\char`\}\space\space\space
        \char`\%\space\space  \LaTeX\space and plain \TeX\space and pure pgf}}
    \index{#1@\protect\texttt{#1} library}%
    \index{Libraries!#1@\protect\texttt{#1}}%
    \pgfmanualentryheadline{{\ttfamily\char`\\usepgflibrary[\declare{#1}]\space\space \char`\%\space\space  Con\TeX t\space and pure pgf}}
    \pgfmanualentryheadline{{\ttfamily\char`\\usetikzlibrary\char`\{\declare{#1}\char`\}\space\space
        \char`\%\space\space  \LaTeX\space and plain \TeX\space when using \tikzname}}
    \pgfmanualentryheadline{{\ttfamily\char`\\usetikzlibrary[\declare{#1}]\space
        \char`\%\space\space  Con\TeX t\space when using \tikzname}}
    \pgfmanualbody
}
{
  \end{pgfmanualentry}
}

\newenvironment{tikzlibrary}[1]{
  \begin{pgfmanualentry}
    \pgfmanualentryheadline{{\ttfamily\char`\\usetikzlibrary\char`\{\declare{#1}\char`\}\space\space \char`\%\space\space  \LaTeX\space and plain \TeX}}
    \index{#1@\protect\texttt{#1} library}%
    \index{Libraries!#1@\protect\texttt{#1}}%
    \pgfmanualentryheadline{{\ttfamily\char`\\usetikzlibrary[\declare{#1}]\space \char`\%\space\space Con\TeX t}}
    \pgfmanualbody
}
{
  \end{pgfmanualentry}
}



\newenvironment{filedescription}[1]{
  \begin{pgfmanualentry}
    \pgfmanualentryheadline{文件 {\ttfamily\declare{#1}}}%
    \index{#1@\protect\texttt{#1} file}%
    \index{Packages and files!#1@\protect\texttt{#1}}%
    \pgfmanualbody
}
{
  \end{pgfmanualentry}
}


\newenvironment{packageoption}[1]{
  \begin{pgfmanualentry}
    \pgfmanualentryheadline{{\ttfamily\char`\\usepackage[\declare{#1}]\char`\{pgf\char`\}}}
    \index{#1@\protect\texttt{#1} package option}%
    \index{Package options for \textsc{pgf}!#1@\protect\texttt{#1}}%
    \pgfmanualbody
}
{
  \end{pgfmanualentry}
}



\newcommand\opt[1]{{\color{black!50!green}#1}}
\newcommand\ooarg[1]{{\ttfamily[}\meta{#1}{\ttfamily]}}

\def\opt{\afterassignment\pgfmanualopt\let\next=}
\def\pgfmanualopt{\ifx\next\bgroup\bgroup\color{black!50!green}\else{\color{black!50!green}\next}\fi}



\def\beamer{\textsc{beamer}}
\def\pdf{\textsc{pdf}}
\def\pgfname{\textsc{pgf}}
\def\tikzname{Ti\emph{k}Z}
\def\pstricks{\textsc{pstricks}}
\def\prosper{\textsc{prosper}}
\def\seminar{\textsc{seminar}}
\def\texpower{\textsc{texpower}}
\def\foils{\textsc{foils}}

{
  \makeatletter
  \global\let\myempty=\@empty
  \global\let\mygobble=\@gobble
  \catcode`\@=12
  \gdef\getridofats#1@#2\relax{%
    \def\getridtest{#2}%
    \ifx\getridtest\myempty%
      \expandafter\def\expandafter\strippedat\expandafter{\strippedat#1}
    \else%
      \expandafter\def\expandafter\strippedat\expandafter{\strippedat#1\protect\printanat}
      \getridofats#2\relax%
    \fi%
  }

  \gdef\removeats#1{%
    \let\strippedat\myempty%
    \edef\strippedtext{\stripcommand#1}%
    \expandafter\getridofats\strippedtext @\relax%
  }

  \gdef\stripcommand#1{\expandafter\mygobble\string#1}
}

\def\printanat{\char`\@}

\def\declare{\afterassignment\pgfmanualdeclare\let\next=}
\def\pgfmanualdeclare{\ifx\next\bgroup\bgroup\color{red!75!black}\else{\color{red!75!black}\next}\fi}


\let\textoken=\command
\let\endtextoken=\endcommand

\def\myprintocmmand#1{\texttt{\char`\\#1}}

\def\example{\par\smallskip\noindent\textit{例： }}
\def\themeauthor{\par\smallskip\noindent\textit{主题作者： }}


\def\indexoption#1{%
  \index{#1@\protect\texttt{#1} option}%
  \index{Graphic options and styles!#1@\protect\texttt{#1}}%
}

\def\itemcalendaroption#1{\item \declare{\texttt{#1}}%
  \index{#1@\protect\texttt{#1} date test}%
  \index{Date tests!#1@\protect\texttt{#1}}%
}



\def\class#1{\list{}{\leftmargin=2em\itemindent-\leftmargin\def\makelabel##1{\hss##1}}%
\extractclass#1@\par\topsep=0pt}
\def\endclass{\endlist}
\def\extractclass#1#2@{%
\item{{{\ttfamily\char`\\documentclass}#2{\ttfamily\char`\{\declare{#1}\char`\}}}}%
  \index{#1@\protect\texttt{#1} class}%
  \index{Classes!#1@\protect\texttt{#1}}}

\def\partname{部分}

\makeatletter
\def\index@prologue{\section*{索引}\addcontentsline{toc}{section}{索引}
  本索引仅包含自动生成的条目。一份好的索引还应包含经过仔细挑选的关键词。本索引并不是一份好的索引。
  \bigskip
}
\c@IndexColumns=2
  \def\theindex{\@restonecoltrue
    \columnseprule \z@  \columnsep 29\p@
    \twocolumn[\index@prologue]%
       \parindent -30pt
       \columnsep 15pt
       \parskip 0pt plus 1pt
       \leftskip 30pt
       \rightskip 0pt plus 2cm
       \small
       \def\@idxitem{\par}%
    \let\item\@idxitem \ignorespaces}
  \def\endtheindex{\onecolumn}
\def\noindexing{\let\index=\@gobble}



\newcommand\symarrow[1]{
  \index{#1@\protect\texttt{#1} arrow tip}%
  \index{Arrow tips!#1@\protect\texttt{#1}}
  \texttt{#1}& 产生粗线
  \begin{tikzpicture}[arrows={#1-#1},thick,baseline]
    \useasboundingbox (0pt,-0.5ex) rectangle (1cm,2ex);
    \draw (0pt,.5ex) -- (1cm,.5ex);
  \end{tikzpicture} 和细线
  \begin{tikzpicture}[arrows={#1-#1},thin,baseline]
    \useasboundingbox (0pt,-0.5ex) rectangle (1cm,2ex);
    \draw (0pt,.5ex) -- (1cm,.5ex);
  \end{tikzpicture}
}
\newcommand\symarrowdouble[1]{
  \index{#1@\protect\texttt{#1} arrow tip}%
  \index{Arrow tips!#1@\protect\texttt{#1}}
  \texttt{#1}& 产生粗线
  \begin{tikzpicture}[arrows={#1-#1},thick,baseline]
    \useasboundingbox (0pt,-0.5ex) rectangle (1cm,2ex);
    \draw (0pt,.5ex) -- (1cm,.5ex);
  \end{tikzpicture}
  和细线
  \begin{tikzpicture}[arrows={#1-#1},thin,baseline]
    \useasboundingbox (0pt,-0.5ex) rectangle (1cm,2ex);
    \draw (0pt,.5ex) -- (1cm,.5ex);
  \end{tikzpicture}，双线
  \begin{tikzpicture}[arrows={#1-#1},thick,baseline]
    \useasboundingbox (0pt,-0.5ex) rectangle (1cm,2ex);
    \draw[double,double equal sign distance] (0pt,.5ex) -- (1cm,.5ex);
  \end{tikzpicture} 和
  \begin{tikzpicture}[arrows={#1-#1},thin,baseline]
    \useasboundingbox (0pt,-0.5ex) rectangle (1cm,2ex);
    \draw[double,double equal sign distance] (0pt,.5ex) -- (1cm,.5ex);
  \end{tikzpicture}
}

\newcommand\sarrow[2]{
  \index{#1@\protect\texttt{#1} arrow tip}%
  \index{Arrow tips!#1@\protect\texttt{#1}}
  \index{#2@\protect\texttt{#2} arrow tip}%
  \index{Arrow tips!#2@\protect\texttt{#2}}
  \texttt{#1-#2}& 产生粗线
  \begin{tikzpicture}[arrows={#1-#2},thick,baseline]
    \useasboundingbox (0pt,-0.5ex) rectangle (1cm,2ex);
    \draw (0pt,.5ex) -- (1cm,.5ex);
  \end{tikzpicture} 和细线
  \begin{tikzpicture}[arrows={#1-#2},thin,baseline]
    \useasboundingbox (0pt,-0.5ex) rectangle (1cm,2ex);
    \draw (0pt,.5ex) -- (1cm,.5ex);
  \end{tikzpicture}
}

\newcommand\sarrowdouble[2]{
  \index{#1@\protect\texttt{#1} arrow tip}%
  \index{Arrow tips!#1@\protect\texttt{#1}}
  \index{#2@\protect\texttt{#2} arrow tip}%
  \index{Arrow tips!#2@\protect\texttt{#2}}
  \texttt{#1-#2}& 产生粗线
  \begin{tikzpicture}[arrows={#1-#2},thick,baseline]
    \useasboundingbox (0pt,-0.5ex) rectangle (1cm,2ex);
    \draw (0pt,.5ex) -- (1cm,.5ex);
  \end{tikzpicture} 和细线
  \begin{tikzpicture}[arrows={#1-#2},thin,baseline]
    \useasboundingbox (0pt,-0.5ex) rectangle (1cm,2ex);
    \draw (0pt,.5ex) -- (1cm,.5ex);
  \end{tikzpicture}，双线
  \begin{tikzpicture}[arrows={#1-#2},thick,baseline]
    \useasboundingbox (0pt,-0.5ex) rectangle (1cm,2ex);
    \draw[double,double equal sign distance] (0pt,.5ex) -- (1cm,.5ex);
  \end{tikzpicture} 和
  \begin{tikzpicture}[arrows={#1-#2},thin,baseline]
    \useasboundingbox (0pt,-0.5ex) rectangle (1cm,2ex);
    \draw[double,double equal sign distance] (0pt,.5ex) -- (1cm,.5ex);
  \end{tikzpicture}
}

\newcommand\carrow[1]{
  \index{#1@\protect\texttt{#1} arrow tip}%
  \index{Arrow tips!#1@\protect\texttt{#1}}
  \texttt{#1}& 对线宽 1ex 产生
  \begin{tikzpicture}[arrows={#1-#1},line width=1ex,baseline]
    \useasboundingbox (0pt,-0.5ex) rectangle (1.5cm,2ex);
    \draw (0pt,.5ex) -- (1.5cm,.5ex);
  \end{tikzpicture}
}
\def\myvbar{\char`\|}
\newcommand\plotmarkentry[1]{%
  \index{#1@\protect\texttt{#1} plot mark}%
  \index{Plot marks!#1@\protect\texttt{#1}}
  \texttt{\char`\\pgfuseplotmark\char`\{\declare{#1}\char`\}} &
  \tikz\draw[color=black!25] plot[mark=#1,mark options={fill=examplefill,draw=black}] coordinates{(0,0) (.5,0.2) (1,0) (1.5,0.2)};\\
}
\newcommand\plotmarkentrytikz[1]{%
  \index{#1@\protect\texttt{#1} plot mark}%
  \index{Plot marks!#1@\protect\texttt{#1}}
  \texttt{mark=\declare{#1}} & \tikz\draw[color=black!25]
  plot[mark=#1,mark options={fill=examplefill,draw=black}]
    coordinates {(0,0) (.5,0.2) (1,0) (1.5,0.2)};\\
}



\ifx\scantokens\@undefined
  \PackageError{pgfmanual-macros}{You need to use extended latex
    (elatex) or (pdfelatex) to process this document}{}
\fi

% Define \find@example such that it doesn't destroy catcodes:
\begingroup
\catcode`|=0
\catcode`[= 1
\catcode`]=2
\catcode`\{=12
\catcode `\}=12
\catcode`\\=12 |gdef|find@example#1\end{codeexample}[|endofcodeexample[#1]]
|endgroup

% define \returntospace.
%
% It should define NEWLINE as {}, spaces and tabs as \space.
\begingroup
\catcode`\^=7
\catcode`\^^M=13
\catcode`\^^I=13
\catcode`\ =13%
\gdef\returntospace{\catcode`\ =13\def {\space}\catcode`\^^I=13\def^^I{\space}\catcode`\^^M=13\def^^M{\pgfkeysvalueof{/codeexample/newline}}}%
\endgroup

\begingroup
\catcode`\%=13
\catcode`\^^M=13
\gdef\commenthandler{\catcode`\%=13\def%{\@gobble@till@return}}
\gdef\@gobble@till@return#1^^M{}
\gdef\@gobble@till@return@ignore#1^^M{\ignorespaces}
\gdef\typesetcomment{\catcode`\%=13\def%{\@typeset@till@return}}
\gdef\@typeset@till@return#1^^M{{\def%{\char`\%}\pgfkeysvalueof{/codeexample/comment font}{\char`\%#1}}\par}
\endgroup

% Define tab-implementation functions
%   \codeexample@tabinit@replacementchars@
% and
%   \codeexample@tabinit@catcode@
%
% They should ONLY be used in case that tab replacement is active.
%
% This here is merely a preparation step.
%
% Idea:
% \codeexample@tabinit@catcode@ will make TAB active
% and
% \codeexample@tabinit@replacementchars@ will insert as many spaces as
% /codeexample/tabsize contains.
{
\catcode`\^^I=13
% ATTENTION: do NOT use tabs in these definitions!!
\gdef\codeexample@tabinit@replacementchars@{%
 \begingroup
 \count0=\pgfkeysvalueof{/codeexample/tabsize}\relax
 \toks0={}%
 \loop
 \ifnum\count0>0
  \advance\count0 by-1
  \toks0=\expandafter{\the\toks0\ }%
 \repeat
 \xdef\codeexample@tabinit@replacementchars@@{\the\toks0}%
 \endgroup
 \let^^I=\codeexample@tabinit@replacementchars@@
}%
\gdef\codeexample@tabinit@catcode@{\catcode`\^^I=13}%
}%

% Called after any options have been set. It assigns
%   \codeexample@tabinit@catcode
% and
%   \codeexample@tabinit@replacementchars
% which are used inside of
%\begin{codeexample}
% ...
%\end{codeexample}
%
% \codeexample@tabinit@catcode  is either \relax or it makes tab
% active.
%
% \codeexample@tabinit@replacementchars is either \relax or it inserts
% a proper replacement sequence for tabs (as many spaces as
% configured)
\def\codeexample@tabinit{%
    \ifnum\pgfkeysvalueof{/codeexample/tabsize}=0\relax
        \let\codeexample@tabinit@replacementchars=\relax
        \let\codeexample@tabinit@catcode=\relax
    \else
        \let\codeexample@tabinit@catcode=\codeexample@tabinit@catcode@
        \let\codeexample@tabinit@replacementchars=\codeexample@tabinit@replacementchars@
    \fi
}

\pgfqkeys{/codeexample}{%
    width/.code=    {\setlength\codeexamplewidth{#1}},
    graphic/.code=    {\colorlet{graphicbackground}{#1}},
    code/.code=    {\colorlet{codebackground}{#1}},
    execute code/.is if=code@execute,
    code only/.code=    {\code@executefalse},
    pre/.code=    {\def\code@pre{#1}},
    post/.code=    {\def\code@post{#1}},
    % #1 is the *complete* environment contents as it shall be
    % typeset. In particular, the catcodes are NOT the normal ones.
    typeset listing/.code=    {#1},
    vbox/.code=    {\def\code@pre{\vbox\bgroup\setlength{\hsize}{\linewidth-6pt}}\def\code@post{\egroup}},
    ignorespaces/.code=    {\let\@gobble@till@return=\@gobble@till@return@ignore},
    leave comments/.code=    {\def\code@catcode@hook{\catcode`\%=12}\let\commenthandler=\relax\let\typesetcomment=\relax},
    comment font/.initial=\textsl,% FIXME : doesn't work!?
    tabsize/.initial=0,
    newline/.initial=^^J,% use 'newline=^^J' to insert a *real* newline
    every codeexample/.style={width=4cm+7pt},
}

\def\code@pre{}
\def\code@post{}
\def\code@catcode@hook{}

\newdimen\codeexamplewidth
\newif\ifcode@execute
\newbox\codeexamplebox
\def\codeexample[#1]{%
  \begingroup%
  \code@executetrue
  \pgfqkeys{/codeexample}{every codeexample,#1}%
  \codeexample@tabinit% assigns \codeexample@tabinit@[catcode,replacementchars]
  \parindent0pt
  \begingroup%
  \par%
  \medskip%
  \let\do\@makeother%
  \dospecials%
  \obeylines%
  \@vobeyspaces%
  \catcode`\%=13%
  \catcode`\^^M=13%
  \code@catcode@hook%
  \codeexample@tabinit@catcode
  \relax%
  \find@example}
\def\endofcodeexample#1{%
  \endgroup%
  \ifcode@execute%
    \setbox\codeexamplebox=\hbox{%
      {%
        {%
          \returntospace%
          \commenthandler%
          \xdef\code@temp{#1}% removes returns and comments
        }%
%\message{Nach returntospace: \meaning\code@temp^^J^^J^^J}%
        \expandafter\pgfmanual@codeexample@generate@graphic\expandafter{\code@temp}%
      }%
    }%
    \ifdim\wd\codeexamplebox>\codeexamplewidth%
      \def\code@start{\par}%
      \def\code@flushstart{}\def\code@flushend{}%
      \def\code@mid{\parskip2pt\par\noindent}%
      \def\code@width{\linewidth-6pt}%
      \def\code@end{}%
    \else%
      \def\code@start{%
        \linewidth=\textwidth%
        \parshape \@ne 0pt \linewidth
        \leavevmode%
        \hbox\bgroup}%
      \def\code@flushstart{\hfill}%
      \def\code@flushend{\hbox{}}%
      \def\code@mid{\hskip6pt}%
      \def\code@width{\linewidth-12pt-\codeexamplewidth}%
      \def\code@end{\egroup}%
    \fi%
    \code@start%
    \noindent%
    \begin{minipage}[t]{\codeexamplewidth}\raggedright
      \hrule width0pt%
      \footnotesize\vskip-1em%
      \code@flushstart\box\codeexamplebox\code@flushend%
      \vskip-1ex
      \leavevmode%
    \end{minipage}%
  \else%
    \def\code@mid{\par}
    \def\code@width{\linewidth-6pt}
    \def\code@end{}
  \fi%
  \code@mid%
  \colorbox{codebackground}{%
    \begin{minipage}[t]{\code@width}%
      {%
        \let\do\@makeother
        \dospecials
        \frenchspacing\@vobeyspaces
        \normalfont\ttfamily\footnotesize
        \typesetcomment%
        \codeexample@tabinit@replacementchars
        \@tempswafalse
        \def\par{%
          \if@tempswa
          \leavevmode \null \@@par\penalty\interlinepenalty
          \else
          \@tempswatrue
          \ifhmode\@@par\penalty\interlinepenalty\fi
          \fi}%
        \obeylines
        \everypar \expandafter{\the\everypar \unpenalty}%
        \pgfmanual@@typeset@listing{#1}%
    }
    \end{minipage}}%
  \code@end%
  \par%
  \medskip
  \end{codeexample}
}

% the purpose of these extra-macros is to plug-in special tex4ht
% instructions for the html version:
\long\def\pgfmanual@codeexample@generate@graphic#1{%
    \colorbox{graphicbackground}{\pgfmanual@codeexample@generate@graphic@{#1}}%
}%
\long\def\pgfmanual@codeexample@generate@graphic@#1{%
    \color{black}\ignorespaces%
    \code@pre%
    \pgfmanual@scantokens{#1}\ignorespaces%
    \code@post\ignorespaces%
}%

\ifluatex
    % LUA has a couple of bugs around \scantokens, see
    % http://tex.stackexchange.com/questions/197634/workaround-for-open-lualatex-bug-in-scantokens
    \def\pgfmanual@scantokens#1{
        \directlua{%
            local s = "\luaescapestring{\code@temp}"
            local splitted = pgfplots.stringsplit(s, "\string\n")
            for i =1,\pgfplotsHASH splitted do
                local line=splitted[i]
              tex.print(line)
            end
        }%
        %\def\pgfmanual@tempfile{\jobname.code.tmp}%
        %\immediate\openout\w@pgf@writea=\pgfmanual@tempfile\relax
        %\immediate\write\w@pgf@writea{#1}%
        %\immediate\closeout\w@pgf@writea
        %\input{\pgfmanual@tempfile}
    }%
\else
    \let\pgfmanual@scantokens=\scantokens
\fi

\long\def\pgfmanual@@typeset@listing#1{%
    \pgfkeysvalueof{/codeexample/typeset listing/.@cmd}{#1}\pgfeov
}%%

\def\endcodeexample{\endgroup}


\makeatother


%%% Local Variables:
%%% mode: latex
%%% TeX-master: "beameruserguide"
%%% End:



\def\pgfplotsifdocpackageuptodate#1#2{%
    \pgfkeysifdefined{/codeexample/prettyprint/word/.@cmd}{#1}{#2}
}%

\pgfplotsiffileexists{pgfmanual.sty}{%
    \RequirePackage{pgfmanual}
    \pgfplotsifdocpackageuptodate{}{%
        \makeatletter
        \input{pgfplotsoldpgfsupp_pgfmanual.code.tex}
        \makeatother
    }%
}{%
    \makeatletter
    \input{pgfplotsoldpgfsupp_pgfmanual.code.tex}
    \makeatother
}%

\makeatletter
\def\pgfplotsmakefilelinkifuseful#1#2{%
    \protect\pgfplotsmakefilelinkifuseful@{#1}{#2}%
}%
\def\pgfplotsmakefilelinkifuseful@#1#2{%
    \edef\temp{#1}%
    \edef\tempb{\jobname}%
    \edef\temp{\meaning\temp}% \meaning normalizes the catcodes.
    \edef\tempb{\meaning\tempb}%
    \ifx\temp\tempb
        % we are processing '#1'. Don't make a link.
        #2%
    \else
        \href{file:#1.pdf}{#2}%
    \fi
}%
\makeatother


\pgfkeys{
    /codeexample/prettyprint/cs arguments/pgfplotscreateplotcyclelist/.initial=2,
    /codeexample/prettyprint/cs/pgfplotscreateplotcyclelist/.code args={#1#2#3}{\pgfmanualpdfref{#1}{#1}\{#2\}\{\pgfmanualprettyprintpgfkeys{#3}\pgfmanualclosebrace},
    /codeexample/prettyprint/cs arguments/tikzset/.initial=1,
    /codeexample/prettyprint/cs/tikzset/.code 2 args={\pgfmanualpdfref{#1}{#1}\{\pgfmanualprettyprintpgfkeys{#2}\pgfmanualclosebrace},
    /codeexample/prettyprint/cs arguments/pgfplotsset/.initial=1,
    /codeexample/prettyprint/cs/pgfplotsset/.code 2 args={\pgfmanualpdfref{#1}{#1}\{\pgfmanualprettyprintpgfkeys{#2}\pgfmanualclosebrace},
    /codeexample/prettyprint/cs arguments/pgfplotstableset/.initial=1,
    /codeexample/prettyprint/cs/pgfplotstableset/.code 2 args={\pgfmanualpdfref{#1}{#1}\{\pgfmanualprettyprintpgfkeys{#2}\pgfmanualclosebrace},
    /codeexample/prettyprint/cs arguments/usepgfplotslibrary/.initial=1,
    /codeexample/prettyprint/cs/usepgfplotslibrary/.code 2 args={\pgfmanualpdfref{#1}{#1}\{\pgfmanualpdfref{#2}{#2}\pgfmanualclosebrace},
    %
    %
    %/codeexample/prettyprint/key value/cycle list/.code 2 args={\pgfmanualprettyprintpgfkeys{#2}},
    /codeexample/prettyprint/key value/xticklabel/.code 2 args={\pgfmanualprettyprintcode{#2}},
    /codeexample/prettyprint/key value/yticklabel/.code 2 args={\pgfmanualprettyprintcode{#2}},
    /codeexample/prettyprint/key value/zticklabel/.code 2 args={\pgfmanualprettyprintcode{#2}},
    /codeexample/prettyprint/key value/nodes near coords/.code 2 args={\pgfmanualprettyprintcode{#2}},
    /codeexample/prettyprint/key value/nodes near coords*/.code 2 args={\pgfmanualprettyprintcode{#2}},
    %
    /codeexample/prettyprint/key value/includegraphics/.code 2 args={\pgfmanualprettyprintpgfkeys{#2}},
    /codeexample/prettyprint/key value/contour gnuplot/.code 2 args={\pgfmanualprettyprintpgfkeys{#2}},
    /codeexample/prettyprint/key value/contour lua/.code 2 args={\pgfmanualprettyprintpgfkeys{#2}},
    /codeexample/prettyprint/key value/contour prepared/.code 2 args={\pgfmanualprettyprintpgfkeys{#2}},
    /codeexample/prettyprint/key value/contour filled/.code 2 args={\pgfmanualprettyprintpgfkeys{#2}},
    /codeexample/prettyprint/key value/hist/.code 2 args={\pgfmanualprettyprintpgfkeys{#2}},
    /codeexample/prettyprint/key value/tieline/.code 2 args={\pgfmanualprettyprintpgfkeys{#2}},
    /codeexample/prettyprint/key value/quiver/.code 2 args={\pgfmanualprettyprintpgfkeys{#2}},
    /codeexample/prettyprint/key value/boxplot/.code 2 args={\pgfmanualprettyprintpgfkeys{#2}},
    /codeexample/prettyprint/key value/boxplot prepared/.code 2 args={\pgfmanualprettyprintpgfkeys{#2}},
    %
    /codeexample/prettyprint/key value/colormap/.code 2 args={\pgfmanualprettyprintpgfkeys{#2}},
    /codeexample/prettyprint/key value/of colormap/.code 2 args={\pgfmanualprettyprintpgfkeys{#2}},
    /codeexample/prettyprint/key value/samples of colormap/.code 2 args={\pgfmanualprettyprintpgfkeys{#2}},
    /codeexample/prettyprint/key value/levels from colormap/.code 2 args={\pgfmanualprettyprintpgfkeys{#2}},
    /codeexample/prettyprint/key value/colors of colormap/.code 2 args={\pgfmanualprettyprintpgfkeys{#2}},
    %
    %
    % whenever an unqualified key is found, the following key prefix
    % list is tried to find a match.
    /pdflinks/search key prefixes in={/pgfplots/table/,/pgfplots/error bars/,/pgfplots/,/pgfplots/plot file/,/tikz/,/pgf/},
    %
    % the link prefix written to the pdf file:
    /pdflinks/internal link prefix=pgfp,
    %
    /pdflinks/warnings=false,
    /pdflinks/codeexample links=true,
    /pdflinks/show labels=false,
}%


% should be used to show something in red which doesn't need to get a
% hyper ref.
%
% Examples are descriptions of key labels.
\def\declaretext#1{\texttt{\declare{#1}}}

% To be used whenever something NEW has been declared.
% In this case, a \pgfmanualpdflabel will be generated using '#1'.
%
% Use '\declaretext' if you only describe something local (for example
% the documentation of key values).
\def\declarelabel#1{%
    \texttt{\declare{#1}}%
    \pgfmanualpdflabel{#1}{}%
}

\def\pgfmanualbar{\char`\|}
\makeatletter

%%%%%%%%%%%%%%%%%%%%%%%%%%%%%%%%%%%%%%%%%%%%%%%%%%%%%%%%%%%%%%%%%%%%%%%%%%%%%%

\usepgfplotslibrary{external}

% use \pgfplotsmanualenableexternalizationofexpensive in the preamble
% to enable externalization of expensive examples.
% This will ONLY externalize expensive examples, i.e. those for which
% \pgfplotsexpensiveexample is written:
%
% \pgfplotsexpensiveexample
% \begin{codeexample}
\def\pgfplotsmanualenableexternalizationofexpensive{%
    \pgfplotsmanual@enable@externalization@for@expensivetrue
    \tikzexternalize[
        prefix=figures/expensiveexamples/,% the 'X' suffix is to avoid confusion in git: previous versions contained the pdfs in git
        figure name={},
        export=false, % needs to be activated for single pictures (i.e. expensive ones)
        mode=list and make,
        verbose=false,
        %xport=true,% FASTER FOR DEBUGGING
    ]
    \tikzifexternalizing{%
        \nofiles
        \pgfkeys{/pdflinks/codeexample links=false}%
    }{}%
}%
%\pgfkeys{/pgf/images/include external/.code={\href{file:#1}{\pgfimage{#1}}}} % FIXME : NOT FOR THE FINAL VERSION

\newif\ifpgfplotsmanual@enable@externalization@for@expensive
\newif\ifpgfplots@example@is@expensive

\pgfkeys{
    /codeexample/every codeexample/.append code={%
        \ifpgfplots@example@is@expensive
            \pgfkeys{/tikz/external/export=true}%
            \global\pgfplots@example@is@expensivefalse
        \fi
    }
}

% Write this macro directly in front of \begin{codeexample} (without arguments):
\def\pgfplotsexpensiveexample{%
    \ifpgfplotsmanual@enable@externalization@for@expensive
        \pgfplots@example@is@expensivetrue
    \else
        \message{[NOTE: I am now about to typeset an expensive example. You will need to ENLARGE YOUR TeX MEMORY CAPACITIES if this fails.]}%
    \fi
}%

%%%%%%%%%%%%%%%%%%%%%%%%%%%%%%%%%%%%%%%%%%%%%%%%%%%%%%%%%%%%%%%%%%%%%%%%%%%%%%

\newif\ifpgfplotsmanualhtmlmode

% MACROS FOR HTML OUTPUT:
% HTML will be compiled in a separate folder. See the Makefile
\tikzset{
    external/html export mode/.style={
        system call={%
        % FIXME : which converter to png ? Why did I use pdftoppm |
        % pnmtopng in the first draft?
            %pdflatex \tikzexternalcheckshellescape -halt-on-error -interaction=batchmode -jobname "\image" "\\def\\pgfplotsmanualhtlatexmode{1}\texsource" && pdftoppm "\image.pdf" | pnmtopng > "\image.png"%
            pdflatex \tikzexternalcheckshellescape -halt-on-error -interaction=batchmode -jobname "\image" "\\def\\pgfplotsmanualhtlatexmode{1}\texsource" && pdftoppm -png "\image.pdf" "\image"%
            %pdflatex \tikzexternalcheckshellescape -halt-on-error -interaction=batchmode -jobname "\image" "\\def\\pgfplotsmanualhtlatexmode{1}\texsource" && convert -density 150 "\image.pdf" "\image.png"%
        },
        /pgf/images/external info,
        /pgf/images/include external/.code={%
            \includegraphics
               [width=\pgfexternalwidth,height=\pgfexternalheight]
               {##1.png}%
        },
        prefix=figures/generated/,
        figure name={manual},
        mode=list and make,
        verbose=false,
    },
}

\def\pgfplotsmanual@configure@for@htlatexmode{%
    % disable this; we will externalize everything now:
    \let\pgfplotsmanualenableexternalizationofexpensive=\relax
    %
    \message{^^Jpgfplots manual: initializing external lib for HTML output (png)^^J}%
    \tikzexternalize[
        html export mode,
    ]
    %
    \@ifpackageloaded{tex4ht}{%
        %\usepackage[html,png,3]{tex4ht}
        \tikzset{external/mode=only graphics}%
        \pgfplotsmanual@configure@for@htlatexmode@
        %\patches@for@htlatex
    }{%
    }%
    %
    %
    \pgfplotsmanualhtmlmodetrue
    %
    \gdef\pgfsys@imagesuffixlist{.png}
    %
    \tikzifexternalizing{%
        %\nofiles
        \pgfkeys{/pdflinks/codeexample links=false}%
    }{}%
}

% explanation of how to customize tex4ht:
% http://www.cvr.cc/tex4ht-low-level-commands/#more-482
%
% some reference:
% http://www.cvr.cc/tex4ht-options/
%
% some tex4ht mailing list discussion:
% http://tug.org/pipermail/tex4ht/2010q4/000246.html
\def\patches@for@htlatex{%
 \let\savecolor\color
 \NewConfigure{color}[2]{\def\a at color{##1}\def\b at color{##2}}
 \def\@@tmp##1{\a at color##1\b at color\savecolor{##1}\aftergroup\endspan}
 \let\color\@@tmp
 \def\endspan{\Tg</span>}
 \Configure{color}{\HCode{<span style="color:}}{\HCode{;">}}
}%

\def\pgfplotsmanual@configure@for@htlatexmode@{%
    \NewConfigure{endofcodeexample}[2]{%
        \def\endofcodeexample@pre{##1}%
        \def\endofcodeexample@post{##2}%
    }%
    \let\endofcodeexample@orig=\endofcodeexample
    \def\endofcodeexample##1{%
        \endofcodeexample@pre
        \endofcodeexample@orig{##1}%
        \endofcodeexample@post
    }%
    \Configure{endofcodeexample}
        {\Tg<div class="img">}
        {\Tg</div>}%
    %
    \NewConfigure{pgfmanual@@typeset@listing}[2]{%
        \def\pgfmanual@@typeset@listing@pre{##1}%
        \def\pgfmanual@@typeset@listing@post{##2}%
    }%
    \let\pgfmanual@@typeset@listing@orig=\pgfmanual@@typeset@listing
    \def\pgfmanual@@typeset@listing##1{%
        \pgfmanual@@typeset@listing@pre
        \pgfmanual@@typeset@listing@orig{##1}%
        \pgfmanual@@typeset@listing@post
    }%
    \Configure{pgfmanual@@typeset@listing}
        {\Tg<div class="texsrc">}
        {\Tg</div>}%
    %
    %
    \NewConfigure{pgfmanual@codeexample@generate@graphic}[2]{%
        \def\pgfmanual@codeexample@generate@graphic@pre{##1}%
        \def\pgfmanual@codeexample@generate@graphic@post{##2}%
    }%
    \def\pgfmanual@codeexample@generate@graphic##1{%
        \pgfmanual@codeexample@generate@graphic@pre
        % the '@' variant skips the colorbox:
        \pgfmanual@codeexample@generate@graphic@{##1}%
        \pgfmanual@codeexample@generate@graphic@post
    }%
    \Configure{pgfmanual@codeexample@generate@graphic}
        {\Tg<span class="graphics">}
        {\Tg</span>}%
}

\@ifpackageloaded{tex4ht}{
    % always active here:
    \def\pgfplotsmanualhtlatexmode{1}%
}{
}

% the macro \pgfplotsmanualhtlatexmode should be set from command line
% (to any value). If it is known, the manual will be translated to
% HTML.
\pgfutil@ifundefined{pgfplotsmanualhtlatexmode}{%
}{%
    \pgfplotsmanual@configure@for@htlatexmode
}%


%%%%%%%%%%%%%%%%%%%%%%%%%%%%%%%%%%%%%%%%%%%%%%%%%%%%%%%%%%%%%%%%%%%%%%%%%%%%%%



\newenvironment{addplotoperation}[3][]{
  \begin{pgfmanualentry}
      {%
    \let\ltxdoc@marg=\marg
    \let\ltxdoc@oarg=\oarg
    \let\ltxdoc@parg=\parg
    \let\ltxdoc@meta=\meta
    \def\marg##1{{\normalfont\ltxdoc@marg{##1}}}%
    \def\oarg##1{{\normalfont\ltxdoc@oarg{##1}}}%
    \def\parg##1{{\normalfont\ltxdoc@parg{##1}}}%
    \def\meta##1{{%
    \ifmmode
        \mathit{##1}%
    \else
        \normalfont\ltxdoc@meta{##1}%
    \fi
}}%
    \pgfmanualentryheadline{\textcolor{gray}{{\ttfamily\char`\\addplot\ }}%
      \declare{\texttt{#2}} \texttt{#3;}}%
      \unskip
     \nobreak
    \pgfmanualentryheadline{\textcolor{gray}{\texttt{\char`\\addplot}\oarg{options} }%
      \declare{\texttt{#2}} \texttt{#3} \textcolor{gray}{\meta{trailing path commands}}\texttt{;}}%
      \unskip
     \nobreak
    \pgfmanualentryheadline{\textcolor{gray}{{\ttfamily\char`\\addplot3}} $\dotsc$}%
    \def\pgfmanualtest{#1}%
    \ifx\pgfmanualtest\@empty%
      \index{#2@\protect\textcolor{gray}{\protect\texttt{plot}}\protect\texttt{ #2}}%
      \index{Plot operations!plot #2@\protect\texttt{plot #2}}%
    \fi%
    \pgfmanualpdflabel{\textbackslash addplot #2}{}%
    \pgfmanualpdflabel{plot #2}{}%
    \pgfmanualpdflabel{#2}{}%
    }%
    \pgfmanualbody
}
{
  \end{pgfmanualentry}
}

\newenvironment{addplot+}{
  \begin{pgfmanualentry}
      {%
    \let\ltxdoc@marg=\marg
    \let\ltxdoc@oarg=\oarg
    \let\ltxdoc@parg=\parg
    \let\ltxdoc@meta=\meta
    \def\marg##1{{\normalfont\ltxdoc@marg{##1}}}%
    \def\oarg##1{{\normalfont\ltxdoc@oarg{##1}}}%
    \def\parg##1{{\normalfont\ltxdoc@parg{##1}}}%
    \def\meta##1{{%
    \ifmmode
        \mathit{##1}%
    \else
        \normalfont\ltxdoc@meta{##1}%
    \fi
}}%
    \pgfmanualentryheadline{{\ttfamily\declare{\char`\\addplot+}\oarg{options} \textcolor{gray}{\dots};}}%
    \index{addplot+@\protect\texttt{\protect\textbackslash addplot+}}%
    \pgfmanualpdflabel{\textbackslash addplot+}{}%
    }%
    \pgfmanualbody
}
{
  \end{pgfmanualentry}
}
\newenvironment{addplot3generic}{
  \begin{pgfmanualentry}
      {%
    \let\ltxdoc@marg=\marg
    \let\ltxdoc@oarg=\oarg
    \let\ltxdoc@parg=\parg
    \let\ltxdoc@meta=\meta
    \def\marg##1{{\normalfont\ltxdoc@marg{##1}}}%
    \def\oarg##1{{\normalfont\ltxdoc@oarg{##1}}}%
    \def\parg##1{{\normalfont\ltxdoc@parg{##1}}}%
    \def\meta##1{{\normalfont\ltxdoc@meta{##1}}}%
    \pgfmanualentryheadline{{\ttfamily\declare{\char`\\addplot3}\oarg{options} \meta{input data} \meta{trailing path commands};}}%
    \index{addplot3@\protect\texttt{\protect\textbackslash addplot3}}%
    \pgfmanualpdflabel{\textbackslash addplot3}{}%
    }%
    \pgfmanualbody
}
{
  \end{pgfmanualentry}
}
\newenvironment{addplot3operation}[3][]{
  \begin{pgfmanualentry}
      {%
    \let\ltxdoc@marg=\marg
    \let\ltxdoc@oarg=\oarg
    \let\ltxdoc@parg=\parg
    \let\ltxdoc@meta=\meta
    \def\marg##1{{\normalfont\ltxdoc@marg{##1}}}%
    \def\oarg##1{{\normalfont\ltxdoc@oarg{##1}}}%
    \def\parg##1{{\normalfont\ltxdoc@parg{##1}}}%
    \def\meta##1{{\normalfont\ltxdoc@meta{##1}}}%
    \pgfmanualentryheadline{\textcolor{gray}{{\ttfamily\char`\\addplot3\ }}%
      \declare{\texttt{#2}} \texttt{#3;}}%
      \unskip
     \nobreak
    \pgfmanualentryheadline{\textcolor{gray}{\texttt{\char`\\addplot3}\oarg{options} }%
      \declare{\texttt{#2}} \texttt{#3} \textcolor{gray}{\meta{trailing path commands}}\texttt{;}}%
    \def\pgfmanualtest{#1}%
    \ifx\pgfmanualtest\@empty%
      \index{#2@\protect\texttt{#2}}%
      \index{Plot operations!addplot3 #2@\protect\texttt{#2}}%
    \fi%
    \pgfmanualpdflabel{\textbackslash addplot3 #2}{}%
    \pgfmanualpdflabel{plot3 #2}{}%
    }%
    \pgfmanualbody
}
{
  \end{pgfmanualentry}
}

\newenvironment{codekey}[1]{%
  \begin{pgfmanualentry}
     \pgfmanualentryheadline{{\ttfamily\declarekey{#1}\textcolor{gray}{/\pgfmanualpdfref{/handlers/.code}{.code}}=\marg{...}}\hfill}%
    \def\mykey{#1}%
    \def\mypath{}%
    \def\myname{}%
    \firsttimetrue%
    \decompose#1/\nil%
    \pgfmanualbody
}
{
  \end{pgfmanualentry}
}
\newenvironment{codeargskey}[2]{%
  \begin{pgfmanualentry}
      {\toks0={#2}%
      \xdef\argpattern{\the\toks0 }%
    }%
      \pgfmanual@command@to@string\argpattern\argpattern
     \pgfmanualentryheadline{{\ttfamily\declarekey{#1}\textcolor{gray}{/\pgfmanualpdfref{/handlers/.code}{.code args}}=\texttt{\{\argpattern\}}\marg{...}}\hfill}%
    \def\mykey{#1}%
    \def\mypath{}%
    \def\myname{}%
    \firsttimetrue%
    \decompose#1/\nil%
    \pgfmanualbody
}
{
  \end{pgfmanualentry}
}
\def\pgfmanual@command@to@string#1#2{%
    \expandafter\pgfmanual@command@to@string@@\meaning#1\pgfmanual@EOI{#2}%
}%
\xdef\pgfmanual@glob@TMPa{\meaning\pgfutil@empty}%
\expandafter\def\expandafter\pgfmanual@command@to@string@@\pgfmanual@glob@TMPa#1\pgfmanual@EOI#2{%
    \def#2{#1}%
}%

\newenvironment{pgfplotscodekey}[1]{%
    \begin{codekey}{/pgfplots/#1}%
}
{
  \end{codekey}
}
\newenvironment{pgfplotscodetwokey}[1]{%
  \begin{pgfmanualentry}
     \pgfmanualentryheadline{{\ttfamily\declarekey{/pgfplots/#1}\textcolor{gray}{/\pgfmanualpdfref{/handlers/.code 2 args}{.code 2 args}}=\marg{...}}\hfill}%
    \def\mykey{/pgfplots/#1}%
    \def\mypath{}%
    \def\myname{}%
    \firsttimetrue%
    \decompose/pgfplots/#1/\nil%
    \pgfmanualbody
}
{
  \end{pgfmanualentry}
}

\newenvironment{pgfplotsxycodekeylist}[1]{%
    \begingroup
    \let\oldpgfmanualentryheadline=\pgfmanualentryheadline
    \def\pgfmanualentryheadline##1{%
        \pgfmanualentryheadline@##1\pgfplots@EOI
    }%
    \def\pgfmanualentryheadline@##1\hfill##2\pgfplots@EOI{%
        \oldpgfmanualentryheadline{{\ttfamily\declarekey{##1}\textcolor{gray}{/\pgfmanualpdfref{/handlers/.code}{.code}}=\marg{...}}\hfill}%
    }
    \begin{pgfplotsxykeylist}{#1}%
}
{
    \end{pgfplotsxykeylist}
    \endgroup
}

\newenvironment{pgfplotsxyexpressionkeylist}[1]{%
    \begingroup
    \let\oldpgfmanualentryheadline=\pgfmanualentryheadline
    \def\pgfmanualentryheadline##1{%
        \pgfmanualentryheadline@##1\pgfplots@EOI
    }%
    \def\pgfmanualentryheadline@##1\hfill##2\pgfplots@EOI{%
        \oldpgfmanualentryheadline{{\ttfamily\declarekey{##1}\textcolor{gray}{/.expression}=\marg{math expression}}\hfill}%
    }
    \begin{pgfplotsxykeylist}{#1}%
}
{
    \end{pgfplotsxykeylist}
    \endgroup
}
\newenvironment{pgfplotsxyappendexpressionkeylist}[1]{%
    \begingroup
    \let\oldpgfmanualentryheadline=\pgfmanualentryheadline
    \def\pgfmanualentryheadline##1{%
        \pgfmanualentryheadline@##1\pgfplots@EOI
    }%
    \def\pgfmanualentryheadline@##1\hfill##2\pgfplots@EOI{%
        \oldpgfmanualentryheadline{{\ttfamily\declarekey{##1}\textcolor{gray}{/.append expression}=\marg{math expression}}\hfill}%
    }
    \begin{pgfplotsxykeylist}{#1}%
}
{
    \end{pgfplotsxykeylist}
    \endgroup
}

\newenvironment{pgfplotskey}[1]{%
  \begin{key}{/pgfplots/#1}%
}
{
  \end{key}
}

\newenvironment{tikzkey}[1]{%
  \begin{key}{/tikz/#1}%
}
{
  \end{key}
}

\def\choicesep{$\vert$}%
\def\choicearg#1{\texttt{#1}}

\newif\iffirstchoice
\newcommand\mchoice[1]{%
    \begingroup
    \let\margold=\marg
    \def\marg##1{{\normalfont\margold{##1}}}%
    \firstchoicetrue
    \foreach \mchoice@ in {#1} {%
        \iffirstchoice
            \global\firstchoicefalse
        \else
            \choicesep
        \fi
        \choicearg{\mchoice@}%
    }%
    \endgroup
}%




% \begin{xykey}{/path/\x label=value}
% \end{xykey}
%
% has same features with 'default', 'initially' etc as key environment
\newenvironment{xykey}[2][]{%
    \begin{pgfmanualentry}
    \def\extrakeytext{}
    \insertpathifneeded{#2}{#1}%
    \expandafter\pgfutil@in@\expandafter=\expandafter{\mykey}%
    \ifpgfutil@in@%
        \expandafter\xykey@eq\mykey\@nil
    \else
        \expandafter\xykey@noeq\mykey\@nil
    \fi
    \pgfmanualbody
}{%
    \end{pgfmanualentry}
}%

% \begin{xystylekey}{/path/\x label=value}
% \end{xystylekey}
%
% has same features with 'default', 'initially' etc as key environment
\newenvironment{xystylekey}[2][]{%
    \begin{pgfmanualentry}
    \def\extrakeytext{style, }
    \insertpathifneeded{#2}{#1}%
    \expandafter\pgfutil@in@\expandafter=\expandafter{\mykey}%
    \ifpgfutil@in@%
        \expandafter\xykey@eq\mykey\@nil
    \else
        \expandafter\xykey@noeq\mykey\@nil
    \fi
    \pgfmanualbody
}{%
    \end{pgfmanualentry}
}%

% \insertpathifneeded{a key}{/pgfplots} -> assign mykey={/pgfplots/a key}
% \insertpathifneeded{/tikz/a key}{/pgfplots} -> assign mykey={/tikz/a key}
%
% #1: the key
% #2: a default path (or empty)
\def\insertpathifneeded#1#2{%
    \def\insertpathifneeded@@{#2}%
    \ifx\insertpathifneeded@@\empty
        \def\mykey{#1}%
    \else
        \insertpathifneeded@#1\@nil
        \ifpgfutil@in@
            \def\mykey{#1}%
        \else
            \def\mykey{#2/#1}%
        \fi
    \fi
}%
\def\insertpathifneeded@#1#2\@nil{%
    \def\insertpathifneeded@@{#1}%
    \def\insertpathifneeded@@@{/}%
    \ifx\insertpathifneeded@@\insertpathifneeded@@@
        \pgfutil@in@true
    \else
        \pgfutil@in@false
    \fi
}%

% \begin{keylist}[default path]
%     {/path/option 1=value,/path/option 2=value2}
% \end{keylist}
\newenvironment{keylist}[2][]{%
    \begin{pgfmanualentry}
    \def\extrakeytext{}%
    \foreach \xx in {#2} {%
        \expandafter\insertpathifneeded\expandafter{\xx}{#1}%
        \expandafter\extractkey\mykey\@nil%
    }%
    \pgfmanualbody
}{%
  \end{pgfmanualentry}
}%

\newenvironment{pgfplotskeylist}[1]{%
    \begin{keylist}[/pgfplots]{#1}%
}{%
    \end{keylist}%
}

\newenvironment{tikzkeylist}[1]{%
    \begin{keylist}[/tikz]{#1}%
}{%
    \end{keylist}%
}

\newenvironment{anchorlist}[1]{
  \begin{pgfmanualentry}
      \foreach \xx in {#1} {%
        \pgfmanualentryheadline{锚点 {\ttfamily\declare{\xx}}}%
        \index{\xx @\protect\texttt{\xx} anchor}%
        \index{Anchors!\xx @\protect\texttt{\xx}}
        \expandafter\pgfmanualpdflabel\expandafter{\xx}{}
    }%
    \pgfmanualbody
}
{
  \end{pgfmanualentry}
}

\newenvironment{coordinatesystemlist}[1]{
  \begin{pgfmanualentry}
      \foreach \xx in {#1} {%
        \pgfmanualentryheadline{坐标系 {\ttfamily\declare{\xx}}}%
        \index{\xx @\protect\texttt{\xx} coordinate system}%
        \index{Coordinate systems!\xx @\protect\texttt{\xx}}
        \expandafter\pgfmanualpdflabel\expandafter{\xx}{}
    }%
    \pgfmanualbody
}
{
  \end{pgfmanualentry}
}
\renewenvironment{coordinatesystem}[1]{
  \begin{pgfmanualentry}
    \pgfmanualentryheadline{坐标系 {\ttfamily\declare{#1}}}%
    \index{#1@\protect\texttt{#1} coordinate system}%
    \index{Coordinate systems!#1@\protect\texttt{#1}}
    \pgfmanualpdflabel{#1}{}
    \pgfmanualbody
}
{
  \end{pgfmanualentry}
}

% \begin{xykeylist}[default path]
%     {/path/option \x1=value,/path/option \x2=value2,/path/option \x3=value}
% \end{xykeylist}
\newenvironment{xykeylist}[2][]{%
    \begin{pgfmanualentry}
    \def\extrakeytext{}
    \foreach \xx in {#2} {%
        \expandafter\insertpathifneeded\expandafter{\xx}{#1}%
        \expandafter\pgfutil@in@\expandafter=\expandafter{\mykey}%
        \ifpgfutil@in@%
            \expandafter\xykey@eq\mykey\@nil
        \else
            \expandafter\xykey@noeq\mykey\@nil
        \fi
    }%
    \pgfmanualbody
}{%
  \end{pgfmanualentry}
}%

\makeatother % FIXME this is almost surely a bug in pgfmanual-en-macros
% \begin{commandlist}
%     {\command1{arg1},\command2{\arg2}}
% \end{commandlist}
\newenvironment{commandlist}[1]{%
    \begin{pgfmanualentry}
    \foreach \xx in {#1} {%
        \expandafter\extractcommand\xx\@@%
    }%
    \pgfmanualbody
}{%
  \end{pgfmanualentry}
}%

\newenvironment{texif}[1]{%
    \begin{pgfmanualentry}
    \pgfmanualentryheadline{\declare{\texttt{\textbackslash if#1}}\meta{true code}\texttt{\textbackslash else}\meta{else code}\texttt{\textbackslash fi}}%
    \index{if#1}%
    \pgfmanualpdflabel{\\if#1}{}%
    \pgfmanualbody
}{%
  \end{pgfmanualentry}
}%
\makeatletter

\newif\ifxykeyfound

\def\pgfmanual@xykey@install@replacements{%
    \def\ { }%
    \def\space{ }%
}%

\def\xykey@eq#1=#2\@nil{%
    \begingroup
    \pgfmanual@xykey@install@replacements
    \def\x{x}%
    \xdef\mykey{#1}%
    \def\xykey@@{#1}%
    \ifx\xykey@@\mykey
        \xykeyfoundfalse
    \else
        \xykeyfoundtrue
    \fi
    \expandafter\extractkey\mykey=#2\@nil%
    \ifxykeyfound
        \def\x{y}%
        \xdef\mykey{#1}%
        \expandafter\extractkey\mykey=#2\@nil%
        \def\x{z}%
        \xdef\mykey{#1}%
        \expandafter\extractkey\mykey=#2\@nil%
    \fi
    \endgroup
}
\def\xykey@noeq#1\@nil{%
    \begingroup
    \pgfmanual@xykey@install@replacements
    \def\x{x}%
    \xdef\mykey{#1}%
    \def\xykey@@{#1}%
    \ifx\xykey@@\mykey
        \xykeyfoundfalse
    \else
        \xykeyfoundtrue
    \fi
    \expandafter\extractkey\mykey\@nil%
    \ifxykeyfound
        \def\x{y}%
        \xdef\mykey{#1}%
        \expandafter\extractkey\mykey\@nil%
        \def\x{z}%
        \xdef\mykey{#1}%
        \expandafter\extractkey\mykey\@nil%
    \fi
    \endgroup
}

% \begin{pgfplotsxykey}{\x label=value}
% \end{pgfplotsxykey}
%
% It introduces the path /pgfplots/ automatically.
%
% has same features with 'default', 'initially' etc as key environment
\newenvironment{pgfplotsxykey}[1]{%
    \begin{xykey}[/pgfplots]{#1}%
}{%
    \end{xykey}%
}


\newenvironment{pgfplotsxykeylist}[1]{%
    \begin{xykeylist}[/pgfplots]{#1}%
}{%
    \end{xykeylist}%
}


% the first, optional argument is the default key path to insert.
\newenvironment{plottype}[2][/tikz]{%
    \begin{keylist}[#1]{#2}%
    \end{keylist}
  \begin{pgfmanualentry}
    \pgfmanualentryheadline{\textcolor{gray}{{\ttfamily\char`\\addplot+[\declare{#2}]}}}%
    \pgfmanualbody
}
{
  \end{pgfmanualentry}
}

\def\index@prologue{\section*{索引}\addcontentsline{toc}{section}{索引}
}

\newenvironment{pgfplotstablecolumnkey}{%
  \begin{pgfmanualentry}
     \pgfmanualentryheadline{{\ttfamily\textcolor{gray}{/pgfplots/table/}\declare{columns/\meta{column name}}\textcolor{gray}{/.style}=\marg{key-value-list}}\hfill}%
    \pgfplotsmanualkeyindex{/pgfplots/table/columns}%
    \pgfmanualbody
}
{
  \end{pgfmanualentry}
}
\newenvironment{pgfplotstabledisplaycolumnkey}{%
  \begin{pgfmanualentry}
     \pgfmanualentryheadline{{\ttfamily\textcolor{gray}{/pgfplots/table/}\declare{display columns/\meta{index}}\textcolor{gray}{/.style}=\marg{key-value-list}}\hfill}%
    \pgfplotsmanualkeyindex{/pgfplots/table/display columns}%
    \pgfmanualbody
}
{
  \end{pgfmanualentry}
}
\newenvironment{pgfplotstablealiaskey}{%
  \begin{pgfmanualentry}
     \pgfmanualentryheadline{{\ttfamily\textcolor{gray}{/pgfplots/table/}\declare{alias/\meta{col name}}\textcolor{gray}{/.initial}=\marg{real col name}}\hfill}%
    \pgfplotsmanualkeyindex{/pgfplots/table/alias}%
    \pgfmanualbody
}
{
  \end{pgfmanualentry}
}


\def\pgfplotsmanualkeyindex#1{%
    \def\mypath{#1}%
    \def\myname{}%
    \firsttimetrue%
    \decompose#1/\nil%
}
\newenvironment{pgfplotstablecreateonusekey}{%
  \begin{pgfmanualentry}
     \pgfmanualentryheadline{{\ttfamily\textcolor{gray}{/pgfplots/table/}\declare{create on use/\meta{col name}}\textcolor{gray}{/.style}=\marg{create options}}\hfill}%
    \def\mykey{/pgfplots/table/create on use}%
    \pgfmanualbody
    \pgfplotsmanualkeyindex{/pgfplots/table/create on use}%
}
{
  \end{pgfmanualentry}
}

\def\pgfplotsassertcmdkeyexists#1{%
    \pgfkeysifdefined{/pgfplots/#1/.@cmd}\relax{%
        \pgfplots@error{DOCUMENTATION ERROR: command key /pgfplots/#1 does not exist!}%
    }%
}%

{
\catcode`\ =12%
\gdef\makespaceexpandable{\def\ { }}}%

\def\pgfplotsassertXYcmdkeyexists#1{%
    {\makespaceexpandable\def\x{x}\edef\pgfplotsassertXYcmdkeyexists@tmp{#1}%
    \pgfkeysifdefined{/pgfplots/\pgfplotsassertXYcmdkeyexists@tmp/.@cmd}\relax{%
        \pgfplots@error{DOCUMENTATION ERROR: command key /pgfplots/#1 does not exist!}%
    }}%
    {\makespaceexpandable\def\x{y}\edef\pgfplotsassertXYcmdkeyexists@tmp{#1}%
    \pgfkeysifdefined{/pgfplots/\pgfplotsassertXYcmdkeyexists@tmp/.@cmd}\relax{%
        \pgfplots@error{DOCUMENTATION ERROR: command key /pgfplots/#1 does not exist!}%
    }}%
}%

\def\pgfplotsshortstylekey #1=#2\pgfeov{%
    \pgfplotsassertcmdkeyexists{#1}%
    \pgfplotsassertcmdkeyexists{#2}%
    \begin{pgfplotskey}{#1=\marg{key-value-list}}
        是 \texttt{\pgfmanualpdfref{#2}{#2}/\pgfmanualpdfref{/handlers/.append style}{.append style}=}\marg{key-value-list} 的简写。

        它把选项追加到已存在的样式 \texttt{\pgfmanualpdfref{#2}{#2}}。
    \end{pgfplotskey}
}
\def\pgfplotsshortxystylekey #1=#2\pgfeov{%
    \pgfplotsassertXYcmdkeyexists{#1}%
    \pgfplotsassertXYcmdkeyexists{#2}%
    \begin{pgfplotsxykey}{#1=\marg{key-value-list}}
        是 {\def\x{x}\texttt{\pgfmanualpdfref{#2}{#2}/\pgfmanualpdfref{/handlers/.append style}{.append style}=}}\marg{key-value-list}
        的简写（或 $y$ 轴对应的样式
            {\def\x{y}\texttt{\pgfmanualpdfref{#2}{#2}/\pgfmanualpdfref{/handlers/.append style}{.append style}=}}\marg{key-{}value-{}list}，
        以及 $z$ 轴
            {\def\x{z}\texttt{\pgfmanualpdfref{#2}{#2}/\pgfmanualpdfref{/handlers/.append style}{.append style}=}}\marg{key-{}value-{}list}）。

        它把选项追加到已存在的样式 {\def\x{x}\texttt{\pgfmanualpdfref{#2}{#2}}}。
    \end{pgfplotsxykey}
}
\def\pgfplotsshortstylekeys #1,#2=#3\pgfeov{%
    \pgfplotsassertcmdkeyexists{#1}%
    \pgfplotsassertcmdkeyexists{#2}%
    \pgfplotsassertcmdkeyexists{#3}%
    \begin{pgfplotskeylist}{%
        #1=\marg{key-value-list},
        #2=\marg{key-value-list}}
        是 \texttt{\pgfmanualpdfref{#3}{#3}/\pgfmanualpdfref{/handlers/.append style}{.append style}=}\marg{key-value-list} 的不同简写。
    \end{pgfplotskeylist}
}
\def\pgfplotsshortxystylekeys #1,#2=#3\pgfeov{%
    \pgfplotsassertXYcmdkeyexists{#1}%
    \pgfplotsassertXYcmdkeyexists{#2}%
    \pgfplotsassertXYcmdkeyexists{#3}%
    \begin{pgfplotsxykeylist}{%
        #1=\marg{key-value-list},
        #2=\marg{key-value-list}}
        是 {\def\x{x}\texttt{\pgfmanualpdfref{#3}{#3}/\pgfmanualpdfref{/handlers/.append style}{.append style}=}}\marg{key-value-list}
        的不同简写（或 $y$ 轴对应的样式
            {\def\x{y}\texttt{\pgfmanualpdfref{#3}{#3}/\pgfmanualpdfref{/handlers/.append style}{.append style}=}\marg{key-{}value-{}list}}，以及 $z$ 轴
            {\def\x{z}\texttt{\pgfmanualpdfref{#3}{#3}/\pgfmanualpdfref{/handlers/.append style}{.append style}=}\marg{key-{}value-{}list}}%
            ）。
    \end{pgfplotsxykeylist}
}


%
% For using the correct form of including libraries in the manual.
%
\newenvironment{pgfplotslibrary}[1]{%
  \begin{pgfmanualentry}
    \pgfmanualentryheadline{{\ttfamily\char`\\usepgfplotslibrary\char`\{\declare{#1}\char`\}\space\space \char`\%\space\space  \LaTeX\space and plain \TeX}}%
    \index{#1@\protect\texttt{#1} library}%
    \index{Libraries!#1@\protect\texttt{#1}}%
    \pgfmanualentryheadline{{\ttfamily\char`\\usepgfplotslibrary[\declare{#1}]\space \char`\%\space\space Con\TeX t}}%
    \pgfmanualentryheadline{{\ttfamily\char`\\usetikzlibrary\char`\{\declare{pgfplots.#1}\char`\}\space\space \char`\%\space\space \LaTeX\space and plain \TeX}}%
    \pgfmanualentryheadline{{\ttfamily\char`\\usetikzlibrary[\declare{pgfplots.#1}]\space \char`\%\space\space Con\TeX t}}%
    \pgfmanualpdflabel{#1}{}%
    \pgfmanualbody
}
{
  \end{pgfmanualentry}
}


%
% Creates and shows a colormap with specification '#1'.
\def\pgfplotsshowcolormapexample#1{%
    \pgfplotscreatecolormap{tempcolormap}{#1}%
    \pgfplotsshowcolormap{tempcolormap}%
}

% Shows the colormap named '#1'.
\def\pgfplotsshowcolormap#1{%
    \pgfplotscolormapifdefined{#1}{\relax}{%
        \pgfplotsset{colormap/#1}%
    }%
    \pgfplotscolormaptoshadingspec{#1}{8cm}\result
    \def\tempb{\pgfdeclarehorizontalshading{tempshading}{1cm}}%
    \expandafter\tempb\expandafter{\result}%
    \pgfuseshading{tempshading}%
}

\makeatother

\def\decompose/#1/#2\nil{%
  \def\test{#2}%
  \ifx\test\empty%
    % aha.
    \index{#1@\protect\texttt{#1} key}%
    \ifx\mypath\empty
    \else
        \index{\mypath#1@\protect\texttt{#1}}%
    \fi
    \def\myname{#1}%
    %\pgfmanualpdflabel{#1}{}% No, its better to use fully qualified keys and search if necessary!
  \else%
    \iffirsttime
        \begingroup
            % also make a pdf link anchor with full key path.
            \def\hyperlabelwithoutslash##1/\nil{%
                \pgfmanualpdflabel{##1}{}%
            }%
            \hyperlabelwithoutslash/#1/#2\nil
        \endgroup
%    CF : disabled for /pgfplots/ prefix.
%        \def\mypath{#1@\protect\texttt{/#1/}!}%
%        \firsttimefalse
        \def\pgfplotslocTMPa{pgfplots}%
        \edef\pgfplotslocTMPb{#1}%
        \ifx\pgfplotslocTMPb\pgfplotslocTMPa
            \def\mypath{}%
        \else
            \def\mypath{#1@\protect\texttt{/#1/}!}%
        \fi
        \firsttimefalse
    \else
      \expandafter\def\expandafter\mypath\expandafter{\mypath#1@\protect\texttt{#1/}!}%
    \fi
    \def\firsttime{}
    \decompose/#2\nil%
  \fi%
}
\def\extracthandler#1#2\@nil{%
  \pgfmanualentryheadline{键处理器 \meta{key}{\ttfamily/\declare{#1}}#2}%
  \index{\gobble#1@\protect\texttt{#1} handler}%
  \index{Key handlers!#1@\protect\texttt{#1}}
  \pgfmanualpdflabel{/handlers/#1}{}%
}
\def\extractcommand#1#2\@@{%
  \pgfmanualentryheadline{\declare{\texttt{\string#1}}#2}%
  \removeats{#1}%
  \index{\strippedat @\protect\myprintocmmand{\strippedat}}%
  \pgfmanualpdflabel{\textbackslash\strippedat}{}%
}
\def\extractenvironement#1#2\@@{%
  \pgfmanualentryheadline{{\ttfamily\char`\\begin\char`\{\declare{#1}\char`\}}#2}%
  \pgfmanualentryheadline{{\ttfamily\ \ }\meta{environment contents}}%
  \pgfmanualentryheadline{{\ttfamily\char`\\end\char`\{\declare{#1}\char`\}}}%
  \index{#1@\protect\texttt{#1} environment}%
  \index{Environments!#1@\protect\texttt{#1}}%
  \pgfmanualpdflabel{#1}{}%
}
\renewenvironment{predefinednode}[1]{
  \begin{pgfmanualentry}
    \pgfmanualentryheadline{预定义节点 {\ttfamily\declare{#1}}}%
    \index{#1@\protect\texttt{#1} node}%
    \index{Predefined node!#1@\protect\texttt{#1}}
    \pgfmanualpdflabel{#1}{}%
    \pgfmanualbody
}
{
  \end{pgfmanualentry}
}


\makeatletter
\@ifpackageloaded{tex4ht}{
}{%
    \IfFileExists{ocg.sty}{%
        \usepackage{ocg}%
    }{%
        \usepackage{pgfplots_ocg_copy}%
    }
}%
\makeatother

\usepackage{nicefrac}

\graphicspath{{figures/}}

\edef\preambleconfig{width=7cm,compat=\pgfkeysvalueof{/pgfplots/compat/mostrecent}}

% generate the following file such that it does not need to be updated
% manually for each release:
\makeatletter
\immediate\openout\w@pgf@writea pgfplots.preliminaries.compatcurrent.tmp
\immediate\write\w@pgf@writea{%
\string\begin{codeexample}[code only]^^J%
\string\usepackage{pgfplots}^^J%
\string\pgfplotsset{compat=\pgfkeysvalueof{/pgfplots/compat/mostrecent}}^^J%
\string\end{codeexample}%
}%
\immediate\closeout\w@pgf@writea


\expandafter\pgfplotsset\expandafter{\preambleconfig}

\pgfplotsset{lua debug=off and silent}

\def\pgfplotsmanualcurlibrary{}

% And now, invoke
%     /codeexample/typeset listing/.add={% Preamble:\pgfplotsset{\preambleconfig}}{}}
% since listings are VERBATIM, I need to do some low-level things
% here to get the correct \catcodes:
\pgfkeys{/codeexample/typeset listing/.add code={%
        \ifcode@execute
            \begingroup
            \def\b@show@header{0}%
            \pgfutil@in@{axis}{#1}%
            \ifpgfutil@in@
                \def\b@show@header{1}%
            \fi
            %
            \pgfutil@in@{addplot}{#1}%
            \ifpgfutil@in@
                \def\b@show@header{1}%
            \fi
            %
            \pgfutil@in@{smithchart}{#1}%
            \ifpgfutil@in@
                \def\b@show@header{1}%
            \fi
            \if1\b@show@header
                {\tiny
                    \% Preamble: \pgfmanualpdfref{\textbackslash pgfplotsset}{\pgfmanual@pretty@backslash pgfplotsset}%
                        \pgfmanual@pretty@lbrace \expandafter\pgfmanualprettyprintpgfkeys\expandafter{\preambleconfig}\pgfmanual@pretty@rbrace
                        \ifx\pgfplotsmanualcurlibrary\pgfutil@empty
                        \else
                            \pgfmanual@pretty@backslash usepgfplotslibrary\pgfmanual@pretty@lbrace \pgfplotsmanualcurlibrary\pgfmanual@pretty@rbrace
                        \fi
                }%
            \fi
            \endgroup
        \fi
    }{},%
    %/codeexample/typeset listing/.show code,
}%
\makeatother

\pgfplotsset{
    %every axis/.append style={width=7cm},
    filter discard warning=false,
}

\pgfqkeys{/codeexample}{%
    every codeexample/.append style={
        width=8cm,
        /pgfplots/legend style={fill=graphicbackground},
        /pgfplots/contour/every contour label/.append style={
            every node/.append style={fill=graphicbackground}
        },
    },
    tabsize=4,
}

\usetikzlibrary{backgrounds,patterns}
% Global styles:
\tikzset{
  shape example/.style={
    color=black!30,
    draw,
    fill=yellow!30,
    line width=.5cm,
    inner xsep=2.5cm,
    inner ysep=0.5cm}
}

\newcommand{\FIXME}[1]{\textcolor{red}{(FIXME: #1)}}

% fuer environment 'sidewaysfigure' bspw
% \usepackage{rotating}

\newcommand\Tikz{Ti\textit kZ}
\newcommand\PGF{\textsc{pgf}}
\newcommand\PGFPlots{\pgfplotsmakefilelinkifuseful{pgfplots}{\textsc{pgfplots}}}
\newcommand\PGFPlotstable{\pgfplotsmakefilelinkifuseful{pgfplotstable}{\textsc{PgfplotsTable}}}

\makeindex

% Fix overfull hboxes automatically:
\tolerance=2000
\emergencystretch=10em

\tikzset{prefix=gnuplot/pgfplots_} % prefix for 'plot function'

\author{%
    Dr.\ Christian Feuersänger\\
    {\footnotesize\texttt{cfeuersaenger@users.sourceforge.net}}}%


\usepackage{todonotes}
\newcommand\todosp[2][]{%        % Stefan Pinnow
    \todo[
        disable,
        color=blue!40,
        #1]{#2}
}

% ---------------------------------------------------------------
% 修复：doc 宏包 v2 的 \ltxdoc@meta 在 TeX Live 2026 中未定义，
% 导致 \marg/\oarg/\parg（内部调用 \meta）写 aux/idx 时报
% "Undefined control sequence"。
% 对策（双保险）：
%  1) 兜底提供 \ltxdoc@meta（doc 的 \meta 依赖它）；
%  2) 在 \AtBeginDocument 末尾覆盖 \meta 为鲁棒命令
%     （doc v2 会在 begin document 时重定义 \meta，必须最后覆盖）。
\makeatletter
\providecommand\ltxdoc@meta[1]{%
    {\normalfont\itshape\textlangle #1\textrangle}%
}
\AtBeginDocument{%
    \DeclareRobustCommand\meta[1]{%
        {\normalfont\itshape\textlangle #1\textrangle}%
    }%
}
\makeatother


\usepackage{array}
\usepackage{colortbl}
\usepackage{booktabs}
\usepackage{eurosym}

\long\def\codeexamplenl{\noexpand\par}%
\pgfqkeys{/codeexample}{%
	every codeexample/.style={
		width=4cm,
		/pgfplots/every axis/.append style={legend style={fill=graphicbackground}}
	},
	narrow/.style={width=7cm},
	tabsize=4,
	%pre={\begin{minipage}{\linewidth}\begingroup},
	%post={\endgroup\end{minipage}},
	vbox,
	newline=\codeexamplenl,
}

\def\API{\textsc{Api}}%

\title{\TeX\ 编程札记}%\\and Library Functions from \PGF\ and \PGFPlots}

\begin{document}
\maketitle
\begin{abstract}%
	本文档收录的札记面向对 \TeX\ 编程感兴趣的读者。它以大量示例帮助初学者入门，也对希望深入了解 \TeX\ 编程细节、经验丰富的 \TeX nician 有参考价值。但它既不是完整的参考手册，也不是 \TeX\ 的完整使用手册。
\end{abstract}
\tableofcontents
\section{引言}
本文档旨在为 \TeX\ 编程（不一定是 \TeX\ 排版）提供一条直接的入门途径。目标读者是对编写宏包或库感兴趣的人。

在撰写本文档之时，它还远未完善。尽管如此，它仍可作为感兴趣读者的良好起点。更多细节请参阅下文给出的文献。

\section{\TeX\ 编程}
\subsection{寄存器中的变量}
\label{sec:variables}
\TeX\ 提供若干种不同的变量及其相关寄存器，可以自由操作。

\begin{command}{\count\meta{num}}
	共有 256 个整数寄存器，提供 32 位整数运算。这些寄存器可以这样使用，例如 |\count0=42 | 或 |\count7=\macro |，其中 |\macro| 展开为一个数字。

	寄存器的值可以用 |\the|\meta{register} 排版输出。
\begin{codeexample}[]
\count0=42
The value is now `\the\count0'.
\def\macro{-123456}
\count0=\macro
The value is now `\the\count0'.
\end{codeexample}

	`|=|' 符号是可选的，可以省略。所有寄存器有一个共同点：形如 |\count0=|\meta{$\cdots$} 的赋值会把其后的一切不断展开，直到展开不再需要更多数字为止——即使后面跟着不止一个宏。
\begin{codeexample}[]
\def\firstmacro{123}
\def\secondmacro{456}
\def\thirdmacro{789}
\count0=\firstmacro\secondmacro\thirdmacro
The value is now `\the\count0'.
\end{codeexample}
 精确的规则可参见~\cite{texbook}，但应当记住此处需要小心。我的代码曾不止一次未能产生预期结果，原因是 \TeX\ 持续展开宏，导致寄存器得到了意外的值。正确的做法如下：
\begin{codeexample}[]
1. \count0=42 % a white space after the number aborts the reading process. It is discarded.
The value is now `\the\count0'.
2. The following code will absorb the `3' of '3.':
\def\macro{1234}
\count0=\macro % a white space after a macro will be absorbed by TeX, so this is wrong.
3. The value is now `\the\count0'.
4. Use \textbackslash relax after an assignment to end scanning:
\count0=\macro\relax
5. The value is now `\the\count0'.
\end{codeexample}
	命令 |\relax| 告诉 \TeX\ ``放松''：它停止扫描记号，但 |\relax| 自身不展开为任何东西。
	\index{relax@\texttt{\textbackslash relax}}%
\end{command}

\begin{command}{\dimen\meta{num}}
	还有 255 个用于定点数的寄存器，其用法与 |\count| 寄存器大致相同——但 |\dimen| 寄存器的赋值需要带单位，如 `|cm|' 或 `|pt|'。

	用 `|\the|' 进行字符串访问的方式与 |\count| 寄存器完全相同。
\begin{codeexample}[]
\dimen0=1pt
The value is now \the\dimen0.
\dimen0=0.0001pt
The value is now \the\dimen0.
\def\macro{1234.5678}
\dimen0=\macro pt
The value is now \the\dimen0.
\end{codeexample}
	赋值之后宏展开的相同规则在此同样适用。

	|\dimen| 寄存器在内部用 32 位缩放整数进行运算，即所谓的 `scaled point'，单位为 `|sp|'。有关系 |1pt=65536sp|=$2^{16}$|sp|。32 位中有一位用作符号位。以 |pt| 计的整个数值范围为 $[-(2^{30}-1)/2^{16}, (2^{30}-1)/2^{16} ] = [-16383.9998,+16383.9998]$\footnote{请注意，这并未覆盖 32 位整数的完整范围，原因我不清楚。}。
\end{command}

\begin{command}{\toks\meta{number}}
\label{cmd:toks}
	还有 255 个记号寄存器，可以把它们看作特殊的字符串变量。当然，每个宏赋值 |\def\macro|\marg{content} 也算是某种字符串变量，但记号寄存器是特殊的：当与 |\the\toks|\meta{number} 一起使用时，其内容不会被展开。这可用于细粒度的展开控制，参见下文节~\ref{sec:expansion:control}。

\begin{codeexample}[]
\toks0={abc}%
\toks1={DEF}%
The value is now \the\toks0 \the\toks1.
\end{codeexample}
	注意 |\the\toks0| 之后的空白：其作用是在 \TeX\ 扫描 |0| 时停止数字解析。该空白会被丢弃。

	记号寄存器还可以包含特殊记号 |#|，而它在宏内部通常具有特殊含义：
\begin{codeexample}[]
\toks0={#1}%
\message{Meaning is \the\toks0}%
\end{codeexample}
	这会在你的日志文件中输出 ``\texttt{Meaning is \#\#1}''。

	记号寄存器主要在细粒度展开控制方面有用，将在节~\ref{sec:expansion:control} 中更深入地讨论。
\end{command}

\subsubsection{分配寄存器}
寄存器的数量非常有限。因此，必须仔细考虑如何分配它们。寄存器的典型用途是临时变量（例如某个中间结果）以及在生成文档或其某一部分时需要不断累积的长期资源。

随意覆盖一个寄存器显然是个坏主意。

\TeX\ 只有一种 ``分配'' 寄存器的方式：
\begin{commandlist}{\newdimen\meta{\textbackslash macroname},\newcount\meta{\textbackslash macroname},\newtoks\meta{\textbackslash macroname}}
	这些宏分配一个新寄存器，之后便可通过 \meta{\textbackslash macroname} 访问它。
\begin{codeexample}[]
\newdimen\variable

\variable=42pt

The value is now \the\variable.
\end{codeexample}

	得到的 \meta{\textbackslash macroname} 现在可以像直接使用寄存器一样使用。事实上，它往往更简单，因为你无需关心寄存器的编号。

	这种分配依赖于某个全局整数变量，每分配一次它就递增一次。这保证了存放在此类已分配变量中的数据不会意外覆盖其他变量的内容。

	请注意，寄存器无法被释放。
\end{commandlist}

虽然分配长期变量完全合理，但应当避免仅仅因为需要一个 ``新的'' 临时变量就去分配一个新变量。

合理的做法是分配若干具名变量，如 |\tempa|、|\tempb| 之类，并在每次临时求值时复用它们。当然，需要小心避免意外覆盖。

也可以如上所述使用记号寄存器。不过，其使用应当借助分组加以保护：
\begin{codeexample}[]
\toks3={Value outside of group}

\begingroup
\toks3={Value inside of group}
toks3 inside of group: \the\toks3
\endgroup

toks3 outside of group: \the\toks3
\end{codeexample}
	分组构成了 \TeX\ 的 ``作用域'' 概念，本文档的其他地方会加以解释。

\subsubsection{使用超过 256 个寄存器}
\TeX\ 本身被限制为 256 个寄存器。不过，你可以通过使用以下命令手动启用 ``extended \TeX\ 模式''：
\begin{codeexample}[code only]
\usepackage{etex}
\end{codeexample}
请把它放在导言区的开头。这其实是个非常好的做法：它允许访问 65536 个寄存器。如今涉及大量宏包的文档实际上都需要 |etex|。

请注意，即使有了 |etex|，也不能因此为了存放临时变量而毫无节制地随意分配寄存器。

如果你想要几乎无限的临时变量，应当把临时量存放在宏中。这当然涉及从数字到字符串的转换，但它是避免寄存器数量受限的唯一安全方法。

\subsection{\TeX\ 中的算术}
\begin{command}{\advance\meta{register}\texttt{ by}\meta{quantity}}
\begin{codeexample}[]
\count0=42
\advance\count0 by 10
The value is now \the\count0.
\end{codeexample}

\begin{codeexample}[]
\dimen0=1pt
\advance\dimen0 by 10pt
The value is now \the\dimen0.
\end{codeexample}
\end{command}

\begin{command}{\multiply\meta{register}\texttt{ by}\meta{integer}}
\begin{codeexample}[]
\count0=42
\multiply\count0 by -10
The value is now \the\count0.
\end{codeexample}

\begin{codeexample}[]
\dimen0=0.5pt
\multiply\dimen0 by 20
The value is now \the\dimen0.
\end{codeexample}
\end{command}

\begin{command}{\divide\meta{register}\texttt{ by}\meta{integer}}
	这允许对 \meta{integer} 做带截断的整数除法。
\begin{codeexample}[]
\count0=5
\divide\count0 by 2
The value is now \the\count0.
\end{codeexample}

	|\dimen| 寄存器的缩放：
\begin{codeexample}[]
\dimen0=10pt
\divide\dimen0 by 20
The value is now \the\dimen0.
\end{codeexample}

	不可能除以某个非整数。
\end{command}

\begin{command}{\dimen\meta{number}\texttt{=}\meta{fixed point number without unit}\textbackslash dimen\meta{number}}
	这允许在 |\dimen| 寄存器中进行定点数乘法。
\begin{codeexample}[]
\dimen1=50pt
\dimen0=0.6\dimen1
The value is now \the\dimen0.
\end{codeexample}
\end{command}

实际上这就是 \TeX\ 所允许的全部。要做一些 ``真正的'' 算术，需要借助强大的宏包，例如带有 |\pgfmathparse|\marg{expression} 的 \PGF。

请注意，这些寄存器有限的数值范围同样适用于任何数值运算的结果。


\subsection{展开控制}
\label{sec:expansion:control}
展开是 \TeX\ 时刻在做的事情。因此，展开控制是理解如何用 \TeX\ 编程的关键概念。

首先要知道的是：\TeX\ 把输入当作一长串 ``记号'' 序列来处理。记号是 \TeX\ 能够理解的最小单位。每个字符在 \TeX\ 第一次见到它时就成为一个记号。每个宏在 \TeX\ 第一次见到它时也成为一个（单个！）记号。

其次要知道的是字符在 \TeX\ 见到它们\emph{之前}是什么样子。尽管这一知识在日常工作中很少用到，但它仍然重要。输入文档中的字符起初只不过是字符而已。即使是那些已知具有特殊含义的字符，如 `|%|'、`|\|' 或花括号 `|{}|'，也都\emph{不}特殊——直到它们被转换为记号为止。转换发生在 \TeX\ 对字符流进行线性处理、第一次遇到它们的时候。记号一旦成为记号，就永远是同一个记号。如果你设法让 \TeX\ 认为 `|\|' 是一个普通字符，而 \TeX\ 只看到了一个反斜杠，那么这个反斜杠就会是一个普通字符记号——即使其后所有反斜杠的含义再次变为特殊。

现在，我们有一长串记号 \meta{token1}\meta{token2}\meta{token3}\meta{token4}\meta{token5}$\cdots$。\TeX\ 按线性顺序逐个处理这些记号。如果 \meta{token1} 是像 `|a|' 这样的字符记号，它就被排版出来。这不是我现在想在这里讨论的；我的重点是如何用 \TeX\ 编程\footnote{当然，排版本身是一门艺术，有大量相关读物。只是不在本札记中讨论。}。因此，本札记中有趣的是 \meta{token1} 是宏的情形。

\subsubsection{宏}
我们上面已经见过宏的一些应用。实际上，愿意读 \TeX\ 编程札记的大多数用户都已见过宏，或许自己也写过一些——例如使用 |\newcommand|（|\newcommand| 是 |\def| 的一个 ``更高级'' 版本，仅在 \LaTeX\ 中使用）。

宏有一个名字，在 \TeX\ 中被当作一个基本记号处理（即使名字非常长）。宏有替换文本。\TeX\ 一遇到宏，就用替换文本替换它出现的位置。此外，宏还可以把其后的一个或多个记号作为参数取走。
\begin{codeexample}[]
\def\macro{This here is actually the replacement text.}
Executing it: `\macro'.
\end{codeexample}
\begin{codeexample}[]
\def\macro#1{replacement with first argument=#1}
Invoking it: \macro{hello!}.
\end{codeexample}
这其实并不令人意外。可能令人意外的是，所接受的参数几乎可以是任何东西。
\begin{codeexample}[]
\def\macro#1-#2.{replacement with arguments: `#1' and `#2'.}
Invoking it: \macro a-sign.
\end{codeexample}
\noindent 最后一个例子中，|\macro| 会遍历紧随 |\macro| 出现位置之后的记号列表。这个记号列表是 ``|a-sign.|''。宏展开是贪婪的，也就是说会使用第一个匹配的模式。我们的 |\macro| 期待先是某种内容，然后是一个减号 `|-|'，再然后是另一个（可能较长的）参数，最后是一个句点 `|.|'。|\macro| 与减号之间的参数作为 |#1| 可用，减号与句点之间的记号作为 |#2| 可用。

\begin{codeexample}[]
\def\macro(#1,#2,#3){I found arguments `#1', `#2' and `#3'.}
\macro(42,43,44)
\end{codeexample}

如我们所见，宏可以通过展开来操纵输入记号：它们取走某些输入参数（可能没有），并把其他记号插入输入记号列表。这些记号将成为接下来要处理的对象。我们很快会学到更多相关内容。

有一个命令有助于理解 \TeX\ 在这里做了什么：

\begin{command}{\meaning\meta{macro}}
	这个命令展开为 \TeX\ 所看到的 \meta{macro} 的内容。
\begin{codeexample}[]
\def\macro{Replacement \textmacro text  \count0=42 \the\count0.}
\message{Debug message: '\meaning\macro'}
\end{codeexample}
结果是，日志文件和终端输出将包含

|Debug message: 'macro:->Replacement \textmacro text \count 0=42 \the \count 0.'|
\end{command}

最后一个例子已经展示了关于 |\def| 的一点：替换文本中仍然可以包含其他宏。

\begin{command}{\def\meta{\textbackslash macroname}\meta{argument pattern}\marg{replacement text}}
	 一个名为 \meta{macroname} 的新宏将被定义（或重新定义）。\marg{replacement text} 是宏体，每当执行该宏时，它就展开为 \marg{replacement text}。\marg{replacement text} 是一个记号列表，可以包含其他宏。在定义之时，\TeX\ \emph{不会}处理（展开）\marg{replacement text}。

	 只有在宏被执行时，\marg{replacement text} 才会被展开。这同样适用于 \marg{replacement text} 内部的任何宏。
\begin{codeexample}[]
\def\macroone{This is macro one}
\def\macrotwo{Macro two contains \macroone.}
Now, I execute it: \macrotwo.
\def\macroone{Redefined macroone}
Now, I execute the second macro again: \macrotwo.
\end{codeexample}

	宏几乎可以在 \TeX\ 文档的任何地方定义。它们也几乎可以在任何地方被调用。

	\meta{argument pattern} 是一个记号列表，可以包含简单字符串、宏参数 `|#|\meta{number}' 或其他宏记号。第一个参数的 \meta{number} 总是 1，第二个必须是 2，依此类推，最多到 9。合法的参数模式有 `|#1#2#3|'、`|(#1,#2,#3)|' 或 `|---\relax|'。当 \TeX\ 执行一个宏时，它会在输入记号列表中搜索 \meta{argument pattern}，直到找到第一个匹配为止。如果找不到匹配，它就会中止并给出一条（或多或少有点用的）错误信息。
\begin{codeexample}[]
\def\macroone abc{\macrotwo}
\def\macrotwo def{\macrothree}
\def\macrothree#1{Got `#1'}
\macroone abcdefg
\end{codeexample}
	最后一个例子包含三个宏定义。然后，\TeX\ 遇到 |\macroone|。此时输入记号列表为

	`|\macroone abcdefg|'.

	|\macroone| 后面的空白会被 \TeX\ 忽略，它们用于界定 \meta{\textbackslash macroname}。现在，\TeX\ 尝试为 \meta{argument pattern} 寻找匹配。它期待 `|abc|'——于是找到了 `|abc|'。这三个记号从输入记号列表中被\emph{移除}，然后 \TeX\ 插入 |\macroone| 的替换文本，即 |\macrotwo|。此时，输入记号列表为

	`|\macrotwo defg|'.

	现在，同样的过程对 |\macrotwo| 继续进行：\TeX\ 搜索所期待的 \marg{argument pattern}，即 `|def|'，把这些记号从输入记号列表中抹去，并转而插入 |\macrotwo| 的替换文本。于是得到

	`|\macrothree g|'.

	最后，|\macrothree| 期待一个参数记号（或用括号括起来的记号列表）。下一个记号是 `|g|'，它从输入记号列表中被取走，然后插入替换文本——`|#1|' 被替换为 `|g|'。于是，记号列表为

	`|Got `g'|'.

	这段文本最终被排版出来（因为它不再进一步展开）。
\end{command}

我们现在看到的是如何用 \TeX\ 宏来修改记号列表。需要明确指出的是：宏展开绝不局限于 \marg{replacement text} 内部所提供的那些记号——如果 \marg{replacement text} 中的最后一个参数是一个需要参数的宏，那么这些参数将从后续记号中取走。通过使用嵌套宏，甚至可以像循环那样处理记号列表的一整段（但我们还不知道如何有条件地影响宏展开，这留待后面讲）。

我们来尝试解决下面这个任务。假设你有一个名为 |\point| 的宏，其 \meta{argument pattern} 为 `|(#1,#2)|'，即

|\def\point(#1,#2){we do something with #1 and #2}|.

\noindent
再假设你想用存放在另一个宏中的内容来调用 |\point|。毕竟，宏也算是一种字符串变量——累积或生成字符串变量、再把它们用作其他宏的输入，是很合理的。假设我们有 |\temp|，且 |\temp| 的内容是 `|(42,1234)|'。调用 |\point| 的第一个想法是用 |\point\temp|。但是：|\point| 搜索的是以 `|(|' 开头的参数模式，而不是以 |\temp| 开头的！这样调用会失败。

\begin{command}{\expandafter\meta{token}\meta{next token}}
	|\expandafter| 命令是一种——乍一看令人困惑的——改变输入记号列表的方法。但是：它解决了我们用 |\point\temp| 遇到的问题！
\begin{codeexample}[]
\def\point(#1,#2){we do something with #1 and #2}
\def\temp{(42,1234)}
\expandafter\point\temp
\end{codeexample}
	为什么这样就行了！？命令 |\expandafter| 会在输入记号列表中扫描 |\expandafter| 之后的那个记号。在我们的例子中它是 |\point|。然后，它再扫描下一个记号，在我们的例子中是 |\temp|（记住：宏被视为基本记号，就像字符 `|a|' 之类一样）。这两个被扫描到的参数从输入记号列表中被移除。接着，|\expandafter| 把 \meta{next token}\emph{展开}一次。在我们的例子中，\meta{next token} 是 |\temp|。|\temp| 的第一层展开是 `|(42,1234)|'。然后，|\expandafter| 把（未展开的）\meta{token} 连同（已展开的）\meta{next token} 的内容一起重新插回输入记号列表。分步骤来看：

	\begin{enumerate}
		\item |\expandafter\point\temp|
		\item 展开 |\expandafter|：接下来的两个记号是 `|\point\temp|'。
		\item 把 |\point| 用作 \meta{token}，把 |\temp| 用作 \meta{next token}。
		\item 把 |\temp| 展开一次，得到记号 `|(42,1234)|'。
		\item 把 \meta{token} 和 \meta{next token} 的展开结果重新插回输入记号列表。于是列表为

			`|\point(42,1234)|'.
		\item 把 |\point| 作为下一个记号展开。
	\end{enumerate}

	再举一个例子：假设我们想调用 |\theimportantmacro|\marg{argument}。但是，\marg{argument} 却存放在另一个宏中！此外，|\theimportantmacro| 被定义为只接受一个参数，而我们想要的参数可能不止一个记号（这意味着需要用花括号把它括起来）。下面的列表可以解决这个问题。
\begin{codeexample}[]
\def\theimportantmacro#1{I got the pre-assembled argument `#1' here.}
\def\temp{xyz}
\expandafter\theimportantmacro\expandafter{\temp}
\end{codeexample}
	这里发生了什么？我们再次逐步套用规则：
	\begin{enumerate}
		\item 在最初的定义之后，记号列表为 |\expandafter\theimportantmacro\expandafter{\temp}|。
		\item \TeX\ 展开 |\expandafter|，把 |\theimportantmacro| 用作 \meta{token}，把第二个 |\expandafter| 用作 \meta{next token}。
		\item 根据规则，\TeX\ 把 \meta{next token} 展开一次。但是：\meta{next token} 又是一个宏，即 |\expandafter|！这会有区别吗？没有：
			\begin{enumerate}
				\item 第二个 |\expandafter| 之后的记号列表是 `|{\temp}|'（3 个记号）。
				\item 因此 \meta{token} 是 `|{|'，\meta{next token} 是 `|\temp|'。
				\item \meta{next token} 的展开结果是 `|xyz|'。
				\item 第二个 |\expandafter| 重新插入它的 \meta{token} 和已展开的 \meta{next token}，即

					`|{xyz|'.

							注意，闭花括号 `|}|' 完全没有被触动，\TeX\ 到目前为止甚至还没看到它。
			\end{enumerate}
			我们从递归中返回。记住：\meta{token} 是 |\theimportantmacro|，而 \meta{next token} 的顶层展开结果——如我们上面所见——是 `|{xyz|'。
		\item \TeX\ 把 \meta{token} 和 \meta{next token} 的展开结果重新插回输入记号列表，于是得到

			`|\theimportantmacro{xyz}|'.

			闭花括号 `|}|' 没有被触动，它就那样留在输入记号列表中。
		\item \TeX\ 展开 |\theimportantmacro|。
	\end{enumerate}

	\meta{next token} 恰好被展开一次。我们已经看到，如果 \meta{next token} 是一个自身会做替换的宏，那么这些替换会递归地进行。但 ``一次'' 究竟是什么意思？我们需要用 |\meaning| 来检查（或用 |\tracingmacros| 工具），因为我们需要看清 \TeX\ 到底做了什么。
\begin{codeexample}[]
\def\macroone{This is macro one \macrotwo}
\def\macrotwo{--2--}
\def\macrothree#1{\def\macrofour{4[#1]}}
\expandafter\macrothree\expandafter{\macroone}%
So far, nothing has been typeset. But now: \macrofour.
\message{We have macrofour = \meaning\macrofour}%
\end{codeexample}
	日志文件（和终端）现在将包含

	`|We have macrofour = macro:->4[This is macro one \macrotwo ]|'.

	发生了什么？我们可以像上一个例子那样来推演。在两次 |\expandafter| 展开之后，\TeX\ 看到的输入记号列表为

	`|\macrothree{This is macro one \macrotwo}|'

	它在执行之后把 |\macrofour| 定义为 `|This is macro one \macrotwo|'。|\macroone| 的顶层展开并没有展开其中嵌套的对 |\macrotwo| 的调用。


	因此，|\expandafter| 是一个可以被展开的普通宏——甚至可以用另一个 |\expandafter| 来展开一个 |\expandafter|。
\end{command}

到目前为止我们看到的是
\begin{enumerate}
	\item |\def| 命令，它把\emph{未展开的}参数存放在一个宏变量中；以及
	\item |\expandafter|，它允许控制宏的顶层展开（它只展开一次）。
\end{enumerate}
\TeX\ 还提供了另外两个用于展开控制的特性：|\edef| 宏和记号寄存器。

\begin{command}{\edef\meta{\textbackslash macroname}\meta{argument pattern}\marg{replacement text}}
	|\edef| 命令与 |\def| 相同之处在于它也定义一个新宏。不同之处在于，它会把 \marg{replacement text} 一直展开，直到只剩下不可展开的记号为止（|\edef| $=$ 展开的定义）。
\begin{codeexample}[]
\def\a{3}
\def\b{2\a}
\def\c{1\b}
\def\d{value=\c}
\message{Macro `d' is defined to be `\meaning\d'}
\edef\d{value=\c}
\message{Macro `d' is e-defined to be `\meaning\d'}
\expandafter\def\expandafter\d\expandafter{\c}
\message{Macro `d' is defined to be `\meaning\d' using expandafter}
\end{codeexample}
	这个列表产生的日志文件输出为

	|Macro `d' is defined to be `macro:->value=\c '|

	|Macro `d' is e-defined to be `macro:->value=123'|

	|Macro `d' is defined to be `macro:->1\b ' using expandafter|

	\noindent 因此，|\def| 完全不展开，|\edef| 一直展开直到不能再展开，而 |\expandafter| 结构把 |\c| 展开一次，并把 |\d| 定义为这次展开的结果。

	虽然可以，但为 |\edef| 指定 \meta{argument pattern} 的情况可能并不常见，因为与 |\def| 不同，|\edef| 是立即展开的。不过它的工作方式相同：位置参数 |#1|、|#2|$,\dotsc,$ |#9| 会被替换为它们各自的实参。

	\marg{replacement text} 的展开方式与 \TeX\ 主记号列表的展开方式相同。

	那么，``一直展开直到只剩下不可展开的记号'' 究竟是什么意思？我们的例子表明，记号 |1|、|2|、|3| 是不可展开的，而宏 |\c|、|\b|、|\a| 是可以展开的。有一大类 \TeX\ 命令是不可展开的：任何赋值操作。例如
\begin{codeexample}[]
\edef\d{\count0=42}
\message{Macro `d' is defined to be `\meaning\d'}
\def\a{1234}
\edef\d{\advance\count0 by\a}
\message{Macro `d' is defined to be `\meaning\d'}
\end{codeexample}
\noindent 产生如下日志信息

|Macro `d' is defined to be `macro:->\count 0=42'| and

|Macro `d' is defined to be `macro:->\advance \count 0 by1234'|.

因此，赋值和算术操作是\emph{不可}展开的，它们作为可执行记号保留在新定义的宏中。这同样适用于 |\let| 和其他赋值操作。

	有趣的是，使用 |\if| $\dotsb$ |\fi| 的条件表达式\emph{是}可展开的，但这留待后面讲。

	还有一种方法可以把一个宏临时转换为不可展开的记号：|\noexpand| 宏。
\end{command}

\begin{command}{\noexpand\meta{expandable token}}
	|\noexpand| 命令只在 |\edef| 命令的 \marg{replacement text} 内部有用。当 |\edef| 一遇到 |\noexpand|，|\noexpand| 就会被移除，而 \meta{expandable token} 会被转换为不可展开的记号。因此，代码
\begin{codeexample}[]
\edef\d{Invoke \noexpand\a another macro}
\message{Macro `d' is defined to be `\meaning\d'}
\end{codeexample}
\noindent 产生终端输出

|Macro `d' is defined to be `macro:->Invoke \a another macro'|

因为 |\noexpand\a| 产生记号 `|\a|'（未展开）\footnote{键 \texttt{\textbackslash noexpand} 实际上被用来实现 \LaTeX\ 命令 \texttt{\textbackslash protect}：\LaTeX\ 的可移动参数概念就是用 \texttt{\textbackslash edef} 实现的。}。
\end{command}

\subsubsection{记号寄存器}
现在我们回到记号寄存器。正如我们在节~\ref{cmd:toks} 中已经看到的，记号寄存器存放一个记号列表。宏也在其 \marg{replacement text} 中存放一个记号列表，那么区别在哪里呢？有两点区别：
\begin{enumerate}
	\item 记号寄存器更快。
	\item 记号寄存器的内容\emph{永远不会}被展开。
\end{enumerate}
第一点我给不出具体数字——我只是在~\cite{texbook} 中读到的。但第二点允许进行展开控制。|\edef| 允许 ``无限'' 展开，而记号寄存器只允许顶层展开，就像 |\expandafter| 一样。但它们可以用比 |\expandafter| 更灵活（而且往往更高效）的方式来使用。

下面的例子演示了第二点。
\begin{codeexample}[]
\toks0={A \token list \a \b \count0=42 will never be expanded}
\edef\d{\the\toks0 }% the space token is important!
\message{Macro `d' is defined to be `\meaning\d'}
\end{codeexample}
\noindent 执行这段代码片段产生如下日志输出

|Macro `d' is defined to be `macro:->A \token list \a \b \count 0=42 will never be expanded'|.

因此，尽管我们用的是 |\edef|，|\toks0| 的内容仍被原样（未展开地）复制到了 |\d| 中。注意，|\the\toks0| 之后的空白记号确实很重要！\TeX\ 用它来界定整数 |0|。如果没有这个空白记号，它就会继续扫描，甚至越过 |\edef| 替换文本的边界（关于这种扫描的细节参见节~\ref{sec:variables}）。

这个例子非常简单，我们也可以像前面那样用 |\expandafter| 完成同样的事。但我们来试一个更难的：我们想拼装一个由不同片段组成的新宏。每个片段存放在一个宏中，而且不管出于什么原因，我们只希望对各个片段做顶层展开。并且：这些片段彼此并不相邻。我们可以用下面的示例列表来拼装目标宏。

\begin{codeexample}[]
\def\piecea{\a{xyz}}
\def\pieceb{\count0=42 }
\def\piecec{string \b}
\toks0=\expandafter{\piecea}
\toks1=\expandafter{\pieceb}
\toks2=\expandafter{\piecec}
\edef\d{I have \the\toks0 and \the\toks1 and \the\toks2}
\message{Macro `d' is defined to be `\meaning\d'}
\end{codeexample}

前三行定义了我们的片段。宏 |\piecea|、|\pieceb|、|\piecec| 各自包含一些在定义 |\d| 时不应被展开的记号。接下来三行把这些片段的顶层展开结果赋给记号寄存器。由于 |\toks0={\piecea}| 会把 `|\piecea|' 本身存进记号寄存器，所以这里需要用 |\expandafter|\footnote{我们本可以直接把一切写进记号寄存器，从而省去这些 \texttt{\textbackslash piece*} 宏。但我认为这个例子更贴近实际。}。然后，我们用 |\the\toks|\meta{number} 把一个记号列表的内容插入某处——在我们的例子中，就是插入宏 |\d| 已展开的替换文本中。因此，整个例子产生如下日志输出

|Macro `d' is defined to be `macro:->I have \a {xyz}and \count 0=42 and string \b '|.

\noindent 用（一大堆）|\expandafter| \emph{确实}也能得到完全相同的结果。不过别去试。


\subsubsection{宏定义命令小结}
除了 |\def| 和 |\edef| 之外，还有一些其他命令也可以用来定义宏（不过主要功能已由 |\def| 和 |\edef| 覆盖）。下面是其余的定义命令。
\begin{command}{\def\meta{\textbackslash macroname}\meta{argument pattern}\marg{replacement text}}
	定义一个名为 |\macroname| 的新宏，不展开 \marg{replacement text}，参见上文。
\end{command}
\begin{command}{\edef\meta{\textbackslash macroname}\meta{argument pattern}\marg{replacement text}}
	定义一个名为 |\macroname| 的新宏，完全展开 \marg{replacement text}（参见上文）。
\end{command}
\begin{command}{\let\meta{\textbackslash newmacro}=\meta{token}}
	定义或重新定义 |\newmacro|，使其等同于 \meta{token}。例如，|\let\a=\b| 会创建宏 |\b| 的一个新副本。该副本名为 |\a|，它将与 |\b| 具有完全相同的 \marg{replacement text} 和 \meta{argument pattern}。

	\meta{token} 也可以不是宏，例如可以是一个具名寄存器或单个字符。
\end{command}

\begin{command}{\gdef\meta{\textbackslash macroname}\meta{argument pattern}\marg{replacement text}}
	|\global\def| 的简写。它全局定义 |\macroname|，与当前作用域无关。

	应当避免某个宏在全局命名空间和某个局部作用域中同时存在、却具有不同含义的情况。节~\ref{sec:scopes} 进一步解释了作用域。
\end{command}
\begin{command}{\xdef\meta{\textbackslash macroname}\meta{argument pattern}\marg{replacement text}}
	|\global\edef| 的简写。它全局定义 |\macroname|，与当前作用域无关。

	应当避免某个宏在全局命名空间和某个局部作用域中同时存在、却具有不同含义的情况。节~\ref{sec:scopes} 进一步解释了作用域。
\end{command}

\begin{command}{\csname\meta{expandable tokens}\textbackslash endcsname}
	这个命令并不是定义宏，而是定义宏的\emph{名字}。``cs'' 表示 ``control sequence''（控制序列）。|\csname| 与 |\endcsname| 这一对用 \meta{expandable tokens} 来定义一个控制序列名（即宏名）。控制序列字符 `|\|' 会由 |\csname| 自动加在前面。\footnote{实际上，这里使用的是 \texttt{\textbackslash escapechar} 的内容。如果它的值为 -1，则不加任何字符。凡是 \TeX\ 命令会插入反斜杠的地方，都同样如此。}
\begin{codeexample}[]
\def\macro{Content}
This here is normal usage: `\macro'.
This here uses csname: `\csname macro\endcsname'.
\end{codeexample}
	\noindent 这个例子表明，|\csname|\meta{expandable tokens}|\endcsname| 实际上就等同于你直接写出 |\|\meta{expandable tokens}——但 |\csname| 结构允许在宏名中使用多得多的记号：
\begin{codeexample}[]
\expandafter\def\csname a01macro with.strange.chars\endcsname{Content}
I use a strange macro. Here is it: `\csname a01macro with.strange.chars\endcsname'.
\end{codeexample}
	\noindent 这个例子用 |\expandafter| 把 |\csname| 展开一次。|\csname| 的顶层展开结果是一个记号，即那个控制序列名。然后，用 |\def| 以这个准备好的宏名来定义一个宏。

	当 |\csname| 被展开时，它会解析直到下一个 |\endcsname| 之前的所有记号。这些记号会被一直展开，直到只剩下不可展开的记号（就像 |\edef| 那样）。得到的字符串将被用来定义一个宏名（前面自动加上控制序列字符 `|\|'）。由于 \meta{expandable tokens} 会被展开，这使得我们可以使用 ``间接的'' 宏名：
\begin{codeexample}[]
\def\macro{onetwothree}
\expandafter\def\csname macro\macro\endcsname{Content}
I have just defined \expandafter\string\csname macro\macro\endcsname
with replacement text `\csname macro\macro\endcsname'.
\end{codeexample}
	\noindent 我想这个例子不言自明，除了下面要讲的 |\string| 命令之外。

	正是由于这种灵活性，|\csname| 被用来实现 \TeX\ 中所有（？）现有的键值宏包。
\end{command}

\begin{command}{\string\meta{\textbackslash macro}}
	这个命令并不定义宏。相反，它把一个宏的名字作为一系列独立记号返回，其中包括控制序列记号 `|\|'。
\begin{codeexample}[]
\def\macro{Content}
I have just defined `\string\macro' using `\string\def'.
\end{codeexample}

	你也可以把 |\string| 用在其他记号上——例如字符。这没什么坏处，该字符会原样返回。
\end{command}

\subsubsection{调试工具——理解并追踪 \TeX\ 的所作所为}
\begin{command}{\message\marg{tokens}}
\end{command}
\begin{command}{\meaning\meta{\textbackslash macro}}
\end{command}
\begin{command}{\tracingmacros=2}
\end{command}
\begin{command}{\tracingcommands=2}
\end{command}
\begin{command}{\tracingrestores=1}
\end{command}

\subsection{变量的作用域}
\label{sec:scopes}
每种编程语言都有作用域的概念：它限制变量或例程的生效范围。然而，\TeX\ 的作用域机制并不是为编程而设计的——\TeX\ 是一种排版语言。许多编程语言，如 |C|、|C++|、|java| 或众多脚本语言，是根据变量定义所在的位置来界定其作用域的。例如下面这段 |C| 代码
\begin{codeexample}[code only]
int i = 42;

{
	++i;
	int i = 5;
}
\end{codeexample}
\noindent 把外层 |i| 的值改为 |43|。内层 |i| 的值是 |5|，但一遇到闭花括号它就会被删除。甚至可以同时访问内层 |i| 变量和外层 |i| 变量的值。

在 \TeX\ 中，花括号也用于界定作用域。但是：\TeX\ 不仅会在作用域结束时销毁该作用域内定义的任何变量（宏），还会\emph{撤销}在该作用域内施加的任何改动。
\begin{codeexample}[]
\def\i{42}
{
	\def\i{43}
	\def\b{2}
}
The value of \textbackslash i is now \i.
\end{codeexample}
\noindent 上面的列表定义了 |\i|，进入一个局部作用域（一个 \TeX\ ``分组''）并修改了 |\i|。然而，根据 \TeX\ 的作用域规则，从局部分组返回后，旧的程序状态会被\emph{完全}恢复！对 |\i| 的改动和对 |\b| 的定义都不会保留。寄存器的改动或其他赋值也同样如此。

\TeX\ 分组可以通过三种方式之一创建：使用花括号\footnote{或者其他具有正确类别的记号，参见~\cite{texbook}。}、使用 |\begingroup| 或使用 |\bgroup|。花括号很少用来界定 \TeX\ 分组，因为其他命令更灵活。如果使用花括号，它们必须配对——不允许出现不匹配的花括号。
\begin{command}{\begingroup}
	开始一个新的 \TeX\ 分组（一个局部作用域）。该作用域将一直有效，直到被 |\endgroup| 关闭。|\endgroup| 命令可以稍后出现在主记号列表中。
\end{command}
\begin{command}{\endgroup}
	结束一个用 |\begingroup| 打开的 \TeX\ 分组。
\end{command}
\begin{command}{\bgroup}
	|\begingroup| 的一种特殊变体，它也可以用来界定传给 |\hbox| 或 |\vbox| 的参数（也就是说，在这种场合下可以省去配对的花括号）。

	|\bgroup| 宏还有一个用处：测试紧随其后的字符是否为开花括号（参见 |\futurelet|）。

	如果只是需要打开一个 \TeX\ 分组，应当优先使用 |\begingroup|。
\end{command}
\begin{command}{\egroup}
	关闭前面的 |\bgroup|。
\end{command}

\TeX\ 不知道如何写入外层作用域的宏——最顶层（全局）作用域除外。如果需要编写复杂结构，这一限制相当棘手：局部变量应当在局部分组内声明，但对结构的改动却应当写回外层分组。没有直接的办法做到这一点（全局变量除外）。

\subsubsection{全局变量}
\TeX\ 只知道 ``全局'' 变量和 ``局部'' 变量。局部变量在其所在分组结束时会被删除。所有局部赋的值也会在分组结束时恢复为旧值。

而全局变量则在\emph{每个}作用域中都保持同一个值。通常，最顶层作用域就是用于全局变量的那个作用域：如果你在 \TeX\ 文档中定义任何东西，就是在全局作用域上添加命令。也可以显式地在全局作用域中进行赋值或定义。

\begin{command}{\global\meta{definition or assignment}}
	紧随 |\global| 之后的定义将在全局范围内进行。
\begin{codeexample}[code only]
{
	\global\def\a{123}
	\global\advance\count0 by3
	\global\toks0={34}
}
\end{codeexample}
\end{command}

\begin{command}{\globaldefs=\mchoice{-1,0,1} (initially 0)}
	我引用~\cite{texbook}：``如果在赋值时 |\globaldefs| 参数为正，则自动隐含 |\global| 前缀；但如果赋值时 |\globaldefs| 为负，则忽略 |\global| 前缀。如果 |\globaldefs| 为零（通常如此），则是否出现 |\global| 决定了是否进行全局赋值。''
\end{command}


\subsubsection{把改动传送到外层分组}
有若干种方法可以把改动 ``传送'' 到外层作用域。有些是复制操作，有些则需要在作用域结束后重新施加这些改动。所有这些方法都可以借助展开控制来实现。

我们从需要跨越分组结束而保留下来的宏定义讲起。我看到的方法有：
\begin{itemize}
	\item 把宏复制到一个全局临时变量（甚至记号寄存器）中，并在作用域结束后取回该值。
\begin{codeexample}[code only]
\def\initialvalue{0}
{
	% do something:
	\def\initialvalue{42}
	\global\let\myglobaltemporary=\initialvalue
}
\let\initialvalue=\myglobaltemporary
\end{codeexample}
	其思路是 |\myglobaltemporary| 只被临时使用；它的值始终是未定义的，可以随时被覆盖。这样就能使用局部变量 |\initialvalue| 了。

	请注意，不要在全局和局部同时使用同一个变量。这会让 \TeX\ 感到困惑，并导致运行时变慢。

	\item 把结果 ``偷偷运出'' 当前分组。我是从 Mark Wibrow 和 Till Tantau 编写的~\cite{tikz} 的实现中得知这一技巧的。其思路是使用若干个 |\expandafter| 和一个 |\def|，在分组结束后立刻重新定义该宏：
\begin{codeexample}[code only]
\def\smuggle#1\endgroup{%
	\expandafter\endgroup\expandafter\def\expandafter#1\expandafter{#1}%
}

\begingroup
	\def\variable{12}
	\edef\variable{\variable34}
	\edef\variable{\variable56}
	\smuggle\variable
\endgroup
\end{codeexample}
	这一技巧依赖于用 |\begingroup| 开始、|\endgroup| 结束的分组，因为 |\def| 不允许不配对的花括号。其效果是，在所有那些 |\expandafter| 之后，\TeX\ 遇到的记号列表为

	|\endgroup\def\variable{123456}|

	即在分组结束时。

	\item 使用 aftergroup 栈。\TeX\ 有一个大小有限的特殊记号栈，可用于在分组结束后重新插入记号。不过，只有当需要传送的记号数量很小且固定时（比如至多三个），这种方法才高效。它的做法是在每个记号前加上 |\aftergroup|，细节参见~\cite{texbook}。
\end{itemize}

有时还需要把其他变量复制到作用域之外。当然，用临时全局变量的那招总是管用的。但也可以定义一个宏，其中包含施加所需改动的命令，然后把这个宏传送出作用域。

\subsection{分支}
\label{sec:branching}

这里我们讨论 \TeX\ 的一些可用分支结构，重点放在涉及数字和记号的条件上。

\begin{command}{ifnum\meta{count/integer number}=\meta{count/integer number}\meta{true-block}\textbackslash else\meta{false-block}\textbackslash fi}
|\ifnum| 比较整数或整数寄存器（|\count| 寄存器），它包含两个分支：条件为真时执行一个，条件为假时执行另一个：
\begin{codeexample}[]
\ifnum1=2 % this space is important.
This is shown if above were true.
\else
This is shown if above results to false.
\fi
\end{codeexample}

	注意，|\else| 及其 \meta{false-block} 是可选的。
\end{command}

\begin{command}{ifdim\meta{dimen/fixed point number}=\meta{dimen/fixed point number}\meta{true-block}\textbackslash else\meta{false-block}\textbackslash fi}
	与 |\ifnum| 类似，|\ifdim| 比较两个定点数或 |\dimen| 寄存器。数字必须带单位。
\begin{codeexample}[]
\ifdim1pt=2pt % this space is important.
This is shown if above were true.
\else
This is shown if above results to false.
\fi
\end{codeexample}
\end{command}

\begin{command}{ifx\meta{token1}\meta{token2}\meta{true-block}\textbackslash else\meta{false-block}\textbackslash fi}
|\ifx| 稍微复杂一些：它比较两个\emph{记号}，只比较到它们的第一层展开。
\begin{codeexample}[]
\def\empty{\empty}
\ifx\empty\empty %
This is shown if the two tokens have equal expansion.
\else
This is shown if the two tokens expand to something different.
\fi
\end{codeexample}

这里，我们定义记号 |\empty| 作为 |\empty| 的替换，然后比较这两个记号在第一层展开下是否相等。注意，这个定义其实是没有意义的。如果 \TeX{} 真的把整个展开都走完——也就是说我们把 |\empty| 放到别处——它就会无限地展开下去。不过，|\ifx| 只做第一层展开并加以比较。因此该语句求值为真。

请看下面的例子：
\begin{codeexample}[]
\def\empty{\relax}
\ifx\empty\relax %
This is shown if the two tokens have equal expansion.
\else
This is shown if the two tokens expand to something different.
\fi
\end{codeexample}

乍一看，这应该求值为真：|\empty| 被定义为 |\relax| 的替换。但它并没有。为什么？

|\empty| 被展开为 |\relax|，然而展开后的 |\relax| 含义不同，即 ``停止扫描''，而不再是 |\relax| 本身。因此它们不同，该语句为假！如果 |\ifx| 中的展开一直进行到最大程度，两者就会相等；但在只比较第一层展开的情况下并非如此。
\end{command}

\begin{command}{if\meta{token1}\meta{token2}\meta{true-block}\textbackslash else\meta{false-block}\textbackslash fi}
	|\if| 比较与 |\ifx| 条件密切相关，但有一个主要区别：它会不断展开记号，直到找到接下来两个不可展开的记号。如果这两个记号相同，就展开为 \meta{true-block}，否则展开为 \meta{false-block}。

	|\if| 条件应当小心使用，因为它可能产生不希望的效果。只有在你清楚自己在做什么时才使用它。

	一个有用的例子是：当你\emph{确知}某个宏至多包含一个字符，而你想测试它是否为某个特定字符时：
\begin{codeexample}[]
\def\choice{a}
\if b\choice
  This is shown for the `b' choice.
\else
  This is shown for all other choices.
\fi
\end{codeexample}
\end{command}

\begin{command}{iftrue\meta{true-block}\textbackslash else\meta{false-block}\textbackslash fi}
	一个 ``条件''，它总是执行 \meta{true-block}。
\end{command}
\begin{command}{iffalse\meta{true-block}\textbackslash else\meta{false-block}\textbackslash fi}
	一个 ``条件''，它总是执行 \meta{false-block}。
\end{command}

\subsubsection{布尔变量}
\begin{command}{\newif\meta{if-name}}
	你可以借助 |\newif\ifsupermanmode| 声明一个新的 ``布尔变量'' |\ifsupermanmode|。之后，你可以用 |\supermanmodetrue| 和 |\supermanmodefalse| 这两个开关来给该布尔变量赋值，并用 |\ifsupermanmode| 来检查它。

	\meta{if-name} 必须以 |\if| 开头（以便支持对嵌套的 |\if...\fi| 配对的扫描，参见下文）。
\end{command}

\subsubsection{条件语句的特殊情况}
每当你使用 |\if|\ldots\ 及其同类命令时，应当了解以下特性：
\begin{enumerate}
	\item |\if...\else...\fi| 是可展开的（包括其中的每一个宏 |\if...|、|\else| 和 |\fi|），这意味着你甚至可以在 |\edef| 内部使用它：
\begin{codeexample}[]
\def\choice{a}
\edef\temp{The choice is \if a\choice `a'\else not `a'\fi}
We have now \texttt{\string temp=\meaning\temp}.
\end{codeexample}

\begin{codeexample}[]
\def\shownexttoken#1{The next token is `\texttt{\string#1}'.}
\def\mymacro{%
 	\ifnum1=1 %
 		\expandafter\shownexttoken%
 	\fi%
}%
\mymacro 23
\end{codeexample}
	这个例子有点棘手。如果没有 |\expandafter|，会发生什么呢？那样的话，|\shownexttoken| 会以 |#1=\fi| 被调用。这会导致错误，因为 |\fi| 会丢失；而且它还会破坏效果，因为我们并不希望 |\fi| 被看到——我们期望的是 |#1=2|。|\expandafter| 先展开 |\fi|（它只是把 |\fi| 移除，没有其他效果），这样在上面的例子中 |\shownexttoken| 就会看到记号 |2|。如果那里是一个 |\else| 分支而不是 |\fi|，这同样有效。

	\item 一般来说，你应当确保与之配对的 |\else| 或 |\fi| 记号是 ``可直接到达的''，也就是无需经过记号展开。

	其背景是：\TeX\ 工作在基于记号的层面上：每当它遇到一条 |\if|\ldots\ 语句，就对它求值并扫描记号，以找到配对的结束部分（一个 |\else| 记号或一个 |\fi| 记号）。但在这个扫描过程中它不会展开记号，不过它会计数嵌套的 |\if...\fi| 配对！因此，如果不小心，它可能会搞混，你的条件语句就会出错。
\end{enumerate}

\subsection{循环}
\label{sec:loops}
如你所见，在 \TeX{} 中我们对记号展开有着非常精细的控制。这使得甚至可以借助递归来构造循环。本质上，一个循环由以下几部分组成：
\begin{itemize}
 \item 计数器，或者更一般地说，条目列表
 \item 递增器，或者更一般地说，下一条目选取器
 \item 阈值，或者更一般地说，一个结束列表标记
 \item 分别对阈值或结束标记进行检查
\end{itemize}

通览上面各节，我们会意识到这些要素其实都已具备：我们知道计数器，也知道分支。剩下要弄清楚的只是如何创建这些循环的具体细节。下面我们将详细展示两种情形：计数循环和对条目列表的循环。

一般来说，对于通过递归实现的循环，我们需要两个定义：一个用于循环开始，另一个用于循环步进。

\subsubsection{计数循环}
\label{sec:counting:loops}
对于计数循环，我们需要一个计数器 |\count0|、一个递增器 |\advance|、一个阈值 |3|，以及一个检查是否达到阈值的判断 |\ifnum\count=10|。


\begin{codeexample}[]
\long\def\countingloop#1 in #2:#3#4{%
	#1=#2 %
	\loopcounter{#1}{#3}{#4}%
}

\long\def\loopcounter#1#2#3{%
	#3%
	\ifnum#1=#2 %
	\else%
		\advance#1 by1 %
		\loopcounter{#1}{#2}{#3}%
	\fi%
}
\countingloop{\count0} in 0:{3}{%
	The current value is `\the\count0'\par
}
\end{codeexample}

上面这个例子有一些微妙之处：
\begin{itemize}
 \item 我们在例子中放了很多 \%。为什么？请注意，每当 \TeX{} 扫描一个数字时——例如在 |#1=#2| 的情形——它会逐个记号、也就是逐位数字地继续扫描，直到确信数字已经结束，即使越过空白也不停，甚至在宏本身可能不再表示数字时还会展开宏。因此，\% 告诉 \TeX{} 停止扫描。一般来说，放置 \% 来告诉 \TeX{} 停止扫描更多数字是个好习惯。不过也有一些例外：在 |\advance#1 by1| 的情形下，中间应保留一个空白，|\ifnum#1=#2| 的情形也是如此。
 \item 我们在 |\countingloop{\count0} in 0:{3}| 中把阈值 |3| 放在花括号里。为什么？否则，如果给出的阈值是例如 |10|，\TeX{} 只会识别记号 |1|，然后被剩下的 ``多余'' 参数 |0| 绊住。这是因为对 \TeX{} 来说一个字母就是一个记号。因此两个字母就是两个记号，若不用分组括起来，就会变成两个参数。这里我们必须把阈值分组，以明确我们的意思。
 \item 最后一点只有在开启调试时才会清楚：由于循环是通过递归实现的，也就是一次接一次地展开，直到达到某个阈值，所以在上面的情形中我们最后会留下一大堆 |\fi|。如果在调用 |\countingloop| 之前放置 |\tracingmacros=2 \tracingcommands=2| 并查看日志文件，这一点就会很明显。这很糟糕，因为 \TeX{} 会为每一个 |\if|\ldots|\fi| 序列保留一个打开的栈帧。如果现在要对一万个条目做循环……
 \item 使用系统计数器之一（这里是 |\count0|）不是个好习惯，因为你永远无法确定它没有被用于别处或在别处被修改。例如，当页面排满时，\TeX{} 会中断当前的记号序列，去处理创建新页面和结束旧页面的事务，在此过程中会改变 |\count0|。因此，我们也应当创建自己的计数器。
\end{itemize}

因此，我们把例子修改如下：
\begin{codeexample}[]
\long\def\countingloop#1 in #2:#3#4{%
	#1=#2 %
	\loopcounter{#1}{#3}{#4}%
}

\long\def\loopcounter#1#2#3{%
	#3%
	\ifnum#1=#2 %
		\let\next=\relax%
	\else
		\advance#1 by1 %
		\def\next{\loopcounter{#1}{#2}{#3}}%
	\fi
	\next
}
\newcount\ourcounter

\countingloop{\ourcounter} in 0:{3}{%
	The current value is `\the\ourcounter'\par
}
\end{codeexample}

原则上，输出没有任何变化。但是请注意，我们引入了宏 |\next|：它要么递归到下一层——但这是在 |\fi| 语句给出之后——要么干脆只包含 |\relax| 来结束递归。此外，我们还声明了一个名为 |\ourcounter| 的新计数器，它是安全无虞的。

最后，我们简要地总结一下具体发生了什么：
\begin{enumerate}
 \item |\countingloop|\ldots 被展开为一个赋值 |#1=#2| 和另一个宏 |\loopcounter|\ldots。
 \item 赋值完成：|\ourcounter| 被设为起始值 |0|。
 \item\label{countingloop:loopstart} 实际的循环宏被展开为命令块——打印当前值——和一条 if 语句。
 \item 打印当前值。
 \item 把 |\ourcounter| 与阈值 |3| 比较，然后……
 \begin{itemize}
	 \item 为假，即该 if 语句被展开为一条 |\advance| 语句，随后把 |\next| 定义为对同一宏循环的又一次调用。
	 \item 为真，即 |\next| 被设为 |\relax|。
 \end{itemize}
 \item 该语句仍为假：|\advance| 把 |\ourcounter| 加一，现在是 |1|。|\next| 被设为循环宏。
 \item 循环宏再次展开，转到步骤~\ref{countingloop:loopstart}。|\ourcounter| 是……|2|……|\ourcounter| 是 |3|。
 \item 现在该语句为真：|\next| 被展开为 |\relax|，什么也不发生。
\end{enumerate}

\subsubsection{对条目列表的循环}
\label{sec:loops:over:list:of:items}
对条目列表做循环非常类似，只是我们要用 |\ifx| 代替 |\ifnum|，并用某个结束标记代替阈值。不过，我们如何指定列表本身呢？我们来构造一个逗号分隔的列表，例如 |{a,b,c,d}|，并把结束标记称为 |\listingloopENDMARKER|。

\begin{codeexample}[]
\def\listingloopENDMARKER{\par \listingloopENDMARKER}
\long\def\listingloop#1in#2#3{%
	\looppicker{#1}{#3}#2,\listingloopENDMARKER,%
}%
\long\def\looppicker#1#2#3,{%
	\def\tempitem{#3}%
	\ifx\tempitem\listingloopENDMARKER
		\let\next=\relax%
	\else
		\def#1{#3}%
		#2%
		\def\next{\looppicker{#1}{#2}}%
	\fi
	\next
}%
\listingloop\x in{a,b,c,,d,e}{%
	The current item is `\x'
}
\end{codeexample}

同样，我们来阐明其中包含的微妙之处：
\begin{itemize}
 \item 我们把 |\listingloopENDMARKER| 定义为替换其自身。这之所以可行，是因为 |\ifx| 只比较第一层展开，参见节~\ref{sec:branching}。
 \item 我们似乎在 \ldots|#1in#2|\ldots 中少了一个空白。不过，记号总是以一个额外的空白结尾，因为 |\xin| 不等于 |\x in|。因此这里不需要空白，而且多于一个空白很可能会被吞掉。
 \item 定义 |\looppicker#1#2#3,|\ldots 有三个参数，但递归调用 |\looppicker{#1}{#2}| 却只给出两个参数！？这正是让这种列表成为可能的奥妙所在！\TeX{} 实际上会越过所给记号的范围继续扫描，以取得第三个参数。实际上，我们是从给定的列表中一口一口、一个条目一个条目地咬下来。这全都是因为我们在 |\looppicker| 宏的定义中声明了一个额外的 |,|——逗号即条目分隔符。循环宏的展开总会从拼接到其末尾的列表中再多取一个条目，直到到达 |\ENDMARKER|。这个结束标记是在循环开始时加到列表最末尾的，循环就在那里停止。
\end{itemize}

\subsection{关于 \TeX 的更多内容}
本文档远未完善。我建议阅读 \cite{schwartz}（德语，在线版）或 \cite{texbook}（实体书）中关于条件表达式的内容。关于循环的提示可在 \PGFPlots 的手册 \cite{pgfplots} 和 \PGF 的手册 \cite{tikz} 中找到。此外，\PGFPlots\ 和 \PGF\ 附带了大量实用函数，它们记录在源代码 |.code.tex| 文件中。

\section{使用 \PGF\ 键的键值处理综述}

对每个 \TeX\ 宏包而言，键值输入都是最重要的事情之一。~\cite{keyvalintro} 对各种键值宏包（其中包括 |xkeyval| 和 |pgfkeys|）做了很好的综述与概述。

除了上面提到的论文以及~\cite{tikz} 中关于 |pdfkeys| 的详尽参考手册之外，我在这里对 |pgfkeys| 做一个简要综述。目标读者主要是宏包作者或宏程序员。本节应能让你为 \PGFPlots\ 和 \PGF\ 定义自己的用户界面和样式。它也应能增进对 |pgfkeys| 及其用法的理解。我还会讨论键过滤的话题，这主要对宏包作者有用。

宏包 |pgfkeys| 可作为独立宏包使用：|\usepackage{pgfkeys}|。不过，我认为你永远不需要显式加载它，因为 \PGF\ 反正会被加载，而 \PGF\ 总是会加载 |pgfkeys|。

它提供两套用户界面。我认为最佳做法是兼取两者之长；不过只用其中一套或许也足够。因此，我会分别讨论这两套，并在后面给出最佳实践的建议。

\subsection{\PGF\ 键的底层（直接）\API}
我们从 \PGF\ 键的底层 \API\ 讲起。它由若干宏组成，允许定义键、赋值以及取回键值。

\begin{command}{\pgfkeyssetvalue\marg{/key path/key name}\marg{value}}
	这个宏（重新）定义一个键。

	它（几乎）等同于某种宏定义

|\expandafter\def\csname key@|\meta{/key path/key name}|\endcsname|\marg{value};

	也就是说，它把 \meta{value} 存进一个新宏，以便能在 \TeX\ 的哈希表中以常数时间查找该键。
	请注意，与 |xkeyval| 等其他键值宏包不同，用于存储值的底层宏名\emph{并不}属于 \PGF Keys \API\footnote{注意，\texttt{key@} 与 \PGF Keys 无关。}——请使用 |\pgfkeysgetvalue| 及其同类命令来访问值。

	可能的键的数量，其唯一限制是 \TeX\ 哈希表的大小（它非常大）。

	你或许好奇斜杠 `|/|' 是什么意思。熟悉 \PGF/\PGFPlots\ 的用户知道，存在某种 ``键路径''，用来限定 \meta{key name}。\meta{key path} 的作用是提供一个命名空间，使得许许多多同名的键可以和平共处、互不干扰——只要使用了正确的 \meta{key path}。它可以看作一条（unix）文件路径：你可以有许多同名文件，只要这些文件位于不同的目录中（即具有不同的路径）。

	你\emph{应当总是}提供键路径，而且强烈建议使用一个不同于仅仅 `|/|' 的键路径。

	\meta{value} 可以是任何东西；它只是被原样存储。它甚至可以包含 |#|。
\end{command}

\begin{command}{\pgfkeysgetvalue\marg{/key path/key name}\marg{\textbackslash macro}}
	正如你可能已经猜到的，这个宏允许取回某个键的值并存入 \meta{\textbackslash macro}。

	既然我们已经了解了 |\pgfkeyssetvalue| 和 |\pgfkeysgetvalue|，也可以给出一个例子：
\begin{codeexample}[]
\pgfkeyssetvalue{/notes/key}{abc}

\pgfkeysgetvalue{/notes/key}\temp

The value of key \texttt{/notes/key} is `\temp'.
\end{codeexample}

	这两个键周围并没有多少奥妙；它就像是一次哈希表访问，只不过键有一些特殊的命名约定（源自键路径）。注意，既然 ``哈希表访问'' 正是 \TeX\ 处理宏时时刻在做的事，我们本可以用 |\def| 和合适的 |\let| 命令（也许再配上 |\csname...\endcsname|）来取代 |\pgfkeyssetvalue|/|\pgfkeysgetvalue| 这一对。而 \PGF Keys 的优势，要等到我们在下一节考察其高层用户界面时才真正显现出来。

	请注意，由于 |\pgfkeyssetvalue| 本质上等同于一个合适的 |\def|，该赋值是当前 \TeX\ 分组局部的。换句话说：该赋值会被下一个闭花括号、或下一个 |\endgroup|、或下一个 |\end|\marg{environment} 撤销。
\end{command}

\begin{command}{\pgfkeyslet\marg{/key path/key name}\marg{\textbackslash macro}}
	它本质上与 |\pgfkeyssetvalue| 相同，区别只在于键的值已经在 \meta{\textbackslash macro} 内部准备好了：

\begin{codeexample}[]
\def\something{abc}
\pgfkeyslet{/notes/key}{\something}

\pgfkeysgetvalue{/notes/key}\temp

The value of key \texttt{/notes/key} is `\temp'.
\end{codeexample}

	正如 |\pgfkeyssetvalue| 归结为 |\def| 一样，|\pgfkeyslet| 归结为 |\let|。
\end{command}

\begin{command}{\pgfkeysvalueof\marg{/key path/key name}}
	它本质上等同于 |\pgfkeysgetvalue|\marg{/key path/key name}\marg{\textbackslash macro} \meta{\textbackslash macro}；也就是说，它展开为存放在某个键中的值。
\begin{codeexample}[]
\pgfkeyssetvalue{/notes/key}{abc}

The value of key \texttt{/notes/key} is `\pgfkeysvalueof{/notes/key}'.
\end{codeexample}

	不过，这个键有一个主要优点：它可以用在 |\edef| 内部（因为它是完全可展开的）：
\begin{codeexample}[]
\pgfkeyssetvalue{/notes/key}{abc}

\edef\temp{--- \pgfkeysvalueof{/notes/key} ---}

The value of key \texttt{/notes/key} along with dashes is \temp.
\end{codeexample}

	它归结为一个合适的 |\csname ... \endcsname|。因此，如果该键恰好未定义，它就展开为 |\relax|（参见下文的 |\pgfkeysifdefined|）。
\end{command}

\begin{command}{\pgfkeysdef\marg{/key path/key name}\marg{macro body}}
	它是 |\pgfkeyssetvalue| 的一个变体。不过，它有一个实质性的区别，而只要我们还在讨论底层 \API，这个区别就显得没有来由。它定义一个所谓的 code-key（代码键）。

	code-key 是可执行的宏。它们接受一个参数，并对它做些处理。给这样一个键 ``赋值''，就相当于以 ``合适的'' 方式调用 \meta{macro body}。

	这个宏调用的结果是一个名为 \meta{/key path/key name/}|.@cmd| 的新键。而那个键又被存为可执行宏。该宏等价于下面这样的定义（当然名字除外）：

	|\def\macro#1\pgfeov|\marg{macro body}.

	这个宏（使用 |\pgfkeyslet|）被存放在 \meta{/key path/key name/}|.@cmd| 之下。

	我们可以用 |\pgfkeysgetvalue| 和/或 |\pgfkeysvalueof| 来访问这个特殊的键\footnote{注意，后缀 |/.@cmd| 是 \PGF Keys 公开 \API\ 的一部分，所以使用它并不算什么 hack。}，尽管它的用处要到本文档后面才会更明显：

\begin{codeexample}[]
\pgfkeysdef{/notes/code key}{Expansion with value #1---X.}%

We ``assign a value'' or ``execute the code key'' (which is equivalent):
\pgfkeysvalueof{/notes/code key/.@cmd}abc\pgfeov
\end{codeexample}
	注意，在这种情况下我们\emph{必须}用 |\pgfeov| 来结束参数列表。我们本可以把参数放进花括号里，但仍必须提供 |\pgfeov|，正如我们必须加上后缀 |/.@cmd| 一样。

\end{command}

\begin{commandlist}{%
	\pgfkeysifdefined\marg{/key path/key name}\marg{true case}\marg{false case},
	\pgfkeysifassignable\marg{/key path/key name}\marg{true case}\marg{false case}}
	这些键提供基于某个键是否存在或其类型的条件判断。细节请参阅~\cite{tikz} 中的参考手册。
\end{commandlist}

\subsection{\PGF Keys 的标准 \API}
\begingroup
\pgfkeys{
	/codeexample/newline=,
}
既然我们已经看到如何在相当低的抽象层次上访问 \PGF Keys 所定义的东西，下面我们用更高的层次来重复同样的事情。本节解释 \PGF Keys 的标准 \API；这是轻松定义和维护键的方式，它同时也是最终用户界面。

\PGF Keys 用其标准 API 解决若干使用场景：
\begin{enumerate}
	\item 简单的键值存储（即放入和取出），
	\item 每当使用某个键时就能执行某些（复杂）操作的 code-key，
	\item 对键值工具本身的配置和修改。
\end{enumerate}
所有这些都可以用同一个宏来完成：

\begin{command}{\pgfkeys\marg{comma--separated key--value pairs}}
	这个键构成 \PGF Keys 的公开 \API。它接受任意数量的键值对，以逗号分隔。
\end{command}

我们从一个例子开始：
\begin{codeexample}[]
% key definition:
\pgfkeys{
	/notes/key/.initial=abc,
}

The value of key \texttt{/notes/key} is `\pgfkeysvalueof{/notes/key}'.

% key usage:
\pgfkeys{
	/notes/key =  efg  ,
}

The value of key \texttt{/notes/key} is `\pgfkeysvalueof{/notes/key}'.
\end{codeexample}

有几点看似不言自明，我这里简要确认一下确实如此：键名前后以及值前后的空白都会被去掉。此外，尾随逗号会被忽略。请注意，尾随逗号是一种最佳实践：总是加上尾随逗号。这会大大简化后续添加键的工作（我已经记不清有多少次加了一个键，却纳闷为什么它没被正确识别，直到发现少了那个逗号）。把加尾随逗号变成习惯就好。另一个好习惯是正确缩进代码，也就是每个新行都缩进一个制表位。每行只写一个键也是个好主意，不过这些都只是可选的。

查看实际代码时，第一个让人觉得奇怪的地方是后缀 `|/.initial|'。这其实是 \PGF Keys 一套一致的新机制：这些后缀允许配置和修改它们所作用的键。它们被称为 ``键处理器''（key handler）。每当你看到 \meta{key path/key name} 后面跟着 `|/.|\meta{handler}'，就可以有把握地断定：\meta{key path/key name} 即将被重新配置或修改。

掌握了键处理器，就掌握了对 \PGF Keys 的控制。下面我简要讨论最重要的几个处理器。

\begin{handler}{{.initial}=\marg{value}}
	键处理器 |/.initial| 定义一个新的 \meta{key}，并赋予其初始 \meta{value}。

	就此而言，它等价于 |\pgfkeyssetvalue|\marg{key}\marg{value}。

	通过 |/.initial| 定义的键，之后随时都可以用一个简单的赋值来设置（参见上面的例子）。

	因此，第一次定义需要这个后缀，之后的所有赋值只需用赋值来设置值即可。
\end{handler}

\begin{handler}{{.code}=\marg{body}}
	这个键处理器定义一个新的 code-key \meta{key}，以 \meta{body} 作为其结果。
\begin{codeexample}[]
\pgfkeys{
	/notes/code key/.code={Expansion with value #1---X.},
}%

用简单 API 执行该键：
\pgfkeysvalueof{/notes/code key/.@cmd}abc\pgfeov

... 用标准 API 中的赋值来执行该键：
\pgfkeys{/notes/code key=abc}
\end{codeexample}
	我们看到，给一个 code key 赋值就意味着执行 \meta{body}，其中 |#1| 被设为在 \API\ 中所赋的值。

	通过 |/.code| 定义的键等价于通过 |\pgfkeysdef| 定义的键。

	注意，参数 \meta{body} 可以用花括号括起来，但并非必须如此：
\begin{codeexample}[code only]
\pgfkeys{
	/notes/code key/.code={Expansion with value #1---X.},
	/notes/code key/.code=Expansion with value #1---X.,
}%
\end{codeexample}
	这是 \PGF Keys 一个常见特性：任何形式的赋值都可以用花括号，但并非必须如此。只有当所赋参数（在我们这里即 \meta{body}）包含 \PGF Keys 的控制字符（即 |=| 或 |,|）时，才需要用花括号。
\end{handler}

\begin{handler}{{.style}=\marg{option list}}
	这个键处理器定义一个新的 code-key \meta{key}，每当它被赋值（使用）时就设置 \meta{option list} 中的所有选项。

	样式的定义方式很简单：它们只是用 \meta{option list} 调用 |\pgfkeys|（嗯，差不多——见下文）。不过，它们在任何应用中都极具表现力。

\begin{codeexample}[]
\pgfkeys{
	/notes/A/.initial=,
	/notes/B/.initial=,
	/notes/my style/.style={
		/notes/A={#1},
		/notes/B={#1},
	},
}%

定义已完成。现在给该样式赋值：
\pgfkeys{
	/notes/my style=42
}
OK. Value of A=\pgfkeysvalueof{/notes/A}, value of B=\pgfkeysvalueof{/notes/B}.
\end{codeexample}
	我们的例子是样式的一个非常简单的应用：它设置了一堆其他选项。

	注意，\meta{option list} 可以依赖于 |#1|。

	到目前为止，本文档总是提供完全限定的键路径。不过 |/.style| 明确支持 ``当前键路径'' 的概念：如果某个 ``当前键路径'' 正在生效，那么 \meta{option list} 将在一个同样使用该当前键路径的上下文中被设置。从技术上讲，这意味着 |/.style| 用 |\pgfkeysalso| 来设置 \meta{option list}，也就是说它并不像上面所说的那样使用 |\pgfkeys|。

\end{handler}

\begin{command}{\pgfkeysalso\marg{comma--separated key--value pairs}}
	这个宏与 |\pgfkeys|\marg{comma--separated key--value pairs}\emph{几乎}相同。事实上，如果其参数中的任何赋值都使用完全限定路径（就像本文档到目前为止所做的那样），这两种调用就是等价的。

	区别在于它们如何处理相对于某个当前键路径\emph{相对}的键，这个概念将在下一小节解释。

	这两个宏的区别在于：|\pgfkeys| 在处理其参数之前会把当前键路径重置为 |/|，而 |\pgfkeysalso| 不改变当前键路径。因此，|\pgfkeysalso| 只在某个 code-key 的宏体内部（例如 |/.style|）才有用。
\end{command}

\subsubsection{当前键路径}
待补充

\subsubsection{键过滤}
待补充

\endgroup


%\subsection{Conditional Expressions}

%\subsection{Loops}

%\subsection{The Problem of Macro--Append Runtime}

%\section{Utility Functions of \PGF}

%\section{Utility Function of \PGFPlots}

\section{特殊技巧}

\subsection{处理参数中的 \#}
我不止一次遇到下面这个难题：我想收集一个含有井号 `|#|' 的参数。这本身并不特别困难，但当要处理得到的参数时，它却可能导致一大堆奇怪的错误信息！请看
\begin{codeexample}[code only]
\def\collectargument#1{%
	\def\collectedcontent{#1}%
	\ifx\collectedcontent\empty
		It is empty.
	\else
		It is not empty, executing it: #1.
	\fi
}%

\collectargument{}% works

\collectargument{something}% works

\collectargument{% does not work!
	\def\something#1{which depends on #1}
}%
\end{codeexample}
这个例子中的代码相当简单：|\collectargument| 宏接受一个参数，并检查它是否为空（使用 |\ifx|，这是一种常见而可靠的判空方法）。如果不为空，就执行它。|\collectargument| 宏在大多数情况下都管用。更确切地说：只要其参数中\emph{没有}井号，它就管用！在我们的例子中，第三次调用失败，报错 ``Illegal parameter number in definition of |\collectedcontent|.''，它发生在 |\def\collectedcontent{#1}| 这一行（由于参数展开的特殊含义，\TeX\ 给出这条信息是有原因的）。

解决办法：用下面的方式重新定义 |\collectargument| 宏

\begin{codeexample}[code only]
\def\collectargument#1{%
	\toks0={#1}%
	\edef\collectedcontent{\the\toks0}%
	\ifx\collectedcontent\empty
		It is empty.
	\else
		It is not empty, executing it: #1.
	\fi
}%

\end{codeexample}
\noindent（你或许想为这个任务分配一个临时记号寄存器）。区别在哪里？是这样的：|\toks0={#1}| 这个赋值不会给井号 |#| 引入任何特殊含义，|\the\toks0| 也不会。不过请注意，这需要用 |\edef\collectedcontent| 而不是 |\def\collectedcontent|，因为 |\the| 语句需要被展开。一切都按预期工作。

\printindex

\bibliographystyle{abbrv} %gerapali} %gerabbrv} %gerunsrt.bst} %gerabbrv}% gerplain}
\nocite{schwartz}
\nocite{pgfplots}
\bibliography{pgfplots}
\end{document}
